% Q-0448: the proofs that Lean checked, section S3 of the supplement (supplement.tex) and an
% appendix of the full document (full.tex). Every proof and statement below was cut unchanged
% from appendix-isotropic.tex by records/affine/split_q0448.py, with the comment lines above
% it; the header of that file has their history. Six statements that are used in these proofs
% only are stated here. `linkage check` reads this file as it reads every .tex file of the
% directory: a statement by its marker, a proof by the \ref of its heading.
\section{Proofs of the machine-checked statements}\label{supp:proofs}

This section has the proofs of the statements of the paper that Lean checked, in the order of
\S\ref{sec:isotropic}. The proofs of the classification are in the paper
(\S\ref{app:classification-proofs}), and so is the proof of
Proposition~\ref{prop:profile-one-dimension-down} (\S\ref{app:print-proofs}). Six lemmas are
stated here and not in the paper:
Lemmas~\ref{lem:positive-transform-origin}, \ref{prop:origin-rate-d},
\ref{lem:threshold-monotone}, \ref{lem:threshold-limit}, \ref{lem:threshold-asymptotic}
and~\ref{lem:filtration-top}. The paper cites them by these numbers, and where a statement or
a figure of its main text needs a formula of one of them, the paper gives the formula.
% Q-0448, the blind draft review of the split (flag 3): "that only these proofs use" was not
% true. The main text of the paper uses five of the six: the Gaussian density g_t of the first
% (displayed before Theorem thm:origin-boundedness-d), the constant R(d,c) and clause (1) of
% the second (displayed, and stated, before Proposition prop:origin-rate-profile-d), the
% criterion of the third (before Definition def:filtration-thresholds), the two bounds of the
% fourth (displayed after Theorem thm:separating-family-critical-values, and drawn in Figure
% fig:thresholds), and the sixth as what the cited fact C6 enters with (section 1.1, with
% what it states). The fifth is clause (3) of that theorem with sharper bounds.

\subsection{The axioms}\label{supp:proofs-axioms}

The proof is the line paper's, carried from the line to $\RR^d$: clause (1) is proved as
\prog{V}{Lem.~4.1} is, with balls in place of intervals, and the passage between the operators
and the kernels follows the proofs of \prog{V}{Lem.~4.2 and Lem.~6.1}, with $\GL(d)$ in place of
the dilations. It is written out in full because those statements are on the line.

% Q-0406: the proof of this paper's Lemma lem:operators-to-kernels, assembled from the proofs of
% record of the two blueprint nodes it renders, in the order of the paper's clauses:
% lem:operators-to-kernels (the two facts used throughout; clause (1) whole; clauses (2), (3), (5))
% and lem:kernels-to-operators (clauses (1) and (3)). It replaces the proof of Q-0334, which had the
% passage from kernels to operators only. The paragraphs are the proofs of record's, with these
% differences, none of which adds or drops a step:
%   - the clauses carry the numbers of the paper's lemma; the node's clause (2), the transport
%     identity, is the paragraph headed "Transport";
%   - lem:operators-to-kernels(5) says "the same argument on the parameter set" as its clause (4):
%     that argument is written out on {0 <= s <= t}, with (I1) and (I2) for (K1) and (K2); so is
%     lem:kernels-to-operators(3), as Q-0334 wrote it;
%   - (I2) is derived from (A7) at every f, the hypothesis of the paper's clause (2), where the node
%     reads (A7) at f = varphi only;
%   - varphi (the standard Gaussian density; gamma before Q-0327, which renamed it at the
%     nodes, and this copy follows) and D_A are defined here (the node defines them in its
%     statement, the paper's statement does not need them); rotation invariance is cited as (A3), its number in
%     Definition def:isotropic-operator-family;
%   - the transport display is broken over two lines (layout only);
%   - the frequency on R^d is xi, as in the paper's clause (2) since Q-0417, where the proofs of
%     record write omega (the letter only, 31 occurrences).
% Pinned against both nodes' proofs.
\begin{proof}[Proof of Lemma~\ref{lem:operators-to-kernels}]
For $A \in \GL(d)$ write $D_Af := |\det A|^{-1}f(A^{-1}\,\cdot)$ on $L^1(\RR^d)$. Then
$R_Q = D_Q$ for $Q \in \Orth(d)$, since $|\det Q| = 1$, and $D_\lambda = D_{\lambda I}$. Let
$\varphi(x) := (2\pi)^{-d/2}e^{-|x|^2/2}$ be the standard Gaussian density, and for $\sigma > 0$
write $\varphi_\sigma(x) := \sigma^{-d}\varphi(x/\sigma)$, the centred Gaussian density with
covariance $\sigma^2I$.

Two facts are used throughout. For a probability measure $\mu$ and $f \in L^1(\RR^d)$, $\mu * f$
is the $L^1(\RR^d)$-valued Bochner integral $\int T_yf\,\mu(dy)$: the map $y \mapsto T_yf$ is
continuous because translation acts continuously on $L^1(\RR^d)$, and it is bounded by
$\|f\|_1$. And with $\widehat h(\xi) := \int e^{i\langle\xi,x\rangle}h(x)\,dx$ for
$h \in L^1(\RR^d)$, Fubini gives the pairing
\[
  \widehat{\mu * f}(\xi) = \widehat\mu(\xi)\,\widehat f(\xi),
  \qquad |\widehat h(\xi)| \le \|h\|_1 ,
\]
and $\widehat\varphi(\xi) = e^{-|\xi|^2/2} > 0$.

\emph{(1), the interchange identity.} For $f, g \in L^1(\RR^d)$, $f * g = \int f(y)\,T_yg\,dy$ as
a Bochner integral, dominated by $|f(y)|\,\|g\|_1$. A bounded operator commutes with an
$L^1(\RR^d)$-valued integral, so (A1) and (A2) give $\Phi(f * g) = \int f(y)\,T_y\Phi g\,dy
= f * \Phi g$.

\emph{The approximants.} Let $B_\varepsilon := \{|x| < \varepsilon\}$, $|B_\varepsilon|$ its
volume, $\rho_\varepsilon := |B_\varepsilon|^{-1}\mathbf 1_{B_\varepsilon}$ and
$h_\varepsilon := \Phi\rho_\varepsilon$. By (A4) and (A5), $h_\varepsilon \ge 0$ with
$\int h_\varepsilon = 1$, so $\nu_\varepsilon := h_\varepsilon\,dx$ is a probability measure on
$\RR^d$. For every $f$, $\|f * \rho_\varepsilon - f\|_1 \le \sup_{|y| < \varepsilon}
\|T_yf - f\|_1 \to 0$ as $\varepsilon \downarrow 0$.

\emph{Tightness.} By the interchange identity $\rho_1 * h_\varepsilon = \Phi(\rho_1 *
\rho_\varepsilon)$, which tends to $\Phi\rho_1$ in $L^1(\RR^d)$ by (A1). Along
$\varepsilon_n := 1/(n+1)$ the functions $\rho_1 * h_{\varepsilon_n}$ therefore have uniformly
small tails: given $\eta > 0$ there is $R$ with $\int_{|x| > R}(\rho_1 * h_{\varepsilon_n}) < \eta$
for every $n$ (choose $N$ with $\|\rho_1 * h_{\varepsilon_n} - \Phi\rho_1\|_1 < \eta/2$ for
$n \ge N$, then $R$ for $\Phi\rho_1$ and the first $N$ terms). Convolution with a nonnegative
unit-mass density carried by the unit ball moves mass by at most $1$: for $|y| > R + 1$ the
translate $\rho_1(\cdot - y)$ is carried by $\{|x| > R\}$, so by Tonelli
\[
  \int_{|x| > R}(\rho_1 * h_\varepsilon)(x)\,dx
  = \int h_\varepsilon(y)\int_{|x| > R}\rho_1(x - y)\,dx\,dy
  \ge \nu_\varepsilon\bigl(\{|y| > R + 1\}\bigr) .
\]
So the sequence $(\nu_{\varepsilon_n})$ is tight, with the closed balls of radius $R + 1$ as
compacts.

\emph{The limit.} By Prokhorov's theorem a subsequence $\nu_{\varepsilon_{n_k}}$ converges weakly
to a probability measure $\mu$ on $\RR^d$.

\emph{Identification.} Let $\Psi \in L^\infty(\RR^d)$, acting on $L^1(\RR^d)$ by
$\Psi(h) := \int\Psi h$, and $f \in L^1(\RR^d)$. The map $y \mapsto \Psi(T_yf)$ is continuous and
bounded by $\|\Psi\|_\infty\|f\|_1$, hence a test function for weak convergence. Since
$f * h_\varepsilon = \int h_\varepsilon(y)\,T_yf\,dy$ and $\Psi$ commutes with the integral,
\[
  \Psi(f * h_{\varepsilon_{n_k}}) = \int\Psi(T_yf)\,\nu_{\varepsilon_{n_k}}(dy)
  \ \xrightarrow{\ k \to \infty\ }\ \int\Psi(T_yf)\,\mu(dy) = \Psi(\mu * f) .
\]
On the other hand $f * h_\varepsilon = \Phi(f * \rho_\varepsilon) \to \Phi f$ in $L^1(\RR^d)$, by
the interchange identity and (A1). Hence $\Psi(\Phi f - \mu * f) = 0$ for every $\Psi$; taking
$\Psi$ the sign of $\Phi f - \mu * f$ gives $\Phi f = \mu * f$.

\emph{Uniqueness.} If $\mu$ and $\nu$ are probability measures with
$\mu * \varphi = \nu * \varphi$, the pairing gives
$\widehat\mu(\xi)\,\widehat\varphi(\xi) = \widehat\nu(\xi)\,\widehat\varphi(\xi)$, so
$\widehat\mu = \widehat\nu$ because $\widehat\varphi > 0$, and $\mu = \nu$ because a finite measure
on $\RR^d$ is determined by its Fourier transform. In particular two measures representing the
same operator agree.

\emph{Converse.} For a probability measure $\mu$, $\|\mu * f\|_1 \le \|f\|_1$ by Minkowski's
integral inequality (A1); $T_a(\mu * f) = \mu * T_af$, both being
$x \mapsto \int f(x - a - y)\,\mu(dy)$ (A2); $\mu * f \ge 0$ when $f \ge 0$ (A4); and
$\int\mu * f = \int f$ for $f \ge 0$ by Tonelli (A5).

\emph{(2), from $L^1$ to the transforms.} Suppose $\mu_n * f \to \mu * f$ in $L^1(\RR^d)$ for
every $f$, and take $f = \varphi$. By the pairing,
$|\widehat\mu_n(\xi) - \widehat\mu(\xi)|\,\widehat\varphi(\xi)
\le \|\mu_n * \varphi - \mu * \varphi\|_1 \to 0$; divide by $\widehat\varphi(\xi) > 0$. The test
density must have a zero-free transform, which is why the indicator of a ball would not do.

\emph{(2), from the transforms to $L^1$: the smoothed measures.} For a probability measure $\nu$
and $f, g \in L^1(\RR^d)$ three facts are used, each by Fubini's theorem:
$\|\nu * f\|_1 \le \|f\|_1$; $\|g * f\|_1 \le \|g\|_1\|f\|_1$; and
$\nu * (g * f) = (\nu * g) * f$, where $\nu * g$ is the integrable function
$x \mapsto \int g(x - y)\,\nu(dy)$. The Gaussian integral gives, for every $z \in \RR^d$,
\[
  \varphi_\sigma(z) = (2\pi)^{-d}\int e^{-i\langle\xi,z\rangle}e^{-\sigma^2|\xi|^2/2}\,d\xi .
\]
Put $z = x - y$ and integrate in $\nu(dy)$. The integrand is bounded by $e^{-\sigma^2|\xi|^2/2}$,
which is integrable for $d\xi \otimes \nu$, so Fubini's theorem gives, for every probability
measure $\nu$ and every $x \in \RR^d$,
\[
  (\nu * \varphi_\sigma)(x) = (2\pi)^{-d}\int e^{-i\langle\xi,x\rangle}\,\widehat\nu(\xi)\,
  e^{-\sigma^2|\xi|^2/2}\,d\xi .
\]
Put $h_n := \mu_n * \varphi_\sigma$ and $h := \mu * \varphi_\sigma$. Since $|\widehat\mu_n| \le 1$ and
$\widehat\mu_n \to \widehat\mu$ pointwise, dominated convergence gives $h_n(x) \to h(x)$ for every
$x$. The $h_n$ and $h$ are nonnegative with integral $1$ (Tonelli), so $\int(h_n - h) = 0$ and
\[
  \|h_n - h\|_1 = 2\int(h - h_n)^+ \xrightarrow{\ n \to \infty\ } 0 ,
\]
by dominated convergence again, since $0 \le (h - h_n)^+ \le h$ and $(h - h_n)^+ \to 0$ pointwise.

\emph{(2), the approximation.} By Minkowski's integral inequality and the substitution $y = \sigma z$,
\[
  \|\varphi_\sigma * f - f\|_1 \le \int\varphi_\sigma(y)\,\|T_yf - f\|_1\,dy
  = \int\varphi(z)\,\|T_{\sigma z}f - f\|_1\,dz ,
\]
which tends to $0$ as $\sigma \downarrow 0$ by dominated convergence: the integrand is bounded by
$2\|f\|_1\varphi(z)$ and tends to $0$ pointwise, translation being continuous on $L^1(\RR^d)$. Now
split $f = (f - \varphi_\sigma * f) + \varphi_\sigma * f$. By the three facts,
$\mu_n * (\varphi_\sigma * f) = h_n * f$ and $\mu * (\varphi_\sigma * f) = h * f$, so
\[
  \|\mu_n * f - \mu * f\|_1 \le 2\,\|f - \varphi_\sigma * f\|_1 + \|(h_n - h) * f\|_1
  \le 2\,\|f - \varphi_\sigma * f\|_1 + \|h_n - h\|_1\,\|f\|_1 .
\]
Given $\varepsilon > 0$, choose $\sigma$ so that the first term is below $\varepsilon$, and then $n_0$ so
that the second is below $\varepsilon$ for $n \ge n_0$.

\emph{Transport.} For a probability measure $\mu$, $A \in \GL(d)$, $f \in L^1(\RR^d)$ and
$x \in \RR^d$, substituting $z = Ay$,
\[
\begin{aligned}
  \bigl((A_*\mu) * D_Af\bigr)(x) &= \int |\det A|^{-1}f\bigl(A^{-1}(x - z)\bigr)\,(A_*\mu)(dz) \\
  &= |\det A|^{-1}\int f(A^{-1}x - y)\,\mu(dy) = D_A(\mu * f)(x) .
\end{aligned}
\]
$D_A$ is invertible on $L^1(\RR^d)$ with inverse $D_{A^{-1}}$, so applying the identity to
$D_A^{-1}g$ gives $D_A\,(\mu * D_A^{-1}g) = (A_*\mu) * g$ for every $g$: the operator
$D_A\,(\mu * \cdot\,)\,D_A^{-1}$ is convolution by $A_*\mu$. It is convolution by a probability
measure $\nu$ if and only if $A_*\mu = \nu$, by the uniqueness in (1).

\emph{(3).} By (1), each $\Phi_{s,t}$ is convolution by a unique probability measure
$\mu_{s,t}$. \emph{(I1).} $\Phi_{t,t} = \mathrm{Id}$ is convolution by $\delta_0$, so
$\mu_{t,t} = \delta_0$ by uniqueness. By Fubini
$\mu_{s,t} * (\mu_{r,s} * f) = (\mu_{s,t} * \mu_{r,s}) * f$, so (A6) says that $\mu_{r,t}$ and
$\mu_{s,t} * \mu_{r,s}$ represent the same operator, and they are equal by uniqueness.
\emph{(I2).} $\{0 \le s \le t\}$ is a subset of $\RR^2$, so continuity on it is sequential. If
$(s_n,t_n) \to (s,t)$ in it, (A7) gives $\mu_{s_n,t_n} * f \to \mu_{s,t} * f$ in $L^1(\RR^d)$ for
every $f$, and (2) gives $\widehat\mu_{s_n,t_n}(\xi) \to \widehat\mu_{s,t}(\xi)$ for every
$\xi$. \emph{(I3).} $R_Q = D_Q$, and (A3) reads $D_Q\Phi_{s,t}D_Q^{-1} = \Phi_{s,t}$; the
transport with $\nu = \mu = \mu_{s,t}$ gives $Q_*\mu_{s,t} = \mu_{s,t}$. \emph{(I4).}
$D_\lambda = D_{\lambda I}$, and (A8) reads
$D_{\lambda I}\Phi_{s,t}D_{\lambda I}^{-1} = \Phi_{S_\lambda s,\,S_\lambda t}$; the transport gives
$(\lambda I)_*\mu_{s,t} = \mu_{S_\lambda s,\,S_\lambda t}$. \emph{(ND).} If $\mu_{s,t} = \delta_0$
with $s < t$ then $\Phi_{s,t} = \mathrm{Id}$, against (ND).

\emph{(4).} Each $\Phi_{s,t}$ is linear and satisfies (A1), (A2), (A4) and (A5) by the converse
in (1). \emph{Cascade.} $\Phi_{t,t} = \delta_0 * \cdot = \mathrm{Id}$ by (I1), and for
$r \le s \le t$, by the associativity of convolution (Fubini) and (I1),
$\Phi_{s,t}\Phi_{r,s}f = \mu_{s,t} * (\mu_{r,s} * f) = (\mu_{s,t} * \mu_{r,s}) * f
= \Phi_{r,t}f$, which is (A6). \emph{Continuity.} $\{0 \le s \le t\}$ is a subset of $\RR^2$, so
continuity on it is sequential. If $(s_n,t_n) \to (s,t)$ in it, (I2) gives
$\widehat\mu_{s_n,t_n}(\xi) \to \widehat\mu_{s,t}(\xi)$ for every $\xi$, and (2) gives
$\Phi_{s_n,t_n}f \to \Phi_{s,t}f$ in $L^1(\RR^d)$ for every $f$, which is (A7). \emph{(ND).} If
$\Phi_{s,t} = \mathrm{Id}$ with $s < t$, then $\mu_{s,t}$ and $\delta_0$ represent the same
operator, so $\mu_{s,t} = \delta_0$ by uniqueness, against (ND) of the family.
\emph{Covariance.} For $Q \in \Orth(d)$, $D_Q = R_Q$, and (I3) with the
transport gives $R_Q\Phi_{s,t}R_Q^{-1} = \Phi_{s,t}$, that is
$\Phi_{s,t}R_Q = R_Q\Phi_{s,t}$, which is (A3). For $\lambda > 0$, $D_\lambda = D_{\lambda I}$, and
(I4) with the transport gives $D_\lambda\Phi_{s,t}D_\lambda^{-1} = \Phi_{S_\lambda s,\,S_\lambda t}$,
which is (A8). \emph{The measures.} $\Phi_{s,t}$ is convolution by $\mu_{s,t}$, so by the
uniqueness in (1) the measure that (1) gives for it is $\mu_{s,t}$.
\end{proof}

\subsection{Reduction to the line}\label{supp:proofs-reduction}

In the next proof the argument for the continuity (A7) is the one by which the line paper
verifies that axiom in the sufficiency half of its classification (the proof of
\prog{V}{Thm.~7.3}).

\begin{proof}[Proof of Lemma~\ref{lem:isotropic-marginal-family}]
\emph{The kernels.} By (I3) each $\mu_{s,t}$ is rotation invariant, so by
Lemma~\ref{lem:radial-marginal}(1) each $\kappa_{s,t}$ is a symmetric probability measure with
$\widehat\mu_{s,t}(\xi) = \widehat\kappa_{s,t}(|\xi|)$.

\emph{(A1), (A2), (A4), (A5).} For a probability measure $\kappa$ on $\RR$ and $f \in L^1(\RR)$,
$(\kappa * f)(x) = \int f(x-y)\,\kappa(dy)$. Tonelli gives $\|\kappa * f\|_1 \le \|f\|_1$, so the
operator is bounded, and Fubini gives $\int \kappa * f = \int f$. A nonnegative $f$ gives a
nonnegative $\kappa * f$, and $\tau_a(\kappa * f) = \kappa * \tau_af$ is a change of variables.

\emph{(A3).} As $\kappa_{s,t}$ is symmetric, $R(\kappa_{s,t} * f)(x) = \int f(-x-y)\,\kappa_{s,t}(dy)
= \int f(-x+y)\,\kappa_{s,t}(dy) = (\kappa_{s,t} * Rf)(x)$.

\emph{(A6).} By (I1) and Lemma~\ref{lem:radial-marginal}(3), $\kappa_{t,t} = \pi_*\delta_0 = \delta_0$
and $\kappa_{r,t} = \pi_*(\mu_{s,t} * \mu_{r,s}) = \kappa_{s,t} * \kappa_{r,s}$. So
$\Phi_{t,t} = \mathrm{Id}$, and $\Phi_{s,t}\Phi_{r,s}f = \kappa_{s,t} * (\kappa_{r,s} * f)
= (\kappa_{s,t} * \kappa_{r,s}) * f = \Phi_{r,t}f$, by Fubini in the middle step.

\emph{(A8).} By (I4) and Lemma~\ref{lem:radial-marginal}(3),
$(\lambda I)_*\kappa_{s,t} = \pi_*\bigl((\lambda I)_*\mu_{s,t}\bigr) = \kappa_{S_\lambda s,S_\lambda t}$.
For a probability measure $\kappa$ on $\RR$, substituting $z = \lambda y$,
\[
  D_\lambda(\kappa * f)(x) = \int \lambda^{-1}f\Bigl(\frac{x - \lambda y}{\lambda}\Bigr)\kappa(dy)
   = \int (D_\lambda f)(x - z)\,\bigl((\lambda I)_*\kappa\bigr)(dz) ,
\]
so $D_\lambda\Phi_{s,t}f = \kappa_{S_\lambda s,S_\lambda t} * D_\lambda f
= \Phi_{S_\lambda s,S_\lambda t}D_\lambda f$.

\emph{(ND).} Let $s < t$. By (ND) of the family $\mu_{s,t} \ne \delta_0$, so
$\kappa_{s,t} \ne \delta_0$ by Lemma~\ref{lem:radial-marginal}(2). If $\kappa_{s,t} * f = f$ for every
$f \in L^1(\RR)$, take $f$ the standard Gaussian density: its transform is positive, and the
transform of the identity reads $\widehat\kappa_{s,t}\widehat f = \widehat f$, so
$\widehat\kappa_{s,t} \equiv 1$ and $\kappa_{s,t} = \delta_0$, a contradiction. So
$\Phi_{s,t} \ne \mathrm{Id}$.

\emph{(A7), the collapsing increment.} For $f \in L^1(\RR)$ put
$\Theta_f(y) := \|\tau_yf - f\|_1$, a function that is bounded by $2\|f\|_1$, continuous, and zero
at $y = 0$, because translation acts continuously on $L^1(\RR)$. For a probability measure $\kappa$,
$\kappa * f - f = \int(\tau_yf - f)\,\kappa(dy)$, so Minkowski's integral inequality gives
\begin{equation}\label{eq:collapsing-increment}
  \|\kappa * f - f\|_1 \le \int \Theta_f\,d\kappa .
\end{equation}
Moreover $\Theta_{\kappa' * g} \le \Theta_g$ for every probability measure $\kappa'$ and every
$g \in L^1(\RR)$, because $\tau_y(\kappa' * g) - \kappa' * g = \kappa' * (\tau_yg - g)$ and
convolution by $\kappa'$ does not increase the norm. Now let pairs $(s_n,t_n)$ with
$s_n \le t_n$ converge to a diagonal point $(t_0,t_0)$. By (I2) at $\xi = \omega e$,
$\widehat\kappa_{s_n,t_n}(\omega) = \widehat\mu_{s_n,t_n}(\omega e) \to \widehat\mu_{t_0,t_0}(\omega e) = 1$
for every $\omega \in \RR$. \Levy's continuity theorem on the line \restatedp{V}{Prop.~2.8}{sec:line-facts} says that if
the transforms of probability measures on $\RR$ converge at every point to the transform of a
probability measure, the measures converge weakly to it. So $\kappa_{s_n,t_n} \to \delta_0$
weakly. Since $\Theta_g$ is bounded and continuous, $\int\Theta_g\,d\kappa_{s_n,t_n} \to
\Theta_g(0) = 0$ for every $g$.

\emph{(A7), the general pair.} Let $(s_n,t_n) \to (s,t)$ with $s_n \le t_n$. If $s = t$ then
$\Phi_{s,t} = \mathrm{Id}$ and \ref{eq:collapsing-increment} with the previous paragraph gives
$\Phi_{s_n,t_n}f \to f$. If $s < t$ put $u_n := \max\{s_n,s\}$ and $v_n := \min\{t_n,t\}$, so
that $u_n \to s$, $v_n \to t$ and, for all large $n$, $s_n \le u_n \le v_n \le t_n$ and
$s \le u_n \le v_n \le t$. By (A6),
\[
  \Phi_{s_n,t_n}f = \Phi_{v_n,t_n}\,\Phi_{u_n,v_n}\,\Phi_{s_n,u_n}f ,
  \qquad
  \Phi_{s,t}f = \Phi_{v_n,t}\,\Phi_{u_n,v_n}\,\Phi_{s,u_n}f .
\]
The left endpoints are moved first, where contraction absorbs them: with
$g_n := \Phi_{u_n,v_n}\Phi_{s_n,u_n}f$ and $h_n := \Phi_{u_n,v_n}\Phi_{s,u_n}f$,
\[
  \|g_n - h_n\|_1 \le \|\Phi_{s_n,u_n}f - f\|_1 + \|f - \Phi_{s,u_n}f\|_1 \to 0 ,
\]
both increments collapsing onto $(s,s)$. The right endpoints are moved last, where
\ref{eq:collapsing-increment} is read at $h_n$ and $\Theta_{h_n} \le \Theta_f$:
\[
  \|\Phi_{v_n,t_n}h_n - \Phi_{v_n,t}h_n\|_1
   \le \int\Theta_f\,d\kappa_{v_n,t_n} + \int\Theta_f\,d\kappa_{v_n,t} \to 0 ,
\]
both increments collapsing onto $(t,t)$. Hence
$\|\Phi_{s_n,t_n}f - \Phi_{s,t}f\|_1 \le \|g_n - h_n\|_1
+ \|\Phi_{v_n,t_n}h_n - \Phi_{v_n,t}h_n\|_1 \to 0$.
% Q-0344 (the author's review of section 2): the proof of record's last paragraph, "What is
% machine-checked where", is not printed. It describes the Lean and what the line paper's checked
% construction rests on; the paper says neither in a proof (WRITING.md section 3; section 1.1 owns the
% trust base). The rest of the proof is the proof of record verbatim.
\end{proof}

\begin{proof}[Proof of Proposition~\ref{prop:isotropic-bridge}]
\emph{(1).} Each $\mu_{s,t}$ is a mixture of probability measures, and $u \mapsto g^{(d)}_u(B)$
is measurable for every Borel $B$, so it is a probability measure. By Fubini and
$\widehat{g^{(d)}_u}(\xi) = e^{-u|\xi|^2/2}$,
\[
  \widehat\mu_{s,t}(\xi) = \int e^{-u|\xi|^2/2}\,\mu^{\mathrm I}_{s^2,t^2}(du)
   = \exp\bigl[-\bigl(F_{\mathrm I}(t^2|\xi|^2/2) - F_{\mathrm I}(s^2|\xi|^2/2)\bigr)\bigr] ,
\]
which is $\exp[-(F(t|\xi|) - F(s|\xi|))]$.
The axioms are read off this transform exactly as in the sufficiency half of
Theorem~\ref{thm:isotropic-classification}, with $\chi = \mathrm{id}$ and $S_\lambda t = \lambda t$. That
half used three properties of $F$, and they hold here. $F$ is continuous on $[0,\infty)$, because
$F_{\mathrm I}$ is continuous by dominated convergence: its integrand is at most
$\max\{1,\Sigma\}(1 \wedge u)\,k_{\mathrm I}(u)/u$ on $\sigma \le \Sigma$, integrable by the two
conditions. $F(0) = F_{\mathrm I}(0) = 0$. And $F \not\equiv 0$, because $F_{\mathrm I} \not\equiv 0$
and $r^2/2$ runs over $[0,\infty)$. So $(\mu_{s,t})$ is an
isotropic cascade family, and Theorem~\ref{thm:isotropic-classification} applies. Its conclusion is met by
$(\mathrm{id},F)$ once $F \in \mathcal{A}_d$, and the argument of its necessity half gives that.
% Q-0406 (item 5): the proof of record restates [B, Lem. 4.2] here after "restated:"; it is
% Proposition prop:line-bridge of section 2, and the restatement is dropped. Here $k_{\mathrm I}$
% and "the two conditions" above are those of Definition def:line-causal. Pinned.
$F$ is admissible by \restated{B}{Lem.~4.2}{prop:line-bridge}. And $\mu_{0,1}$ has transform $e^{-F(|\cdot|)}$ and, by
(I1), residuals $\mu_{\lambda,1}$ in \ref{eq:self-decomposable}, as $\widehat\mu_{0,\lambda}(\xi)
= e^{-F(\lambda|\xi|)} = \widehat\mu_{0,1}(\lambda\xi)$. So $F \in \mathcal{A}_d$, and $d$ was
arbitrary.

\emph{(2) and (3)} are the exponent identities $\gamma\log(1 + r^2/2)$ and
$(r^2/2)^\alpha = 2^{-\alpha}r^{2\alpha}$, read at $r = t|\xi|$ through (1). At $\alpha = 1$,
$\widehat\mu_{s,t}(\xi) = e^{-(t^2-s^2)|\xi|^2/2}$, the transform of $g^{(d)}_{t^2-s^2}$.

\emph{(4).} The mixture is a probability measure as in (1), and by Fubini its transform is
$E\,e^{-t^2T_1|\xi|^2/2}$. On the line, \prog{B}{Prop.~5.10(2)} says that the law of the mixture of
$g^{(1)}_u$ against $T_1$ has transform
$\tfrac{2^{1-\vartheta}}{\Gamma(\vartheta)}|\omega|^\vartheta K_\vartheta(|\omega|)$. By the
same Fubini that transform is $E\,e^{-T_1\omega^2/2}$. Reading it at $\omega = t|\xi|$ gives the
display. When $\mu^{\mathrm I}_{0,1}$ is the law of $T_1$, by (1) and its Laplace transform at
$t^2|\xi|^2/2$, $\widehat\mu_{0,t}(\xi) = e^{-F_{\mathrm I}(t^2|\xi|^2/2)} = E\,e^{-t^2T_1|\xi|^2/2}$,
which is the transform of the mixture. Two probability measures on $\RR^d$ with the same
transform are equal, so $\mu_{0,t}$ is the mixture, and its exponent is in $\mathcal{A}_d$ by (1).
\end{proof}

\subsection{The kernel at its centre}\label{supp:proofs-origin}

% Q-0414: Feller, vol. II (2nd ed.), XV.3, Corollary, p. 510, read at the page image
% (READING-feller2009introduction.md). The lemma is an identity and one implication, in R^d;
% the equivalence is the lemma with C1, said after it. How Feller proves his corollary is not
% retold.
The proof of Theorem~\ref{thm:origin-boundedness-d} rests on two lemmas. The first goes back to an
observation of Feller, on the line: a nonnegative characteristic function is integrable exactly
when the law has a bounded density (\sscite{feller2009introduction}, XV.3, Corollary, p.~510).
The lemma is one half of it in $\RR^d$. It gives the value at the centre of the law blurred by a
Gaussian as an integral of the transform, and from it the direction that the unbounded side of
the theorem uses: a bound on the law near the origin makes the transform integrable.

% shared with blueprint lem:positive-transform-origin
\begin{lemma}[a nonnegative transform: bounded near the origin only if integrable]\label{lem:positive-transform-origin}
Let $d \ge 1$ and let $\mu$ be a probability measure on $\RR^d$ whose characteristic function
$\widehat\mu(\xi) = \int e^{i\langle\xi,x\rangle}\mu(dx)$ is real and nonnegative at every
$\xi \in \RR^d$. For $t > 0$ let $g_t(x) = (2\pi t)^{-d/2}e^{-|x|^2/2t}$ be the heat kernel.
\begin{enumerate}
\item For every $t > 0$,
\[
  \int_{\RR^d} g_t\,d\mu = (2\pi)^{-d}\int_{\RR^d}\widehat\mu(\xi)\,e^{-t|\xi|^2/2}\,d\xi ,
\]
and as $t \downarrow 0$ this increases to $(2\pi)^{-d}\int_{\RR^d}\widehat\mu(\xi)\,d\xi \in (0,\infty]$.
\item If there are $\delta > 0$ and $M < \infty$ with $\mu(B) \le M|B|$ for every Borel set
$B \subseteq \{|x| < \delta\}$, then $\int_{\RR^d}\widehat\mu(\xi)\,d\xi \le (2\pi)^d M$. The
hypothesis holds in particular when $\mu$ has a density with a representative bounded by $M$ on
$\{0 < |x| < \delta\}$.
\end{enumerate}
\end{lemma}

With C1 (Theorem~\ref{thm:cited-A1}), which gives the other direction, the lemma is the form in
$\RR^d$ of Feller's equivalence: a law with a nonnegative transform has a density bounded near
the origin exactly when its transform is integrable. The law of every isotropic origin datum has
the positive transform $e^{-F_d}$. The second lemma, Lemma~\ref{lem:isotropic-exponent-log},
bounds that transform from below.

\begin{proof}[Proof of Lemma~\ref{lem:positive-transform-origin}]
\emph{Clause (1).} The one-dimensional Gaussian integral
$\int_{\RR}e^{-isy}e^{-ts^2/2}\,ds = (2\pi/t)^{1/2}e^{-y^2/2t}$, multiplied over the $d$
coordinates, gives $g_t(x) = (2\pi)^{-d}\int_{\RR^d}e^{-i\langle\xi,x\rangle}e^{-t|\xi|^2/2}\,d\xi$.
The integrand has modulus $e^{-t|\xi|^2/2}$, which is integrable against $d\xi \otimes \mu(dx)$
because $\mu$ is finite, so Fubini's theorem gives
\[
  \int g_t\,d\mu = (2\pi)^{-d}\int e^{-t|\xi|^2/2}\int e^{-i\langle\xi,x\rangle}\mu(dx)\,d\xi
   = (2\pi)^{-d}\int e^{-t|\xi|^2/2}\,\widehat\mu(-\xi)\,d\xi ,
\]
and the substitution $\xi \mapsto -\xi$ gives the display. Since $\widehat\mu \ge 0$, the integrand
increases as $t$ decreases, and monotone convergence gives the limit. It is positive because
$\widehat\mu$ is continuous with $\widehat\mu(0) = 1$.

\emph{Clause (2).} Put $U = \{|x| < \delta\}$. The hypothesis gives
$\int_U h\,d\mu \le M\int_U h\,dx$ for every nonnegative Borel $h$, for indicator functions of
Borel subsets of $U$ directly and in general by simple functions and monotone convergence. With
$h = g_t$, $\int_U g_t\,d\mu \le M\int_{\RR^d}g_t\,dx = M$. Off $U$,
$g_t \le (2\pi t)^{-d/2}e^{-\delta^2/2t}$, so
$\int_{\RR^d\setminus U}g_t\,d\mu \le (2\pi t)^{-d/2}e^{-\delta^2/2t}$, which tends to $0$ as
$t \downarrow 0$. By clause (1), $(2\pi)^{-d}\int\widehat\mu = \lim_{t\downarrow0}\int g_t\,d\mu \le
M$. For the last sentence: if $q$ is a density of $\mu$ with $q \le M$ on $U \setminus \{0\}$, then
$\mu(B) = \int_{B\setminus\{0\}}q \le M|B|$ for every Borel $B \subseteq U$, since the origin is a
Lebesgue-null set.
\end{proof}

\begin{proof}[Proof of Lemma~\ref{lem:isotropic-exponent-log}]
Since the Gaussian coefficient is $0$, \ref{eq:isotropic-exponent} and \ref{eq:polar-decomposition}
give, with $\sigma_{d-1}$ the surface measure of $S^{d-1}$,
\[
  F_d(\xi) = \frac{1}{|S^{d-1}|}\int_{S^{d-1}}\varphi\bigl(\langle\xi,\eta\rangle\bigr)\,\sigma_{d-1}(d\eta),
  \qquad
  \varphi(s) := \int_0^\infty\bigl(1 - \cos(rs)\bigr)\,k_d(r)\,\frac{dr}{r} \quad (s \in \RR).
\]
$\varphi$ is even and finite, because $1 - \cos(rs) \le \min\{2, r^2s^2/2\}$ and the two
integrability conditions of Definition~\ref{def:isotropic-origin-datum} hold. As $k_d$ is nonincreasing with
right limit $c$ at the origin, $k_d \le c$ on $(0,\infty)$. And $k_d(r) \to 0$ as $r \to \infty$,
since a nonincreasing $k_d \ge \ell > 0$ would make $\int_1^\infty k_d(r)\,dr/r$ diverge.

\emph{The logarithmic bounds on $\Cin$.} Write $\Cin(X) = \int_0^X(1-\cos s)\,s^{-1}ds$ for
$X \ge 0$, a nondecreasing function. For $X \ge 1$,
$\Cin(X) = \Cin(1) + \log X - \int_1^X s^{-1}\cos s\,ds$. Here $0 \le \Cin(1) \le \int_0^1 \tfrac s2\,ds
= \tfrac14$, because $1 - \cos s \le s^2/2$. Integrating by parts,
$\bigl|\int_1^X s^{-1}\cos s\,ds\bigr| = \bigl|[s^{-1}\sin s]_1^X + \int_1^X s^{-2}\sin s\,ds\bigr|
\le 1 + \sin 1 + 1 \le 3$. For $0 < X < 1$ the lower bound below holds because its right side is
negative. Hence
\begin{equation}\label{eq:cin-log}
  \log X - 3 \le \Cin(X) \quad (X > 0),
  \qquad
  \Cin(X) \le \log X + \tfrac{13}{4} \quad (X \ge 1).
\end{equation}

\emph{Clause (1).} Write $\rho = |\xi| \ge 1$. For $|s| \le \rho$, split $\varphi(s)$ at $r = 1$.
On $(0,1)$ use $k_d \le c$, and on $[1,\infty)$ use $1 - \cos \le 2$:
\[
  \varphi(s) \le c\,\Cin(|s|) + 2C_k \le c\,\Cin(\rho) + 2C_k .
\]
On the sphere $|\langle\xi,\eta\rangle| \le \rho$, and the normalized surface measure has mass $1$.
So $F_d(\xi) \le c\,\Cin(\rho) + 2C_k$, and the upper bound of \ref{eq:cin-log} gives the bound on
$F_d$. The bound on $e^{-F_d}$ follows.

\emph{Clause (2), one variable.} Let $a \ge 1$ and $s > 0$. Substituting $u = ar$ in
$\varphi(as)$,
\[
  \varphi(as) - \varphi(s) = \int_0^\infty \bigl(1 - \cos(us)\bigr)\,h_a(u)\,du ,
  \qquad h_a(u) := \frac{k_d(u/a) - k_d(u)}{u} \ge 0 ,
\]
where both terms on the left are finite. For $0 < \varepsilon < N$,
\[
  \int_\varepsilon^N h_a(u)\,du = \int_{\varepsilon/a}^{N/a}k_d(v)\,\frac{dv}{v}
     - \int_\varepsilon^N k_d(v)\,\frac{dv}{v}
   = \int_{\varepsilon/a}^{\varepsilon}k_d(v)\,\frac{dv}{v} - \int_{N/a}^{N}k_d(v)\,\frac{dv}{v} ,
\]
and every integral here is finite because $k_d \le c$. The first term lies between
$k_d(\varepsilon)\log a$ and $c\log a$, and tends to $c\log a$ as $\varepsilon \downarrow 0$. The
second lies between $0$ and $k_d(N/a)\log a$, and tends to $0$ as $N \to \infty$. By monotone
convergence $\int_0^\infty h_a = c\log a$, so $h_a$ is integrable on $(0,\infty)$, and the
Riemann--Lebesgue lemma gives $\int_0^\infty\cos(us)\,h_a(u)\,du \to 0$ as $s \to \infty$. Hence
\[
  \varphi(as) - \varphi(s) \to c\log a \quad (s \to \infty),
  \qquad 0 \le \varphi(as) - \varphi(s) \le 2c\log a \quad (s > 0).
\]
For $0 < a < 1$ apply this with $1/a$ at the point $as$, which tends to infinity with $s$:
$\varphi(as) - \varphi(s) = -\bigl(\varphi(a^{-1}\cdot as) - \varphi(as)\bigr) \to c\log a$, and
$|\varphi(as) - \varphi(s)| \le 2c|\log a|$.

\emph{Clause (2), the sphere average.} Write $\xi = \rho\widehat\xi$ with $|\widehat\xi| = 1$.
Since $\varphi$ is even,
\[
  F_d(a\xi) - F_d(\xi) = \frac{1}{|S^{d-1}|}\int_{S^{d-1}}
    \Bigl[\varphi\bigl(a\rho|\langle\widehat\xi,\eta\rangle|\bigr)
      - \varphi\bigl(\rho|\langle\widehat\xi,\eta\rangle|\bigr)\Bigr]
    \sigma_{d-1}(d\eta) .
\]
For every $\eta$ with $\langle\widehat\xi,\eta\rangle \ne 0$ the bracket tends to $c\log a$ as
$\rho \to \infty$. The remaining $\eta$ form a $\sigma_{d-1}$-null set: a great subsphere when
$d \ge 2$, and no point at all when $d = 1$. The bracket is bounded by $2c|\log a|$. By dominated
convergence $F_d(a\rho\widehat\xi) - F_d(\rho\widehat\xi) \to c\log a$ as $\rho \to \infty$. The
left side does not depend on $\widehat\xi$, because $F_d$ is rotation invariant, so the limit holds
as $|\xi| \to \infty$. For the equivalent form, $g(au) - g(u) = F_d(\sqrt a\,\sqrt u\,e) -
F_d(\sqrt u\,e) \to c\log\sqrt a = \tfrac c2\log a$.
\end{proof}

\begin{proof}[Proof of Theorem~\ref{thm:origin-boundedness-d}]
\emph{Clause (1).} Write $\rho = |\xi|$, and let $\Cin$ be as in the proof of
Lemma~\ref{lem:isotropic-exponent-log}. Since $a = 0$, \ref{eq:isotropic-exponent} and
\ref{eq:polar-decomposition} give
\[
  -\log\bigl|\widehat\mu(\xi)\bigr| = F_d(\xi)
   = \frac{1}{|S^{d-1}|}\int_{S^{d-1}}\varphi\bigl(\langle\xi,\eta\rangle\bigr)\,\sigma_{d-1}(d\eta),
  \qquad
  \varphi(s) := \int_0^\infty\bigl(1 - \cos(rs)\bigr)\,k_d(r)\,\frac{dr}{r} .
\]
$F_d$ is radial: if $|\xi| = |\xi'|$, the reflection $Q$ in $(\xi - \xi')^\perp$ carries $\xi$ to
$\xi'$, preserves Lebesgue measure and the radial density of $\nu_d$, so substituting $x = Qy$ in
\ref{eq:isotropic-exponent} gives $F_d(\xi') = F_d(\xi)$. It is therefore enough to bound
$F_d(\rho e_1)$. Fix a finite $\beta < c$. If $\beta \le 0$ the bound holds with $C_\beta = 1$,
since $|\widehat\mu| \le 1 \le (1+|\xi|)^{-\beta}$. Let $\beta > 0$ and choose $\beta_1$ with
$\beta < \beta_1 < c$. As $k_d$ is nonincreasing with right limit $c$ at the origin, there is
$r_\beta \in (0,1]$ with $k_d \ge \beta_1$ on $(0,r_\beta)$. The integrand of $\varphi$ is
nonnegative, so discarding $(r_\beta,\infty)$ and substituting gives
$\varphi(s) \ge \beta_1\,\Cin(r_\beta|s|)$.

For $\varepsilon > 0$ let $K_\varepsilon := \{\eta \in S^{d-1} : |\langle e_1,\eta\rangle| \ge
\varepsilon\}$. As $\varepsilon \downarrow 0$ these caps increase to the complement of the great
subsphere $\{\langle e_1,\eta\rangle = 0\}$, which is $\sigma_{d-1}$-null: $\sigma_{d-1}(E)$ is $d$
times the Lebesgue measure of the cone $(0,1)\cdot E$, and the cone over the great subsphere lies in
the hyperplane $e_1^\perp$. So there is $\varepsilon \in (0,1]$ with
$p := \sigma_{d-1}(K_\varepsilon)/|S^{d-1}| \ge \beta/\beta_1$. Put
$L := 3 - \log(\varepsilon r_\beta) \ge 0$. For $\eta \in K_\varepsilon$,
$r_\beta|\langle\rho e_1,\eta\rangle| \ge \rho\varepsilon r_\beta$, so the lower bound of
\ref{eq:cin-log} gives $\varphi(\langle\rho e_1,\eta\rangle) \ge \beta_1(\log\rho - L)$. Discarding
the rest of the sphere, $F_d(\rho e_1) \ge p\,\beta_1(\log\rho - L) \ge \beta(\log\rho - L)$ when
$\log\rho \ge L$, and the outer inequality holds as well when $\log\rho < L$, because $F_d \ge 0$.
Hence $|\widehat\mu(\xi)| \le e^{\beta L}\rho^{-\beta}$ for every $\xi \ne 0$. For $\rho \ge 1$,
$1 + \rho \le 2\rho$ gives $\rho^{-\beta} \le 2^\beta(1+\rho)^{-\beta}$; for $\rho < 1$,
$|\widehat\mu(\xi)| \le 1 \le 2^\beta(1+\rho)^{-\beta}$. This is the stated bound with
$C_\beta = 2^\beta(1 + e^{\beta L})$.

\emph{Clause (2).} Let $c > d+N$ and choose a finite $\beta$ with $d + N < \beta < c$. By clause
(1),
\[
  \int_{\RR^d}\bigl|\widehat\mu(\xi)\bigr|\,|\xi|^N\,d\xi
   \le C_\beta\int_{\RR^d}(1+|\xi|)^{-\beta}|\xi|^N\,d\xi < \infty .
\]
The integral converges at infinity because $\beta - N > d$, and at the origin because $N \ge 0$.
That is the hypothesis of \ledger{A1}, whose conclusion is the clause: a density $g \ge 0$ of
class $C^N$ whose partial derivatives of orders $0,\dots,N$ tend to $0$ at infinity.
At $N = 0$ the density is continuous and tends to $0$ at infinity, so it is bounded: it is at most
$1$ outside a compact set, and continuous on that set. For $c = \infty$ every $N$ is admissible,
and the densities obtained for different $N$ coincide: two continuous densities of $\mu$ agree
almost everywhere, hence everywhere, since a nonempty open set has positive Lebesgue measure. So the
density at $N = 0$ is of class $C^N$ for every $N$, that is, of class $C^\infty$.

\emph{Clause (3).} First, $c = 2\,k_{\mathrm I}(0{+})$. As $r \downarrow 0$ the integrand of
\ref{eq:polar-profile-d} increases to $v^{d/2-1}e^{-v}k_{\mathrm I}(0{+})$, because $k_{\mathrm I}$
is nonincreasing. So monotone convergence gives
$k_d(0{+}) = 2\,\Gamma(d/2)^{-1}\Gamma(d/2)\,k_{\mathrm I}(0{+})$. Next, \prog{B}{Prop.~4.9(2)},
restated, which is Lemma~\ref{lem:delay-polar-profile} for the delay datum $(0,k_{\mathrm I})$: the datum's exponent is $F_d(\xi) = F_{\mathrm I}(|\xi|^2/2)$. The law of $W$ has Laplace
transform $e^{-F_{\mathrm I}}$ by hypothesis. So
$\widehat\mu(\xi) = e^{-F_{\mathrm I}(|\xi|^2/2)} = \int_0^\infty e^{-u|\xi|^2/2}\Pr(W \in du)$,
which is the transform of the law of $\sqrt{W}Z$, and $\mu$ is that law. Since
$k_{\mathrm I}(0{+}) = c/2 > 0$, clause (1) read along a ray, at $|\xi| = \sqrt{2\sigma}$, gives
$e^{-F_{\mathrm I}(\sigma)} \to 0$, so
$\Pr(W = 0) = \lim_{\sigma\to\infty}e^{-F_{\mathrm I}(\sigma)} = 0$. So $W > 0$ almost surely,
and conditioning on $W$ gives \ref{eq:mixture-density}. For $x \ne 0$ the integrand is bounded in
$u$, so $p(x) < \infty$. Each integrand is radial and nonincreasing in $|x|$, and so is $p$.
Monotone convergence as $x \to 0$ gives
\[
  p_0 := \lim_{x\to0}p(x) = \int_0^\infty (2\pi u)^{-d/2}\Pr(W \in du) \in (0,\infty] .
\]
The two other forms of $p_0$ in the clause follow from Tonelli. Since $\widehat\mu \ge 0$, its
integral is defined in $[0,\infty]$, and $\int_{\RR^d}e^{-u|\xi|^2/2}d\xi = (2\pi/u)^{d/2}$ gives
\[
  (2\pi)^{-d}\int_{\RR^d}\widehat\mu(\xi)\,d\xi
   = (2\pi)^{-d}\int_0^\infty\Bigl(\frac{2\pi}{u}\Bigr)^{d/2}\Pr(W \in du) = p_0 .
\]
And $\int_0^\infty\sigma^{d/2-1}e^{-\sigma u}d\sigma = \Gamma(d/2)u^{-d/2}$ gives
\[
  \int_0^\infty\sigma^{d/2-1}e^{-F_{\mathrm I}(\sigma)}\,d\sigma
   = \Gamma(d/2)\int_0^\infty u^{-d/2}\Pr(W \in du) = \Gamma(d/2)\,(2\pi)^{d/2}\,p_0 .
\]

It remains to decide when $p_0$ is finite, and the argument reads the transform $\widehat\mu$
itself. If $c > d$, pick a finite $\beta \in (d,c)$. By clause (1),
$\widehat\mu(\xi) \le C_\beta(1+|\xi|)^{-\beta}$, which is integrable over $\RR^d$ because
$\beta > d$; so $p_0 < \infty$. If $c \le d$, then $c$ is finite, and
Lemma~\ref{lem:isotropic-exponent-log}(1) gives $\widehat\mu(\xi) \ge C_0|\xi|^{-c}$ for $|\xi| \ge 1$,
with $C_0 = e^{-13c/4 - 2C_k}$. In polar coordinates
$\int_{\RR^d}\widehat\mu \ge C_0|S^{d-1}|\int_1^\infty\rho^{d-1-c}\,d\rho = \infty$, because
$d - 1 - c \ge -1$. So $p_0 = \infty$, the threshold $c = d$ included.

That makes \emph{every} representative unbounded near the origin, by comparing measures rather than
functions, as in \restated{B}{Cor.~3.9}{cor:line-threshold}. If a representative $q$ of $\mu$'s density had $q \le M$ on
$\{0 < |x| < \delta\}$, then $\mu(B) = \int_B q \le M|B|$ for every Borel
$B \subseteq \{|x| < \delta\}$, the origin being a Lebesgue-null set, and
Lemma~\ref{lem:positive-transform-origin}(2) would give $\int_{\RR^d}\widehat\mu \le (2\pi)^dM < \infty$,
a contradiction.

\emph{Clause (4).} By Definition~\ref{def:isotropic-origin-datum}, $\widehat\mu = e^{-F_d}$ with $F_d$ real
and finite, so $\widehat\mu > 0$ everywhere and Lemma~\ref{lem:positive-transform-origin} applies. By
Lemma~\ref{lem:isotropic-exponent-log}(1), $\widehat\mu(\xi) \ge C_0|\xi|^{-c}$ for $|\xi| \ge 1$, with
$C_0 = e^{-13c/4 - 2C_k} > 0$ and $C_k$ as in that lemma. In polar coordinates,
\[
  \int_{\RR^d}\widehat\mu(\xi)\,d\xi \ge C_0\,|S^{d-1}|\int_1^\infty \rho^{d-1-c}\,d\rho = \infty ,
\]
because $d - 1 - c \ge -1$. If there were $\delta$ and $M$ with $\mu(B) \le M|B|$ for every Borel
$B \subseteq \{|x| < \delta\}$, Lemma~\ref{lem:positive-transform-origin}(2) would make this integral at
most $(2\pi)^dM$. So there are none. A representative $q$ of a density with $q \le M$ on
$\{0 < |x| < \delta\}$ would give such $\delta$ and $M$, by the last sentence of
Lemma~\ref{lem:positive-transform-origin}(2). So there is none.

For the rate, Lemma~\ref{lem:positive-transform-origin}(1) and the same lower bound give, for
$0 < t \le 1$,
\[
  \int g_t\,d\mu \ge (2\pi)^{-d}C_0\,|S^{d-1}|\int_1^\infty \rho^{d-1-c}e^{-t\rho^2/2}\,d\rho .
\]
If $c < d$, the substitution $\rho = \sigma/\sqrt t$ turns the last integral into
$t^{-(d-c)/2}\int_{\sqrt t}^\infty\sigma^{d-1-c}e^{-\sigma^2/2}\,d\sigma \ge
t^{-(d-c)/2}\int_1^\infty\sigma^{d-1-c}e^{-\sigma^2/2}\,d\sigma$, a positive multiple of
$t^{-(d-c)/2}$. If $c = d$, it is at least
$e^{-1/2}\int_1^{t^{-1/2}}\rho^{-1}\,d\rho = \tfrac12 e^{-1/2}\log(1/t)$, since
$t\rho^2 \le 1$ on that range. For the upper bound let $0 < \beta < c$. Clause (1) and
$(1+|\xi|)^{-\beta} \le |\xi|^{-\beta}$ give, with $\xi = y/\sqrt t$,
\[
  \int g_t\,d\mu \le (2\pi)^{-d}C_\beta\int_{\RR^d}|\xi|^{-\beta}e^{-t|\xi|^2/2}\,d\xi
   = (2\pi)^{-d}C_\beta\,t^{-(d-\beta)/2}\int_{\RR^d}|y|^{-\beta}e^{-|y|^2/2}\,dy ,
\]
and the last integral is finite because $\beta < c \le d$. Now take logarithms and divide by
$\log(1/t)$. By the lower bound, the quotient $\log\int g_t\,d\mu/\log(1/t)$ is at least
$(d-c)/2 - \varepsilon$ for every $\varepsilon > 0$ once $t$ is small enough. At $c = d$ this
uses $\log\log(1/t)/\log(1/t) \to 0$. By the upper bound, the quotient is at most
$(d-\beta)/2 + \varepsilon$ for every $\beta < c$ and $\varepsilon > 0$ once $t$ is small enough. So it
tends to $(d-c)/2$.

\emph{Clause (5).} If $c > d$, clause (2) at $N = 0$ gives $\mu$ a continuous density that tends
to $0$ at infinity, hence a bounded one. If $c \le d$, clause (4) excludes every bound
$\mu(B) \le M|B|$ on the Borel subsets $B$ of a ball $\{|x| < \delta\}$, and with it every
representative of a density bounded near the origin.
\end{proof}

% Q-0409: the proof of record of cor:origin-profile-exponent, transcribed (records/affine/
% edits_q0409_paper.py): Theorem thm:origin-boundedness-d(1) for the blueprint's split child
% lem:origin-decay-exponent, and without the proof of record's last sentence on what is
% machine-checked, which is section 1.1's. Pinned.
\begin{proof}[Proof of Corollary~\ref{cor:origin-profile-exponent}]
Let $\nu_d$ be the datum's \Levy{} measure and put
$G(\xi) := \int_{\RR^d}(1 - \cos\langle\xi,x\rangle)\,\nu_d(dx)$, so that
$F_d(\xi) = \tfrac12 a|\xi|^2 + G(\xi)$ by \ref{eq:isotropic-exponent}. Then $G \ge 0$, and $G$ is the
exponent of the datum $(0,k_d)$, which has the same polar profile and the same profile at the origin
$c$. As in the proof of Lemma~\ref{lem:isotropic-exponent-log}, \ref{eq:polar-decomposition} gives
\[
  G(\xi) = \frac{1}{|S^{d-1}|}\int_{S^{d-1}}\varphi\bigl(\langle\xi,\eta\rangle\bigr)\,\sigma_{d-1}(d\eta),
  \qquad
  \varphi(s) := \int_0^\infty\bigl(1 - \cos(rs)\bigr)\,k_d(r)\,\frac{dr}{r} .
\]

\emph{Clause (1).} Fix $\tau \in (0,1]$ and put $I(\tau) := \int_0^\tau k_d(r)\,r\,dr$ and
$T(\tau) := \int_\tau^\infty k_d(r)\,dr/r$. Both are finite: $I(\tau)$ is at most the first integral
of Definition~\ref{def:isotropic-origin-datum}, and
$T(\tau) \le k_d(\tau)\log(1/\tau) + \int_1^\infty k_d(r)\,dr/r$ because $k_d$ is nonincreasing.
Since $1 - \cos(rs) \le r^2s^2/2$ on $(0,\tau)$ and $1 - \cos \le 2$ on $[\tau,\infty)$,
$\varphi(s) \le \tfrac12 s^2\,I(\tau) + 2\,T(\tau)$. On the sphere
$|\langle\xi,\eta\rangle| \le |\xi|$, and the normalized surface measure has mass $1$. So
\[
  0 \le G(\xi) \le \tfrac12|\xi|^2\,I(\tau) + 2\,T(\tau), \qquad \xi \in \RR^d .
\]
Let $\varepsilon > 0$. The integrand of $I(1)$ is integrable on $(0,1)$, so $I(\tau) \to 0$ as
$\tau \downarrow 0$, and there is a $\tau$ with $I(\tau) \le \varepsilon$. For that $\tau$,
$G(\xi)/|\xi|^2 \le \tfrac12\varepsilon + 2T(\tau)/|\xi|^2 \le \varepsilon$ once
$|\xi|^2 \ge 4T(\tau)/\varepsilon$. Hence $G(\xi)/|\xi|^2 \to 0$, and
$F_d(\xi)/|\xi|^2 = \tfrac12 a + G(\xi)/|\xi|^2 \to \tfrac12 a$.

\emph{Clause (2).} The quotient is $G(\xi)/\log|\xi|$, and it is nonnegative for $|\xi| > 1$.
\emph{From below.} Let $\beta$ be finite with $0 < \beta < c$. By Theorem~\ref{thm:origin-boundedness-d}(1)
for the datum $(0,k_d)$ there is $C_\beta$ with $e^{-G(\xi)} \le C_\beta(1+|\xi|)^{-\beta}$ for
every $\xi$, and $C_\beta \ge 1$ by the bound at $\xi = 0$. That clause states the bound for
$|\widehat\mu|$, $\mu$ being a law of the datum $(0,k_d)$. Its proof bounds $e^{-G}$, which is
$|\widehat\mu|$ by the definition of the law, so the bound holds whether or not a law is given.
Taking logarithms, $G(\xi) \ge \beta\log|\xi| - \log C_\beta$ for $|\xi| \ge 1$. \emph{From above.}
Let $c < \infty$. By Lemma~\ref{lem:isotropic-exponent-log}(1) for the datum $(0,k_d)$,
$G(\xi) \le c\log|\xi| + \tfrac{13}{4}c + 2C_k$ for $|\xi| \ge 1$, with
$C_k = \int_1^\infty k_d(r)\,dr/r$. \emph{The limit.} Divide by $\log|\xi|$, which tends to
infinity. For every $\varepsilon > 0$ the quotient is at least $\beta - \varepsilon$ once $|\xi|$ is
large enough, for every finite $\beta$ with $0 < \beta < c$, and at most $c + \varepsilon$ once
$|\xi|$ is large enough when $c < \infty$. If $c = \infty$ the first bound makes the quotient tend to
$+\infty$. If $0 < c < \infty$ the two bounds give the limit $c$. If $c = 0$ the quotient lies
between $0$ and $\varepsilon$, and tends to $0$. For the last sentence of the clause, $a = 0$ gives
$F_d = G$, and $|\widehat\mu| = e^{-F_d}$ by the definition of the law.

\emph{Clause (3).} Since $d, d' \ge 1$ there are unit vectors $e \in \RR^d$ and $e' \in \RR^{d'}$,
and $F_d(re) = F(r) = F'_{d'}(re')$ for every $r > 0$. Clause (1) for the two data, read along
$\xi = re$ and $\xi' = re'$, gives $F(r)/r^2 \to \tfrac12 a$ and $F(r)/r^2 \to \tfrac12 a'$ as
$r \to \infty$. So $a' = a$. Clause (2), read along the same rays, then gives
$\bigl(F(r) - \tfrac12 a r^2\bigr)/\log r \to k_d(0{+})$ and
$\bigl(F(r) - \tfrac12 a r^2\bigr)/\log r \to k'_{d'}(0{+})$ as $r \to \infty$. A limit in
$[0,\infty]$ is unique, so $k'_{d'}(0{+}) = k_d(0{+})$.
\end{proof}

\begin{proof}[Proof of Proposition~\ref{prop:origin-lower-bound-d}]
Write $|E|$ for the Lebesgue measure of $E$, $B(z,r)$ and $\bar B(z,r)$ for the open and the closed
ball, and $\nu_d$ for the datum's \Levy{} measure. As $k_d$ is nonincreasing with right limit $c$ at
the origin, $k_d \le c$ on $(0,\infty)$ and $k_d(r) \uparrow c$ as $r \downarrow 0$, so
$\varepsilon_0$ exists. Every integral of $\nu_d$ below is computed by \ref{eq:polar-decomposition},
applied to a nonnegative Borel function vanishing at the origin.

\emph{The decomposition at scale $\varepsilon$.} Fix $\varepsilon \in (0,\varepsilon_0]$ and let
$\nu_1$ and $\nu_{\ge}$ be the restrictions of $\nu_d$ to $\{0 < |x| < \tau\varepsilon\}$ and to
$\{|x| \ge \tau\varepsilon\}$. Each has a radial density, so each is symmetric, and
$\psi_1(\xi) := \int(1 - \cos\langle\xi,x\rangle)\,\nu_1(dx)$ and
$\psi_{\ge}(\xi) := \int(1 - \cos\langle\xi,x\rangle)\,\nu_{\ge}(dx)$ are finite with
$F_d = \psi_1 + \psi_{\ge}$, $F_d$ being \ref{eq:isotropic-exponent} at $a = 0$. The mass of
$\nu_{\ge}$ is finite: as $\tau\varepsilon \le 1$ and $k_d \le c$,
\[
  m := \nu_{\ge}(\RR^d) = \int_{\tau\varepsilon}^{\infty}k_d(r)\,\frac{dr}{r}
   \le c\log\frac{1}{\tau\varepsilon} + C_k ,
  \qquad\text{so}\qquad e^{-m} \ge \tau^c\varepsilon^c e^{-C_k} .
\]
For a finite symmetric measure $\lambda$ of mass $m_\lambda$ put
$\pi_\lambda := e^{-m_\lambda}\sum_{n\ge0}\lambda^{*n}/n!$, with $\lambda^{*0} = \delta_0$; it is a
probability measure, and, summing the series of transforms by dominated convergence
($|\widehat\lambda^n| \le m_\lambda^n$), its transform is
$\exp(\widehat\lambda(\xi) - m_\lambda) = \exp\bigl(-\int(1 - \cos\langle\xi,x\rangle)\,\lambda(dx)\bigr)$,
the sine part of $\widehat\lambda$ vanishing by symmetry. Put $\pi := \pi_{\nu_{\ge}}$, of transform
$e^{-\psi_{\ge}}$. The profile $k_d\mathbf 1_{(0,\tau\varepsilon)}$ is nonincreasing and satisfies the
two integrability conditions of Definition~\ref{def:isotropic-origin-datum}, so
$(0, k_d\mathbf 1_{(0,\tau\varepsilon)})$ is an isotropic origin datum, with \Levy{} measure $\nu_1$
and exponent $\psi_1$; let $\mu_1$ be its law (\ledger{A3}(iii)), so $\widehat\mu_1 = e^{-\psi_1}$.
Then $\mu_1 * \pi$ has transform $e^{-F_d} = \widehat\mu$, and by the uniqueness of the Fourier
transform
\[
  \mu = \mu_1 * \pi .
\]
Keeping the term $n = 2$ of the series, $\pi \ge \tfrac12 e^{-m}\,\nu_{\ge} * \nu_{\ge}$, as measures.

\emph{The twofold convolution fills a ball.} Let $S := \{\varepsilon/2 \le |x| < \varepsilon\}$, which
lies in $\{|x| \ge \tau\varepsilon\}$ because $\tau \le 1/2$. On $S$ the density
$k_d(|x|)\,(|S^{d-1}||x|^d)^{-1}$ of $\nu_{\ge}$ is at least
$a := c\,(2|S^{d-1}|\varepsilon^d)^{-1}$, because $k_d(|x|) \ge k_d(\varepsilon_0) \ge c/2$ and
$|x|^d < \varepsilon^d$; so $\nu_{\ge}(A) \ge a\,|A \cap S|$ for every Borel $A$. Let
$|z| \le 3\varepsilon/2$. As $d \ge 2$ there is a unit vector $w$ orthogonal to $z$; put
$y_0 := z/2 + sw$ with $s := (9\varepsilon^2/16 - |z|^2/4)^{1/2}$, so that
$|y_0| = |z - y_0| = 3\varepsilon/4$. Every $y \in B(y_0,\varepsilon/8)$ has $|y|$ and $|z - y|$ in
$(5\varepsilon/8, 7\varepsilon/8)$, inside $S$, so
$|\{y \in S : z - y \in S\}| \ge |B(0,\varepsilon/8)|$. For a Borel $G \subseteq \bar B(0,3\varepsilon/2)$,
Tonelli's theorem and the substitution $x = z - y$ give
\[
\begin{aligned}
  (\nu_{\ge} * \nu_{\ge})(G) &\ge a^2\iint\mathbf 1_S(x)\,\mathbf 1_S(y)\,\mathbf 1_G(x + y)\,dx\,dy \\
   &= a^2\int_G\bigl|\{y \in S : z - y \in S\}\bigr|\,dz
   \ge a^2\,\bigl|B(0,\varepsilon/8)\bigr|\,|G| = \kappa_0\,\varepsilon^{-d}|G| .
\end{aligned}
\]
Hence $\pi(G) \ge \tfrac12 e^{-m}\kappa_0\varepsilon^{-d}|G|$ for every Borel
$G \subseteq \bar B(0, 3\varepsilon/2)$.

\emph{The small jumps move little.} By $1 - \cos u \le u^2/2$ and \ref{eq:polar-decomposition},
\[
  \psi_1(\xi) \le \frac{|\xi|^2}{2}\int|x|^2\,\nu_1(dx)
   = \frac{|\xi|^2}{2}\int_0^{\tau\varepsilon}k_d(r)\,r\,dr
   \le \frac{c\,\tau^2\varepsilon^2}{4}\,|\xi|^2 .
\]
Let $\gamma$ be the centred Gaussian law on $\RR^d$ with covariance $\varepsilon^{-2}I$, whose
transform at $w$ is $e^{-|w|^2/2\varepsilon^2}$ and whose second moment is
$\int|\xi|^2\,\gamma(d\xi) = d\,\varepsilon^{-2}$. By Fubini's theorem, $\widehat\mu_1$ being real,
and by $1 - e^{-\psi} \le \psi$,
\[
  \int\bigl(1 - e^{-|w|^2/2\varepsilon^2}\bigr)\,\mu_1(dw)
   = \int\bigl(1 - \widehat\mu_1(\xi)\bigr)\,\gamma(d\xi)
   \le \int\psi_1\,d\gamma
   \le \frac{c\,\tau^2\varepsilon^2}{4}\cdot\frac{d}{\varepsilon^2} = \frac{c\,\tau^2 d}{4} .
\]
The integrand on the left is at least $1 - e^{-1/2} \ge 1/3$ where $|w| \ge \varepsilon$, and
$c\tau^2 d/4 \le 3/(16d) \le 3/32$ by the choice of $\tau$ and $d \ge 2$, so
\[
  \mu_1\bigl(\{|w| \ge \varepsilon\}\bigr) \le \tfrac{9}{32} \le \tfrac12 .
\]

\emph{One scale.} Let $E \subseteq \bar B(0,\varepsilon/2)$ be Borel. For $|w| < \varepsilon$ the
translate $E - w$ lies in $\bar B(0,3\varepsilon/2)$ and has measure $|E|$. So, the integrand being
nonnegative,
\[
  \mu(E) = \int\pi(E - w)\,\mu_1(dw)
   \ge \int_{\{|w|<\varepsilon\}}\pi(E - w)\,\mu_1(dw)
   \ge \tfrac12 e^{-m}\kappa_0\varepsilon^{-d}|E|\cdot\tfrac12 .
\]
With $e^{-m} \ge \tau^c\varepsilon^c e^{-C_k}$ this is
\[
  \mu(E) \ge \kappa_1\,\varepsilon^{c-d}\,|E| ,\qquad
  \kappa_1 := \tfrac14\kappa_0\tau^ce^{-C_k},\qquad
  E \subseteq \bar B(0,\varepsilon/2),\ 0 < \varepsilon \le \varepsilon_0 .
\]

\emph{All scales.} For $n \ge 0$ put $\varepsilon_n := 2^{-n}\varepsilon_0$ and
$A_n := \{\varepsilon_n/4 < |x| \le \varepsilon_n/2\}$; these sets are disjoint and cover
$B(0,\delta)\setminus\{0\}$. For $x \in A_n$, $2|x| \le \varepsilon_n < 4|x|$, so
$\varepsilon_n^{c-d} \ge \theta|x|^{c-d}$ with $\theta := \min\{2^{c-d},4^{c-d}\}$ (the first bound
when $c \ge d$, the second when $c < d$). For a Borel $E \subseteq B(0,\delta)$, the origin being a
Lebesgue-null set, the previous step at $\varepsilon = \varepsilon_n$ gives
\[
  \mu(E) \ge \sum_{n\ge0}\mu(E \cap A_n)
   \ge \kappa_1\sum_{n\ge0}\varepsilon_n^{c-d}|E \cap A_n|
   \ge \kappa_1\theta\sum_{n\ge0}\int_{E\cap A_n}|x|^{c-d}\,dx
   = c_1\int_E|x|^{c-d}\,dx ,
\]
which is \ref{eq:origin-lower-bound}.

\emph{The consequences.} Write $\mu = \mu_{\mathrm{ac}} + \mu_{\mathrm s}$ for the Lebesgue
decomposition and let $N$ be a Borel set with $\mu_{\mathrm s}(N) = 0$ and
$|\RR^d \setminus N| = 0$. For a Borel $E \subseteq B(0,\delta)$,
$\int_E q\,dx \ge \int_{E\cap N} q\,dx = \mu_{\mathrm{ac}}(E \cap N) = \mu(E \cap N)
\ge c_1\int_{E\cap N}|x|^{c-d}\,dx = c_1\int_E|x|^{c-d}\,dx$. A function whose integral over every
Borel subset of $B(0,\delta)$ is at least that of the measurable $c_1|x|^{c-d}$ is at least
$c_1|x|^{c-d}$ almost everywhere on $B(0,\delta)$, both functions being nonnegative and measurable
and Lebesgue measure being $\sigma$-finite: on a set $E \subseteq B(0,\delta)$ of finite measure
where $q + 1/n \le c_1|x|^{c-d} \le n$, the inequality gives
$\int_E q\,dx + |E|/n \le \int_E q\,dx$, the integral being finite as $q \le n$ on $E$, so $|E| = 0$, and such sets exhaust the set
where $q < c_1|x|^{c-d}$. That is clause (1); for a density $q$ of $\mu$ the
same argument runs with $\int_E q\,dx = \mu(E)$, and $c_1|x|^{c-d} > 0$ off the origin, which is
null. As $x \to 0$, $|x| \to 0$ through positive values, and $|x|^{c-d} \to \infty$ when $c < d$.
Clause (2) is \ref{eq:origin-lower-bound} at $E = \{|x| < r\}$, in polar coordinates:
$\int_{\{|x|<r\}}|x|^{c-d}dx = |S^{d-1}|\int_0^r\rho^{c-1}d\rho = |S^{d-1}|\,r^c/c$.
% Q-0377: the proof of record's last paragraph, "What is machine-checked", is not printed, as Q-0344
% decided for lem:kernels-to-operators: section 1.1 owns what the formal proofs rest on. The rest of
% the proof is the proof of record verbatim (the marker's + pin records the difference).
\end{proof}

The last paragraph of the next proof shows that the uniform laws of C5 form a Borel set. Its
argument is a standard one in the space of the compact convex sets with the Hausdorff metric, by
Blaschke's selection theorem and the support function, and the facts are cited there from
\sscite{schneider2014convex}.

\begin{proof}[Proof of Proposition~\ref{prop:isotropic-radial-monotone}]
The argument reads four properties of a law $\mu$ on $\RR^d$: it is rotation invariant, it is
self-decomposable, its transform is real, and $\widehat\mu(\xi) \to 0$ as $|\xi| \to \infty$. The
first two paragraphs check them for the two clauses. The rest derives the density from them.

\emph{Clause (1): the four properties.} Put $\mu := \mu_{0,t}$. By
Theorem~\ref{thm:isotropic-classification}, $\widehat\mu(\xi) = e^{-F(\chi(t)|\xi|)}$ for a gauge $\chi$
and an $F \in \mathcal{A}_d$ that does not vanish identically on $[0,\infty)$, and
$s := \chi(t) > 0$. By the definition of $\mathcal{A}_d$ there is a self-decomposable law $\mu'$
on $\RR^d$ with $\widehat{\mu'}(\xi) = e^{-F(|\xi|)}$, so $\widehat\mu(\xi) = \widehat{\mu'}(s\xi)$.
So $\widehat\mu$ is real and radial: for $Q \in \Orth(d)$,
$\widehat{Q_*\mu}(\xi) = \widehat\mu(Q^{\tp}\xi) = \widehat\mu(\xi)$, and a law is determined by its
transform, so $\mu$ is rotation invariant. For $\lambda \in (0,1)$ let $\rho_\lambda$ be a residual
of $\mu'$, $\widehat{\mu'}(\xi) = \widehat{\mu'}(\lambda\xi)\,\widehat\rho_\lambda(\xi)$. Then
\[
  \widehat\mu(\xi) = \widehat{\mu'}(s\xi) = \widehat{\mu'}(\lambda s\xi)\,\widehat\rho_\lambda(s\xi)
   = \widehat\mu(\lambda\xi)\,\widehat{(sI)_*\rho_\lambda}(\xi) ,
\]
which is \ref{eq:self-decomposable} with residual $(sI)_*\rho_\lambda$. It remains to see that
$e^{-F(r)} \to 0$ as $r \to \infty$. $F$ is admissible, of the form \ref{eq:admissible-line} with
Gaussian coefficient $a_{\mathrm V}$ and folded profile $k$. The pair $(2a_{\mathrm V}, k)$ is an
isotropic origin datum in dimension $1$: its two integrability conditions are those of
\ref{eq:admissible-line}, and by \ref{eq:polar-decomposition} at $d = 1$, where $|S^0| = 2$, its
exponent at $\omega$ is
$a_{\mathrm V}\omega^2 + \int_{\RR}(1 - \cos\omega x)\,k(|x|)\,(2|x|)^{-1}dx = F(|\omega|)$. Its
law is the image of $\mu'$ under $x \mapsto \langle e_1, x\rangle$, whose transform is
$\widehat{\mu'}(\omega e_1) = e^{-F(|\omega|)}$, as in Lemma~\ref{lem:radial-marginal}(1). The datum is
not trivial: a trivial datum has exponent $0$, and $F$ does not vanish identically on $[0,\infty)$.
So the last two sentences of the next paragraph, at $d = 1$, give $e^{-F(r)} \to 0$, and
$\widehat\mu(\xi) \to 0$ as $|\xi| \to \infty$.

\emph{Clause (2): the four properties.} Let $\mu$ be the datum's law, so that
$\widehat\mu = e^{-F_d}$ is real, positive and radial by \ref{eq:isotropic-exponent}. For
$Q \in \Orth(d)$, $\widehat{Q_*\mu}(\xi) = \widehat\mu(Q^{\tp}\xi) = \widehat\mu(\xi)$, and a law is
determined by its transform, so $\mu$ is rotation invariant. Put $F(r) := F_d(re_1)$. Then
$\xi \mapsto F(|\xi|)$ is the exponent $F_d$ of the datum, so by
Proposition~\ref{prop:isotropic-polar-reading}(1), in its converse direction, it is the exponent of a
self-decomposable law, which is $\mu$. Finally $F_d(\xi) \ge \tfrac12 a|\xi|^2$, which tends to
infinity if $a > 0$. If $a = 0$, then $k_d$ is not almost everywhere $0$, so
$c = k_d(0{+}) > 0$, $k_d$ being at most $c$ on $(0,\infty)$, and
Theorem~\ref{thm:origin-boundedness-d}(1) with a finite $\beta \in (0,c)$ gives
$\widehat\mu(\xi) \le C_\beta(1+|\xi|)^{-\beta}$.

\emph{No atom at the origin.} Let $\mu$ have the four properties. For $\varepsilon > 0$, multiply
the identity of Lemma~\ref{lem:positive-transform-origin}(1) at $t = \varepsilon$ by
$(2\pi\varepsilon)^{d/2}$ and substitute $\xi = y/\sqrt\varepsilon$:
\[
  \int_{\RR^d} e^{-|x|^2/2\varepsilon}\,\mu(dx)
   = (2\pi)^{-d/2}\int_{\RR^d}\widehat\mu\bigl(y/\sqrt\varepsilon\bigr)\,e^{-|y|^2/2}\,dy .
\]
Let $\varepsilon \downarrow 0$. On the left the integrand tends to the indicator of the origin, and
on the right $\widehat\mu(y/\sqrt\varepsilon) \to 0$ for every $y \ne 0$. The integrand on the left is bounded
by $1$, which is integrable against $\mu$, and the one on the right by $e^{-|y|^2/2}$, which is
integrable against $dy$, as $|\widehat\mu| \le 1$; so by dominated convergence $\mu(\{0\}) = 0$.

\emph{The mixture.} The law $\mu$ is symmetric, by rotation invariance at $Q = -I$, and
self-decomposable. By \ledger{A8} it is symmetric unimodal: there is a probability measure
$\Lambda$ on the uniform laws $u_K$ of the compact convex symmetric sets $K \subset \RR^d$ with
\[
  \mu(A) = \int u_K(A)\,\Lambda(du_K) \qquad \text{for every Borel } A \subseteq \RR^d .
\]
Here $u_K$ is uniform for the Lebesgue measure $\mathrm{Vol}_V$ of the linear span $V$ of $K$, whose
dimension $k$ may be anything from $0$ to $d$, and $K$ is nonempty. $\Lambda$ is a probability
measure on the measurable space of the probability measures on $\RR^d$, whose $\sigma$-algebra is
generated by the evaluations $\nu \mapsto \nu(A)$, and almost every measure under it is a $u_K$;
nothing below asks the set $W_d$ of the $u_K$ to be measurable. The source puts its mixing measure
on $W_d$ itself, with the Borel sets of weak convergence, which on the probability measures of
$\RR^d$ are the sets of that $\sigma$-algebra. Reading it as almost every measure being a $u_K$
uses that $W_d$ is such a set. It is one, being a countable union of weakly compact sets, as the
paragraph on $W_d$ at the end of the proof shows. The formal statement puts
$\Lambda$ on all measures on $\RR^d$, which loses nothing: the probability measures form a
measurable set, and almost every measure under $\Lambda$, being a $u_K$, lies in it. A set with
$k < d$ carries no density in $\RR^d$, and the next two paragraphs deal with that.

\emph{The radial part of one uniform law.} Let $\omega_d$ be the volume of the unit ball and
$\varphi(x) := \omega_d|x|^d$ the volume of the ball of radius $|x|$. For $v \ge 0$ put
$r_v := (v/\omega_d)^{1/d}$ and $B(v) := \varphi^{-1}((0,v]) = \{0 < |x| \le r_v\}$, the
punctured closed ball of volume $v$, and put $G_K(v) := u_K(B(v))$. We show that $G_K$ is concave
on $[0,\infty)$. If $k = 0$, then $K = \{0\}$, $u_K = \delta_0$ and $G_K \equiv 0$, as $B(v)$
misses the origin. Let $k \ge 1$, let $S_V$ be the unit sphere of $V$ and $\sigma_V$ its surface
measure, so that $\mathrm{Vol}_V(A) = \int_{S_V}\sigma_V(d\theta)\int_0^\infty
\mathbf{1}_A(s\theta)\,s^{k-1}\,ds$ for every Borel $A \subseteq V \setminus \{0\}$ (polar
coordinates in $V$). Fix $\theta \in S_V$. The set $J_\theta := \{s > 0 : s\theta \in K\}$ contains
$(0, s]$ with each of its points $s$: $K$ contains $0 = \tfrac12 y + \tfrac12(-y)$ for any
$y \in K$, so $s\theta \in K$ and $0 < s' < s$ give
$s'\theta = \tfrac{s'}{s}\,s\theta + \bigl(1 - \tfrac{s'}{s}\bigr)0 \in K$ by convexity. It is
bounded, $K$ being compact. With $\rho_\theta := \sup J_\theta$ ($\rho_\theta := 0$ if
$J_\theta$ is empty), $(0,\rho_\theta) \subseteq J_\theta \subseteq (0,\rho_\theta]$, and for
$s > 0$, $s\theta \in B(v)$ exactly when $s \le r_v$. So the ray in the direction $\theta$
contributes
\[
  \int_0^\infty\mathbf{1}_{K \cap B(v)}(s\theta)\,s^{k-1}\,ds
   = \int_0^{\min\{\rho_\theta,\,r_v\}}s^{k-1}\,ds
   = \frac1k\min\Bigl\{\rho_\theta^k,\ \Bigl(\frac{v}{\omega_d}\Bigr)^{k/d}\Bigr\} .
\]
As $k \le d$, $v \mapsto (v/\omega_d)^{k/d}$ is concave on $[0,\infty)$, and so is its minimum with
a constant. Integrating over $\theta$ against $\sigma_V$ and dividing by
$\mathrm{Vol}_V(K) \in (0,\infty)$, $G_K(v) = \mathrm{Vol}_V(K \cap B(v))/\mathrm{Vol}_V(K)$ is
concave in $v$. The volume is positive because $K$ spans $V$, so it contains a basis
$v_1,\dots,v_k$ of $V$ and, being convex and symmetric, the convex hull of the $\pm v_i$, which has
positive volume in $V$; it is finite because $K$ is compact. (Concavity survives multiplication by
any nonnegative constant, so the formal proof does not need this; that $u_K$ is a probability
measure is checked separately.) In words: each ray of $K$ contributes the minimum of a constant and
a concave power of the volume, the power being $k/d \le 1$ because $K$ has dimension at most $d$.

\emph{The density.} Put $G(v) := \mu(B(v)) = \int G_K(v)\,\Lambda(du_K)$. A mixture of concave
functions is concave, the inequality $aG_K(v) + bG_K(w) \le G_K(av + bw)$ integrating against
$\Lambda$, so $G$ is concave on $[0,\infty)$. Let $m$ be the image of $\mu$ under $\varphi$, a
finite measure on $\RR$ with $m((0,v]) = G(v)$ and $m((-\infty,0]) = \mu(\{0\}) = 0$. $G$ is
nondecreasing; let $D(v)$ be its right derivative at $v > 0$. It is finite, nonnegative and
nonincreasing, and $G$ is continuous on $(0,\infty)$, so by the fundamental theorem of calculus
for right derivatives $m((a,b]) = G(b) - G(a) = \int_a^b D(w)\,dw$ for $0 < a \le b$. Letting
$a \downarrow 0$ by continuity from below, and as $m$ does not charge $(-\infty,0]$, the measures
$m$ and $D(w)\,dw$ on $(0,\infty)$ agree on every interval $(a,b]$, so $m = D(w)\,dw$ on
$(0,\infty)$. Put $q(r) := D(\omega_d r^d)$, a nonincreasing function on $(0,\infty)$ with values
in $[0,\infty)$, so that $q(|x|) = D(\varphi(x))$. The image of Lebesgue measure on $\RR^d$ under
$\varphi$ is Lebesgue measure on $(0,\infty)$, because $\varphi^{-1}((-\infty,v])$ is the closed
ball of volume $v$. So the image of $q(|x|)\,dx$ under $\varphi$ is $D(w)\,dw = m$, the image of
$\mu$, and as $|x| = (\varphi(x)/\omega_d)^{1/d}$ the two measures have the same image under
$x \mapsto |x|$. Both are finite and rotation invariant ($q(|x|)\,dx$ because Lebesgue measure is and
the density is radial), and both vanish at the origin. A measure with these properties is
determined by its image under $x \mapsto |x|$. For let $m$ be one of them and
$\varepsilon > 0$, and let $m_\varepsilon$ be its restriction to
$U_\varepsilon := \{|x| > \varepsilon\}$, a finite measure invariant under every $Q \in \Orth(d)$,
as $U_\varepsilon$ is. Then $\widehat{m_\varepsilon}(\xi) = \widehat{m_\varepsilon}(|\xi|\theta)$
for every unit vector $\theta$, a reflection carrying $\xi$ to $|\xi|\theta$. Averaging over
$\theta$ against $\sigma_{d-1}/|S^{d-1}|$ and exchanging the integrals gives
\[
  \widehat{m_\varepsilon}(\xi)
   = \frac{1}{|S^{d-1}|}\int_{U_\varepsilon}\Phi(|\xi|,|x|)\,m(dx) ,
  \qquad
  \Phi(s,r) := \int_{S^{d-1}} e^{i s r\langle e_1,\theta\rangle}\,\sigma_{d-1}(d\theta) ,
\]
since $\int e^{i s\langle x,\theta\rangle}\sigma_{d-1}(d\theta)$ depends on $x$ only through $|x|$,
by the invariance of $\sigma_{d-1}$. The right side depends on $m$ only through its image under
$x \mapsto |x|$. So the two measures have restrictions to $U_\varepsilon$ with the same transform,
hence equal; letting $\varepsilon \downarrow 0$ along $1/(n+1)$, they agree off the origin, and as
neither charges it, they are equal. Hence $\mu = q(|x|)\,dx$, and $p(x) := q(|x|)$ is a density of
$\mu$.

% Q-0424 (the external review, round 1; records/affine/reviews/VERIFY-round1-mathematics.md, item
% 6(a)): the paragraph below and the three sentences on the source's mixing measure in the
% paragraph on the mixture are the proof of record's, changed at the blueprint and mirrored here
% verbatim (the marker is re-pinned). They replace the assertion that W_d is analytic, hence
% universally measurable, which had no proof and no reference.
% Q-0430: citations only, changed at the blueprint and mirrored here verbatim (the marker is
% re-pinned): Schneider, Convex Bodies, 2nd ed., Thm. 1.8.7, p. 63, with p. 8; Lemma 1.8.14,
% p. 66; p. 45, each read at the page image (records/affine/READING-schneider2014convex.md).
\emph{The set $W_d$ is a Borel set.} This paragraph proves the reading of the source made in the
paragraph on the mixture, and no other step uses it. Let $\bar B$ be the closed unit ball of
$\RR^d$, and give the nonempty compact convex sets $K = -K$ of $\RR^d$ the Hausdorff metric. For
integers $k \ge 1$ and $n \ge 1$ let $\mathcal{K}_{k,n}$ be the set of those $K$ whose span $V$
has dimension $k$ and which satisfy
\[
  \tfrac1n\bar B \cap V \subseteq K \subseteq n\bar B .
\]
Every $K$ of dimension $k \ge 1$ is in $\mathcal{K}_{k,n}$ for some $n$. It is bounded, and it
contains the convex hull of the points $\pm v_i$ of the paragraph on one uniform law, which
contains a ball of $V$ round the origin.

Each $\mathcal{K}_{k,n}$ is compact. Let $K_j \in \mathcal{K}_{k,n}$, with spans $V_j$. By
Blaschke's selection theorem (\sscite{schneider2014convex}, Thm.~1.8.7, p.~63, where a convex
body is any nonempty compact convex set, p.~8), the nonempty compact convex subsets of a ball
form a compact space in the Hausdorff metric. So a subsequence of $(K_j)$ converges to a nonempty compact convex set
$K$, and $K = -K$. Let $P_j$ be the orthogonal projection onto $V_j$. Along a further subsequence
the matrices $P_j$ converge, and their limit $P$ is the orthogonal projection onto a subspace $V'$
of dimension $k$, since $P_j^2 = P_j = P_j^{\tp}$ and $\tr P_j = k$ pass to the limit. A point
$x \in K$ is a limit of points $x_j \in K_j$, and $x_j = P_jx_j \to Px$, so $x \in V'$. Hence
$K \subseteq V' \cap n\bar B$. For $y \in V'$ with $|y| \le 1/n$ the points $P_jy$ lie in
$\tfrac1n\bar B \cap V_j \subseteq K_j$ and converge to $y$, so $y \in K$. Hence the span of $K$
is $V'$, and $K \in \mathcal{K}_{k,n}$.

The map $K \mapsto u_K$ is continuous on $\mathcal{K}_{k,n}$, for the weak convergence of measures.
Let $K_j \to K$ in $\mathcal{K}_{k,n}$, with spans $V_j$ and $V$. Choose a basis $y_1,\dots,y_k$ of
$V$ in $K$ and points $y_{i,j} \in K_j$ with $y_{i,j} \to y_i$ as $j \to \infty$. For large $j$
the points $y_{1,j},\dots,y_{k,j}$ are linearly independent, their Gram determinant converging to
that of the $y_i$, which is positive. So they are a basis of $V_j$. Let $T_j$ be the linear map
from $V$ onto $V_j$ with $T_jy_i = y_{i,j}$, and let $\varepsilon_j$ be the supremum of $|T_jz - z|$
over the $z \in V$ with $|z| \le 1$. Then $\varepsilon_j \to 0$. A linear bijection carries the
uniform law of a set to the uniform law of its image, so $u_{K_j}$ is the image of $u_{K'_j}$ under
$T_j$, where $K'_j := T_j^{-1}K_j \subseteq V$. For $w \in K'_j$, $|T_jw - w| \le \varepsilon_j|w|$
and $|T_jw| \le n$, so $|w| \le n/(1-\varepsilon_j)$, and the Hausdorff distance of $K'_j$ and $K_j$
is at most $\varepsilon_jn/(1-\varepsilon_j)$. So $K'_j \to K$. Let $\eta_j$ be the Hausdorff
distance of $K'_j$ and $K$, and put $r := 1/n$ and $B_V := \bar B \cap V$, so that
$rB_V \subseteq K$. Then
\[
  K'_j \subseteq K + \eta_jB_V \subseteq \bigl(1 + \eta_j/r\bigr)K ,
\]
$K$ being convex. In the other direction, write $h_A(e) := \sup_{x \in A}\langle e,x\rangle$ for the
support function of a compact convex set $A \subseteq V$ at a unit vector $e$ of $V$. Then
$h_K \le h_{K'_j} + \eta_j$, the Hausdorff distance of two compact convex sets being the supremum
over the unit vectors of the absolute difference of their support functions
(\sscite{schneider2014convex}, Lemma~1.8.14, p.~66), and $h_K \ge r$, so
$h_{K'_j} \ge (1 - \eta_j/r)\,h_K$. A compact
convex set contains every compact convex set with a smaller support function, and
$h_{\lambda A} = \lambda h_A$ for $\lambda \ge 0$ (\sscite{schneider2014convex}, p.~45), so
$(1 - \eta_j/r)K \subseteq K'_j$ once $\eta_j < r$. The two inclusions hold for $K$ in place of
$K'_j$ as well. So $K'_j$ and $K$ differ by a set of volume at most
$\bigl((1+\eta_j/r)^k - (1-\eta_j/r)^k\bigr)\mathrm{Vol}_V(K)$, which tends to $0$, and
$u_{K'_j}(A) \to u_K(A)$ uniformly in the Borel set $A$. For a bounded continuous function
$\varphi$ on $\RR^d$,
\[
  \int\varphi\,du_{K_j} = \int\varphi(T_jw)\,u_{K'_j}(dw) ,
\]
and $\varphi(T_jw) \to \varphi(w)$ uniformly for $w$ in the bounded set $2K$, which contains $K'_j$
for large $j$. So $\int\varphi\,du_{K_j} \to \int\varphi\,du_K$.

So each set $\{u_K : K \in \mathcal{K}_{k,n}\}$ is the continuous image of a compact set. It is
weakly compact, hence weakly closed. With $\{\delta_0\}$, the case $k = 0$, these countably many
sets make up $W_d$. So $W_d$ is a Borel set for weak convergence.
% Q-0377: the proof of record's last paragraph, "What is machine-checked", is not printed, as Q-0344
% decided for lem:kernels-to-operators: section 1.1 owns what the formal proofs rest on. The rest of
% the proof is the proof of record verbatim (the marker's + pin records the difference).
\end{proof}

% Q-0415: the proof of record of prop:origin-upper-bound-d, transcribed (records/affine/
% edits_q0415_paper.py), after that of prop:isotropic-radial-monotone, which it reads, and
% without the proof of record's last paragraph on what is machine-checked, which is section
% 1.1's. Pinned.
\begin{proof}[Proof of Proposition~\ref{prop:origin-upper-bound-d}]
Write $B_r := \{|x| < r\}$ for the open ball and $|E|$ for the Lebesgue measure of $E$, so that
$|B_r| = V_d\,r^d$.

\emph{Clause (1).} Fix $\beta$ with $0 < \beta < c$, and let $C'_\beta$ be a constant with
$\int g_t\,d\mu \le C'_\beta\,t^{-(d-\beta)/2}$ for all $t > 0$, $g_t$ being the heat kernel of
Lemma~\ref{lem:positive-transform-origin}. There is one by Theorem~\ref{thm:origin-boundedness-d}(4), since
$c \le d$. Let $r > 0$. For $|x| < r$,
\[
  g_{r^2}(x) = (2\pi r^2)^{-d/2}e^{-|x|^2/2r^2} \ge (2\pi)^{-d/2}e^{-1/2}\,r^{-d} .
\]
As $g_{r^2} \ge 0$,
\[
  \int_{\RR^d} g_{r^2}\,d\mu \ge \int_{B_r} g_{r^2}\,d\mu
   \ge (2\pi)^{-d/2}e^{-1/2}\,r^{-d}\,\mu(B_r) .
\]
The left side is at most $C'_\beta\,(r^2)^{-(d-\beta)/2} = C'_\beta\,r^{\beta-d}$. Hence
$\mu(B_r) \le (2\pi)^{d/2}e^{1/2}\,C'_\beta\,r^{\beta}$, which is the bound with the stated
$A_\beta$.

\emph{Clause (2).} As $c = \sup_{r>0}k_d(r) > 0$, there is $r_0 > 0$ with $k_d(r_0) > 0$, and
$k_d \ge k_d(r_0)$ on $(0,r_0]$ because $k_d$ is nonincreasing. So $k_d$ is not almost everywhere
$0$, the datum is not trivial, and Proposition~\ref{prop:isotropic-radial-monotone}(2) gives a density
$p(x) = q(|x|)$ with $q$ nonincreasing. Now let $q : (0,\infty) \to [0,\infty)$ be any
nonincreasing function such that $x \mapsto q(|x|)$ is a density of $\mu$, and let $r > 0$. For
$0 < |x| < r$, $q(|x|) \ge q(r)$. The origin is a Lebesgue-null set, so
\[
  \mu(B_r) = \int_{B_r} q(|x|)\,dx \ge q(r)\,|B_r| = q(r)\,V_d\,r^d .
\]
With clause (1), $q(r) \le A_\beta V_d^{-1}\,r^{\beta-d}$.

\emph{Clause (3), the lower bound at every radius.} Let $d \ge 2$. Since $c \le d$, $c$ is finite,
and Proposition~\ref{prop:origin-lower-bound-d} applies. By its clause (1) the density $x \mapsto q(|x|)$ of
$\mu$ is at least $c_1|x|^{c-d}$ at almost every $x \in B_\delta$. Suppose that
$q(r_0) < c_1 r_0^{c-d}$ for some $r_0$ with $0 < r_0 < \delta$. The function
$r \mapsto c_1 r^{c-d}$ is continuous at $r_0$, so there is $r_1$ with $r_0 < r_1 \le \delta$ and
$q(r_0) < c_1 r^{c-d}$ for $r_0 < r < r_1$. For such $r$, $q(r) \le q(r_0) < c_1 r^{c-d}$,
because $q$ is nonincreasing. So $q(|x|) < c_1|x|^{c-d}$ on the shell $\{r_0 < |x| < r_1\}$. The
shell is a nonempty open subset of $B_\delta$, hence of positive Lebesgue measure, and this
contradicts the bound at almost every point. So $q(r) \ge c_1 r^{c-d}$ for every $0 < r < \delta$.

\emph{Clause (3), the limits.} Both follow from one remark. Let $L$ be real and let $f$ be a
function on $(0,\infty)$ such that, for small $r > 0$, $f(r) \ge C r^{-L}$ with a $C > 0$ and, for
every $\varepsilon > 0$, $f(r) \le C_\varepsilon r^{-(L+\varepsilon)}$ with a
$C_\varepsilon > 0$. Then $\log f(r)/\log(1/r) \to L$ as $r \downarrow 0$. Indeed, for
$0 < r < 1$ small, $f(r) > 0$ and $\log(1/r) > 0$, and taking logarithms gives
\[
  L + \frac{\log C}{\log(1/r)} \le \frac{\log f(r)}{\log(1/r)}
   \le L + \varepsilon + \frac{\log C_\varepsilon}{\log(1/r)} .
\]
As $\log(1/r) \to \infty$, the quotient is at least $L - \varepsilon$ and at most
$L + 2\varepsilon$ once $r$ is small enough, for every $\varepsilon > 0$.

For the density take $f = q$ and $L = d - c$. The lower bound is the one just proved,
$q(r) \ge c_1 r^{-(d-c)}$ for $0 < r < \delta$. For the upper bound let $\varepsilon > 0$ and
choose $\beta$ with $0 < \beta < c$ and $\beta \ge c - \varepsilon$. By clause (2), for
$0 < r < 1$,
$q(r) \le A_\beta V_d^{-1}\,r^{\beta-d} \le A_\beta V_d^{-1}\,r^{-(d-c+\varepsilon)}$, as
$\beta - d \ge -(d - c + \varepsilon)$ and $r < 1$. So $\log q(r)/\log(1/r) \to d - c$, and with
$r = |x|$ this is the first limit, since $|x| \downarrow 0$ through positive values as $x \to 0$.

For the mass of a ball take $f(r) = \mu(B_r)$ and $L = -c$. By
Proposition~\ref{prop:origin-lower-bound-d}(2), $\mu(B_r) \ge c_1|S^{d-1}|\,c^{-1}\,r^{c}$ for
$0 < r \le \delta$. By clause (1), with $\beta$ as before,
$\mu(B_r) \le A_\beta\,r^{\beta} \le A_\beta\,r^{c-\varepsilon}$ for $0 < r < 1$. So
$\log\mu(B_r)/\log(1/r) \to -c$, and $\log(1/r) = -\log r$ gives the second limit.
\end{proof}

% Q-0328, Q-0424: the closed form of the constant of the upper bound, computed in the text.
\paragraph{The constant of Proposition~\ref{prop:origin-upper-bound-d} in closed form.} This
paragraph computes a constant from the two proofs above. It is not part of the formal development
(\S\ref{sec:trust-base}). Let $0 < \beta < c \le d$, put $\beta_1 := (\beta + c)/2$, and let
$r_\beta \in (0,1]$ satisfy $k_d(r_\beta) \ge \beta_1$. There is such an $r_\beta$, because
$k_d(0{+}) = c > \beta_1$, and $k_d \ge \beta_1$ on $(0,r_\beta)$, because $k_d$ is nonincreasing.
So $\beta_1$ and $r_\beta$ can be used in the proof of clause (1) of
Theorem~\ref{thm:origin-boundedness-d}. That proof also asks for an $\varepsilon \in (0,1]$ for
which the cap $K_\varepsilon = \{\eta \in S^{d-1} : |\langle e_1,\eta\rangle| \ge \varepsilon\}$
has at least the share $\beta/\beta_1$ of the sphere. Take
\[
  \varepsilon := \frac{c - \beta}{(c + \beta)\sqrt d} .
\]
Then $1 - \beta/\beta_1 = \varepsilon\sqrt d$, so it is enough that the band
$\{|\langle e_1,\eta\rangle| < \varepsilon\}$ has at most the share $\varepsilon\sqrt d$. For
$d = 1$ the band is empty. For $d = 2$ its share is $(2/\pi)\arcsin\varepsilon \le \varepsilon$.
For $d \ge 3$ the first coordinate of a uniform point of the sphere has the density
$\Gamma(\tfrac d2)\,\bigl(\sqrt\pi\,\Gamma(\tfrac{d-1}2)\bigr)^{-1}(1 - u^2)^{(d-3)/2}$ on $(-1,1)$.
It is at most $\sqrt{(d-1)/(2\pi)}$, since $\Gamma(x + \tfrac12) \le \Gamma(x)\sqrt x$ by the
logarithmic convexity of $\Gamma$. So the share of the band is at most
$2\varepsilon\sqrt{(d-1)/(2\pi)} \le \varepsilon\sqrt d$. With these choices the number $L$ of
that proof satisfies $e^{L} = e^3/(\varepsilon r_\beta) = e^3\sqrt d\,(c+\beta)/((c-\beta)\,r_\beta)$,
and clause (1) holds with $C_\beta = 2^\beta(1 + e^{\beta L})$. The proof of clause (4) gives
\[
  C'_\beta = (2\pi)^{-d}\,C_\beta\int_{\RR^d}|y|^{-\beta}e^{-|y|^2/2}\,dy ,
  \qquad
  \int_{\RR^d}|y|^{-\beta}e^{-|y|^2/2}\,dy
   = |S^{d-1}|\,2^{(d-\beta)/2-1}\,\Gamma\bigl(\tfrac{d-\beta}{2}\bigr) ,
\]
the second in polar coordinates. With $|S^{d-1}| = 2\pi^{d/2}/\Gamma(\tfrac d2)$, the constant
$A_\beta = (2\pi)^{d/2}e^{1/2}\,C'_\beta$ of the proposition is
\[
  A_\beta = \sqrt e\;2^{-\beta/2}\,C_\beta\,
    \frac{\Gamma\bigl(\tfrac{d-\beta}{2}\bigr)}{\Gamma\bigl(\tfrac d2\bigr)}
  = \sqrt e\;2^{\beta/2}\,\frac{\Gamma\bigl(\tfrac{d-\beta}{2}\bigr)}{\Gamma\bigl(\tfrac d2\bigr)}
    \Bigl(1 + \Bigl(\frac{e^3\sqrt d\,(c + \beta)}{(c - \beta)\,r_\beta}\Bigr)^{\beta}\Bigr) .
\]

% Q-0448 (the author's decision of 2026-10-13: the rate at the centre is stated once in the
% paper, through the polar profile): the rate through the transform, Proposition 3.16 of the
% second review build, is a lemma of the supplement. The statement is the blueprint node
% prop:origin-rate-d, verbatim, and the label is the node's. Its attribution is in the table of
% the sources of the paper, once.
The rate through the polar profile (Proposition~\ref{prop:origin-rate-profile-d}) is proved
from the rate through the transform, which is the next lemma. For a law whose transform is a
power times a slowly varying factor its clause (1) is a known Tauberian theorem, and
\S\ref{sec:isotropic-sources} says where it is in print. For a finite $c$ the transform of a
kernel has that form, $|\xi|^{-c}L(|\xi|)$ with $L$ slowly varying
(Lemma~\ref{lem:isotropic-exponent-log}(2)). In the statement $s$ is a distance from the
centre, not a scale, $e$ is a unit vector, and $\widehat\mu(e/s)$ is the transform at a
frequency of length $1/s$.

% shared with blueprint prop:origin-rate-d+8e1d7317599c
\begin{lemma}[the rate at the centre through the transform]\label{prop:origin-rate-d}
Let $d \ge 1$, let $(0, k_d)$ be an isotropic origin datum in dimension $d$ with Gaussian
coefficient $0$, let $\mu$ be its law, and suppose that $c := k_d(0{+})$ satisfies $0 < c \le d$.
Let $e$ be a unit vector of $\RR^d$, let $|S^{d-1}|$ be the surface measure of the unit sphere and
$V_d$ the volume of the unit ball. The law $\mu$ has a density $p$ with $p(x) = q(|x|)$ for
$x \ne 0$, where $q : (0,\infty) \to [0,\infty)$ is nonincreasing. Let $q$ be any such function.
All limits are as $s \downarrow 0$.
\begin{enumerate}
\item \emph{(Below the threshold.)} Let $c < d$ and put
\[
  R(d,c) := \frac{\Gamma\bigl(\tfrac{d-c}{2}\bigr)}{2^{c}\,\pi^{d/2}\,\Gamma\bigl(\tfrac c2\bigr)} .
\]
Then
\[
  \frac{\mu(\{|x| \le s\})}{R(d,c)\,|S^{d-1}|\,c^{-1}\,\widehat\mu(e/s)} \to 1
  \qquad\text{and}\qquad
  \frac{q(s)}{R(d,c)\,s^{-d}\,\widehat\mu(e/s)} \to 1 .
\]
\item \emph{(At the threshold.)} Let $c = d$ and put
$I(s) := \int_{\{|\xi| \le 1/s\}}\widehat\mu(\xi)\,d\xi$. Then
\[
  \frac{\mu(\{|x| \le s\})}{(2\pi)^{-d}\,V_d\,s^{d}\,I(s)} \to 1
  \qquad\text{and}\qquad
  \frac{q(s)}{(2\pi)^{-d}\,I(s)} \to 1 .
\]
Moreover $I(s) \ge |S^{d-1}|\,e^{-13d/4 - 2C_k}\log(1/s)$ for $0 < s \le 1$, with
$C_k = \int_1^\infty k_d(r)\,dr/r$. So for every
$C < (2\pi)^{-d}\,|S^{d-1}|\,e^{-13d/4 - 2C_k}$ there is $s_0 > 0$ with
$q(s) \ge C\log(1/s)$ for $0 < s < s_0$.
\end{enumerate}
\end{lemma}

% A paragraph of paper prose stands between the statement and its proof: `linkage check` reads a
% proof that directly follows a statement as that statement's proof, heading included.
So for $c < d$ the density $p(x)$ is asymptotic, as $x \to 0$, to
$R(d,c)\,|x|^{c-d}L(1/|x|)$. Clause (2) is the boundary case. The density at distance $s$ is
asymptotic to $(2\pi)^{-d}$ times the integral of the transform cut at the frequency $1/s$,
which depends on the law and grows at least like $\log(1/s)$.

\begin{proof}[Proof of Lemma~\ref{prop:origin-rate-d}]
As $c = \sup_{r>0}k_d(r) > 0$ and $k_d$ is nonincreasing, $k_d$ is positive on an interval
$(0,r_0]$, so the datum is not trivial and Proposition~\ref{prop:isotropic-radial-monotone}(2) gives a density
$p(x) = q(|x|)$ of $\mu$ with $q$ nonincreasing. From here on $q$ is any nonincreasing function
on $(0,\infty)$ such that $x \mapsto q(|x|)$ is a density of $\mu$.

Write $\chi(\rho) := \widehat\mu(\rho e) = e^{-F_d(\rho e)}$ for $\rho > 0$. As $F_d$ is rotation
invariant, $\widehat\mu(\xi) = \chi(|\xi|)$ for $\xi \ne 0$. Three properties of $\chi$ are used.

(a) $0 < \chi \le 1$, and $\chi$ is continuous, because $F_d \ge 0$ and $\widehat\mu$ is a
characteristic function.

(b) $\chi$ is nonincreasing. In the notation of the proof of
Lemma~\ref{lem:isotropic-exponent-log}, $F_d(\xi)$ is the average over the sphere of
$\varphi(\langle\xi,\eta\rangle)$, and the substitution $u = ar$ made there gives
$\varphi(as) = \int_0^\infty(1 - \cos(us))\,k_d(u/a)\,du/u \ge \varphi(s)$ for $a \ge 1$ and every
real $s$, because $k_d(u/a) \ge k_d(u)$. So $F_d(a\xi) \ge F_d(\xi)$ for $a \ge 1$.

(c) $\chi(a\rho)/\chi(\rho) \to a^{-c}$ as $\rho \to \infty$, for every $a > 0$. This is
Lemma~\ref{lem:isotropic-exponent-log}(2), which applies because $c \le d$ is finite.

\emph{Step 1: a bound on the ratios of $\chi$.} Let $\gamma > c$. By (c) at $a = 2$,
$\chi(2y)/\chi(y) \to 2^{-c} > 2^{-\gamma}$, so there is $X \ge 1$ with
$\chi(y) \le 2^{\gamma}\chi(2y)$ for every $y \ge X$. Let $X \le y \le x$, and let $n \ge 0$ be the
least integer with $2^n y \ge x$, so that $n < \log_2(x/y) + 1$. Iterating the inequality and
using (b),
\[
  \chi(y) \le 2^{\gamma n}\,\chi(2^n y) \le 2^{\gamma n}\,\chi(x)
   \le 2^{\gamma}\,(x/y)^{\gamma}\,\chi(x) .
\]
It follows that, with $C := 2^{\gamma}/\chi(X)$,
\begin{equation}\label{eq:rate-ratio-bound}
  \frac{\chi(Ru)}{\chi(R)} \le 1 + C\,u^{-\gamma} \qquad (R \ge X,\ u > 0) .
\end{equation}
Indeed, for $u \ge 1$ the quotient is at most $1$ by (b). For $u < 1$ and $Ru \ge X$ the display
with $y = Ru$ and $x = R$ gives the bound $2^{\gamma}u^{-\gamma}$, and $2^{\gamma} \le C$ because
$\chi(X) \le 1$. For $u < 1$ and $Ru < X$ the numerator is at most $1$, and the display with
$y = X$ and $x = R$ gives $1/\chi(R) \le 2^{\gamma}(R/X)^{\gamma}/\chi(X) \le C\,u^{-\gamma}$.

\emph{Step 2: the Gaussian average of the law.} For $\lambda > 0$ put
$M(\lambda) := \int_{\RR^d}e^{-\lambda|x|^2}\mu(dx)$. At $t = 1/(2\lambda)$ the heat kernel is
$g_t(x) = (\lambda/\pi)^{d/2}e^{-\lambda|x|^2}$, and Lemma~\ref{lem:positive-transform-origin}(1) gives
$M(\lambda) = (\pi/\lambda)^{d/2}(2\pi)^{-d}\int_{\RR^d}\widehat\mu(\xi)\,e^{-|\xi|^2/4\lambda}\,d\xi$.
In polar coordinates, at $\lambda = R^2$,
\begin{equation}\label{eq:rate-gaussian-average}
  M(R^2) = (4\pi)^{-d/2}\,|S^{d-1}|\,R^{-d}\,T(R),
  \qquad
  T(R) := \int_0^\infty \chi(\rho)\,\rho^{d-1}\,e^{-\rho^2/4R^2}\,d\rho ,
\end{equation}
and the substitution $\rho = Ru$ gives
$T(R) = R^{d}\int_0^\infty \chi(Ru)\,u^{d-1}\,e^{-u^2/4}\,du$.

\emph{Step 3: the Gaussian average for $c < d$.} Let $c < d$ and take $\gamma := (c+d)/2$ in
Step 1, so that $c < \gamma < d$. By the second form of $T(R)$ in Step 2,
\[
  \frac{T(R)}{R^{d}\chi(R)} = \int_0^\infty \frac{\chi(Ru)}{\chi(R)}\,u^{d-1}\,e^{-u^2/4}\,du .
\]
For each $u > 0$ the quotient under the integral tends to $u^{-c}$ as $R \to \infty$, by (c). For
$R \ge X$ the integrand is at most $(1 + C\,u^{-\gamma})\,u^{d-1}e^{-u^2/4}$ by
\ref{eq:rate-ratio-bound}, which is integrable on $(0,\infty)$ because $\gamma < d$. By dominated
convergence
\[
  \frac{T(R)}{R^{d}\chi(R)} \to \int_0^\infty u^{d-1-c}\,e^{-u^2/4}\,du
   = 2^{d-c-1}\,\Gamma\bigl(\tfrac{d-c}{2}\bigr) \qquad (R \to \infty),
\]
and with $|S^{d-1}| = 2\pi^{d/2}/\Gamma(d/2)$,
\begin{equation}\label{eq:rate-abelian}
  \frac{M(R^2)}{\chi(R)} \to I_c := (4\pi)^{-d/2}\,|S^{d-1}|\,2^{d-c-1}\,\Gamma\bigl(\tfrac{d-c}{2}\bigr)
   = \frac{2^{-c}\,\Gamma\bigl(\tfrac{d-c}{2}\bigr)}{\Gamma\bigl(\tfrac d2\bigr)} .
\end{equation}

\emph{Step 3$'$: the Gaussian average for $c = d$.} Let $c = d$, and put
$P(y) := \int_0^y\chi(\rho)\,\rho^{d-1}\,d\rho$ for $y > 0$ and $L(\rho) := \rho^{d}\chi(\rho)$.
By (b), for $0 < x \le y$,
\begin{equation}\label{eq:rate-P-increments}
  \chi(y)\,\frac{y^{d} - x^{d}}{d} \le P(y) - P(x) \le \chi(x)\,\frac{y^{d} - x^{d}}{d} .
\end{equation}
\emph{(i) $L(R)/P(R) \to 0$ as $R \to \infty$.} The left inequality of
\ref{eq:rate-P-increments} at $x = y/2$ reads $P(y) - P(y/2) \ge \kappa\,L(y)$ with
$\kappa := (1 - 2^{-d})/d$. For each integer $j \ge 0$,
$L(R/2^j)/L(R) = 2^{-jd}\chi(2^{-j}R)/\chi(R) \to 1$ as $R \to \infty$, by (c) with $c = d$. Let
$n \ge 1$. For all large $R$, $L(R/2^j) \ge L(R)/2$ for every $j < n$, and then
\[
  P(R) \ge \sum_{j=0}^{n-1}\bigl(P(R/2^j) - P(R/2^{j+1})\bigr) \ge \kappa\sum_{j=0}^{n-1}L(R/2^j)
   \ge \tfrac{n\kappa}{2}\,L(R) .
\]
As $n$ is arbitrary, $L(R)/P(R) \to 0$.
\emph{(ii) $P(aR)/P(R) \to 1$ as $R \to \infty$, for every $a > 0$.} For $a \ge 1$ the right
inequality of \ref{eq:rate-P-increments} gives
$0 \le P(aR) - P(R) \le (a^{d} - 1)\,L(R)/d$, and (i) gives the limit. For $a < 1$ the quotient is
the reciprocal of $P(a^{-1}R')/P(R')$ at $R' = aR$, which tends to infinity with $R$.
\emph{(iii) $T(R)/P(R) \to 1$ as $R \to \infty$.} Let $\varepsilon > 0$. On $(0, \varepsilon R]$
the Gaussian factor of $T(R)$ is at least $e^{-\varepsilon^2/4}$, so
$T(R) \ge e^{-\varepsilon^2/4}P(\varepsilon R)$. On $(0, R]$ it is at most $1$, and on $(R,\infty)$,
$\chi(\rho) \le \chi(R)$ by (b), so
\[
  T(R) \le P(R) + \chi(R)\int_0^\infty\rho^{d-1}e^{-\rho^2/4R^2}\,d\rho
   = P(R) + 2^{d-1}\Gamma\bigl(\tfrac d2\bigr)\,L(R) .
\]
Hence
$e^{-\varepsilon^2/4}\,P(\varepsilon R)/P(R) \le T(R)/P(R) \le 1 + 2^{d-1}\Gamma(d/2)\,L(R)/P(R)$.
By (ii) the left side tends to $e^{-\varepsilon^2/4}$ and by (i) the right side tends to $1$. As
$\varepsilon$ is arbitrary, $T(R)/P(R) \to 1$. With \ref{eq:rate-gaussian-average},
\begin{equation}\label{eq:rate-abelian-critical}
  \frac{M(R^2)}{(4\pi)^{-d/2}\,|S^{d-1}|\,R^{-d}\,P(R)} \to 1 \qquad (R \to \infty) .
\end{equation}

\emph{Step 4: the mass of a ball, by Karamata's theorem.} Let $\nu$ be the image of $\mu$ under
$x \mapsto |x|^2$. It is a probability measure on $[0,\infty)$ with
$\int_{[0,\infty)}e^{-\lambda r}\,\nu(dr) = M(\lambda)$ and $\nu([0,r]) = \mu(\{|x| \le \sqrt r\})$.

Let $c < d$ and put $\ell(\lambda) := \lambda^{c/2}\chi(\sqrt\lambda)$ for $\lambda > 0$. It is
positive and continuous, and
$\ell(a\lambda)/\ell(\lambda) = a^{c/2}\chi(\sqrt a\sqrt\lambda)/\chi(\sqrt\lambda) \to 1$ as
$\lambda \to \infty$ for every $a > 0$, by (c). By \ref{eq:rate-abelian} at $R = \sqrt\lambda$,
$M(\lambda)/(I_c\,\lambda^{-c/2}\,\ell(\lambda)) \to 1$. Karamata's Tauberian theorem
(Theorem~\ref{thm:cited-A11}, with the index $c/2 > 0$ and the constant $I_c > 0$) gives
$\nu([0,r]) \sim I_c\,r^{c/2}\,\ell(1/r)/\Gamma(1 + c/2)$ as $r \downarrow 0$, and
$r^{c/2}\ell(1/r) = \chi(r^{-1/2})$. At $r = s^2$,
\begin{equation}\label{eq:rate-ball}
  \frac{\mu(\{|x| \le s\})}{B_c\,\chi(1/s)} \to 1 \quad (s \downarrow 0),
  \qquad B_c := \frac{I_c}{\Gamma(1 + c/2)} = \frac{R(d,c)\,|S^{d-1}|}{c} ,
\end{equation}
the last equality by $\Gamma(1 + c/2) = (c/2)\Gamma(c/2)$ and $|S^{d-1}| = 2\pi^{d/2}/\Gamma(d/2)$.
This is the first limit of clause (1).

Let $c = d$ and put $\ell(\lambda) := P(\sqrt\lambda)$. It is positive and continuous, and
$\ell(a\lambda)/\ell(\lambda) \to 1$ by Step 3$'$(ii). By \ref{eq:rate-abelian-critical},
$M(\lambda)/\bigl((4\pi)^{-d/2}|S^{d-1}|\,\lambda^{-d/2}\,\ell(\lambda)\bigr) \to 1$. Karamata's
theorem with the index $d/2$ gives, at $r = s^2$,
$\mu(\{|x| \le s\}) \sim (4\pi)^{-d/2}\,|S^{d-1}|\,s^{d}\,P(1/s)/\Gamma(1 + d/2)$. In polar
coordinates $|S^{d-1}|\,P(1/s) = I(s)$, and
$(4\pi)^{-d/2}/\Gamma(1 + d/2) = (2\pi)^{-d}\,\pi^{d/2}/\Gamma(1 + d/2) = (2\pi)^{-d}\,V_d$. This is
the first limit of clause (2).

\emph{Step 5: from the mass of balls to the density.} Write $m(s) := \mu(\{|x| \le s\})$. Let
$W$ be a positive function on $(0,\infty)$, $\varkappa > 0$ and $B > 0$, and suppose that
$W(as)/W(s) \to a^{\varkappa}$ for every $a > 0$ and $m(s)/W(s) \to B$, as $s \downarrow 0$. Then
\begin{equation}\label{eq:rate-monotone-density}
  \frac{q(s)\,s^{d}}{W(s)} \to \frac{B\,\varkappa}{|S^{d-1}|} \qquad (s \downarrow 0) .
\end{equation}
Indeed, let $0 < a < b$. On the shell $\{as < |x| \le bs\}$, whose volume is
$V_d\,s^{d}(b^{d} - a^{d})$, the density $q(|x|)$ lies between $q(bs)$ and $q(as)$, because $q$ is
nonincreasing. So
\[
  q(bs)\,V_d\,s^{d}\,(b^{d} - a^{d}) \le m(bs) - m(as) \le q(as)\,V_d\,s^{d}\,(b^{d} - a^{d}) ,
\]
and $(m(bs) - m(as))/W(s) \to B\,(b^{\varkappa} - a^{\varkappa})$. With $b = 1$ and $a < 1$ the left
inequality shows that, for every $\varepsilon > 0$, $q(s)s^{d}/W(s) \le
B(1 - a^{\varkappa})/(V_d(1 - a^{d})) + \varepsilon$ for all small $s$. With $a = 1$ and $b > 1$ the
right one shows $q(s)s^{d}/W(s) \ge B(b^{\varkappa} - 1)/(V_d(b^{d} - 1)) - \varepsilon$ for all small
$s$. Now $(t^{\varkappa} - 1)/(t^{d} - 1) \to \varkappa/d$ as $t \to 1$, the quotient of the
derivatives of $t^{\varkappa}$ and $t^{d}$ at $1$. So the first bound tends to
$B\varkappa/(V_d\,d) + \varepsilon$ as $a \uparrow 1$ and the second to
$B\varkappa/(V_d\,d) - \varepsilon$ as $b \downarrow 1$. As $\varepsilon$ is arbitrary and
$V_d\,d = |S^{d-1}|$, this proves \ref{eq:rate-monotone-density}. It is the argument of the
monotone density theorem (\sscite{bingham1987regular}, Thm.~1.7.2b, pp.~39--40), on spherical
shells.

For $c < d$ take $W(s) := \chi(1/s)$, $\varkappa := c$ and $B := B_c$: the first hypothesis is (c)
at the ratio $1/a$, the second is \ref{eq:rate-ball}. So
$q(s)\,s^{d}/\chi(1/s) \to B_c\,c/|S^{d-1}| = R(d,c)$, the second limit of clause (1). For $c = d$
take $W(s) := s^{d}I(s)$, $\varkappa := d$ and $B := (2\pi)^{-d}V_d$: the first hypothesis is
Step 3$'$(ii), since $I(s) = |S^{d-1}|P(1/s)$, and the second is the first limit of clause (2).
So $q(s)/I(s) \to (2\pi)^{-d}\,V_d\,d/|S^{d-1}| = (2\pi)^{-d}$, the second limit of clause (2).

\emph{The logarithmic bound of clause (2).} Let $c = d$ and $0 < s \le 1$. By
Lemma~\ref{lem:isotropic-exponent-log}(1), $\chi(\rho) \ge e^{-13d/4 - 2C_k}\rho^{-d}$ for $\rho \ge 1$.
In polar coordinates
\[
  I(s) \ge |S^{d-1}|\int_1^{1/s}\chi(\rho)\,\rho^{d-1}\,d\rho
   \ge |S^{d-1}|\,e^{-13d/4 - 2C_k}\int_1^{1/s}\frac{d\rho}{\rho}
   = |S^{d-1}|\,e^{-13d/4 - 2C_k}\log(1/s) .
\]
Let $C < C_* := (2\pi)^{-d}|S^{d-1}|e^{-13d/4-2C_k}$. If $C \le 0$ there is nothing to prove, as
$q(s) > 0$ for small $s$ by the second limit. Otherwise $q(s) \ge (C/C_*)(2\pi)^{-d}I(s)$ for all
small $s$, by the same limit, and the bound on $I(s)$ gives $q(s) \ge C\log(1/s)$.
\end{proof}

% Q-0428: the proofs of the two machine-checked nodes of Q-0425 on the rate through the polar
% profile, written by records/affine/transcribe-q0428-nodes.py from the blueprint and pinned
% (its header says what the transcription changes). The proof of the third node,
% lem:exponent-profile-constant, which is proved in print only, is in Appendix A of the paper.
The proofs of Propositions~\ref{prop:isotropic-exponent-profile}
and~\ref{prop:origin-rate-profile-d} follow.

\begin{proof}[Proof of Proposition~\ref{prop:isotropic-exponent-profile}]
Write $k := k_d$ and $\bar\sigma := \sigma_{d-1}/|S^{d-1}|$ for the normalized surface measure.
Three things are taken from the proof of Lemma~\ref{lem:isotropic-exponent-log}. First,
\[
  F_d(\xi) = \int_{S^{d-1}}\varphi\bigl(\langle\xi,\eta\rangle\bigr)\,\bar\sigma(d\eta),
  \qquad
  \varphi(s) := \int_0^\infty\bigl(1 - \cos(rs)\bigr)\,k(r)\,\frac{dr}{r} \quad (s \in \RR),
\]
with $\varphi$ even and finite. Second, $k \le c$ on $(0,\infty)$ and $k(r) \to 0$ as
$r \to \infty$. Third, the bounds \ref{eq:cin-log}, and $0 \le \Cin(X) \le \tfrac14$ for
$0 \le X \le 1$, because $\Cin$ is nondecreasing and $\Cin(1) \le \tfrac14$.

\emph{Clause (1).} For $0 < u \le 1$ the integrand $(c - k(u))/u$ lies in $[0, c/u]$. So
$\int_x^1 (c - k(u))\,du/u$ lies in $[0, c\log(1/x)]$, it is continuous and nonincreasing in $x$,
and it vanishes at $x = 1$. Hence $K$ is continuous and nonincreasing with
$1 = K(1) \le K(x) \le x^{-c}$. Let $a > 0$, and let $0 < x \le \min\{1, 1/a\}$, so that $x$ and
$ax$ lie in $(0,1]$. Between $x$ and $ax$ the integrand is at most $(c - k(\max\{a,1\}x))/u$,
because $k$ is nonincreasing. So
\[
  \bigl|\log K(ax) - \log K(x)\bigr| = \Bigl|\int_{ax}^{x}\bigl(c - k(u)\bigr)\frac{du}{u}\Bigr|
   \le \bigl(c - k(\max\{a,1\}\,x)\bigr)\,|\log a| ,
\]
which tends to $0$ as $x \downarrow 0$ because $k(0{+}) = c$. By monotone convergence
$\int_x^1 (c - k(u))\,du/u$ increases to $\int_0^1 (c - k(u))\,du/u$ as $x \downarrow 0$. If this
integral is finite, $K(x)$ increases to its exponential, which is $K(0{+})$. If it is infinite,
$K(x) \to \infty$.

\emph{Clause (2), the limit $g_\infty$.} Put $\log^+ X := \max\{\log X, 0\}$ for $X > 0$,
$\log^+ 0 := 0$, and
\[
  G(X) := \Cin(X) - \log^+ X \qquad (X \ge 0) .
\]
Then $-3 \le G \le \tfrac{13}{4}$ on $[0,\infty)$: for $X \ge 1$ this is \ref{eq:cin-log}, and for
$0 \le X < 1$, $G(X) = \Cin(X) \in [0,\tfrac14]$. For $X \ge 1$,
$\Cin(X) = \Cin(1) + \log X - \int_1^X u^{-1}\cos u\,du$, so
\[
  G(X) = \Cin(1) - \int_1^X \frac{\cos u}{u}\,du
       = \Cin(1) - \Bigl[\frac{\sin u}{u}\Bigr]_1^X - \int_1^X \frac{\sin u}{u^2}\,du .
\]
The bracket tends to $-\sin 1$ as $X \to \infty$, and $u^{-2}\sin u$ is absolutely integrable on
$[1,\infty)$. So $G(X)$ has a limit as $X \to \infty$. As $G(X) = \Cin(X) - \log X$ for $X \ge 1$,
this limit is $g_\infty$, and $-3 \le g_\infty \le \tfrac{13}{4}$.

\emph{Clause (2), the mean of the logarithm over the sphere.} Let $B$ be the closed unit ball of
$\RR^d$ and $e_1$ the first basis vector. Lebesgue measure is the image of
$\sigma_{d-1} \otimes r^{d-1}dr$ under $(\eta, r) \mapsto r\eta$, so for every Borel
$h : S^{d-1} \to [0,\infty]$,
\[
  \int_B h\bigl(x/|x|\bigr)\,dx = \int_0^1 r^{d-1}\,dr\int_{S^{d-1}}h\,d\sigma_{d-1}
   = \frac1d\int_{S^{d-1}}h\,d\sigma_{d-1} .
\]
Take $h(\eta) := -\log|\langle e_1,\eta\rangle| \in [0,\infty]$. For $x \in B$ with $x_1 \ne 0$,
$h(x/|x|) = -\log|x_1| + \log|x| \le -\log|x_1|$, because $|x| \le 1$. As $B$ lies in the cube
$[-1,1]^d$, on which $-\log|x_1| \ge 0$,
\[
  \frac1d\int_{S^{d-1}}h\,d\sigma_{d-1} \le \int_{[-1,1]^d}\bigl(-\log|x_1|\bigr)\,dx
   = 2^{d-1}\int_{-1}^{1}\bigl(-\log|t|\bigr)\,dt = 2^{d} .
\]
So $\log|\langle e_1,\eta\rangle|$ is $\sigma_{d-1}$-integrable, and its mean is at most $0$
because $|\langle e_1,\eta\rangle| \le 1$. For a unit vector $e$ take $Q \in \Orth(d)$ with
$Qe_1 = e$. Then $\langle e,\eta\rangle = \langle e_1, Q^{\tp}\eta\rangle$, and $\sigma_{d-1}$ is
invariant under $Q^{\tp}$, so the integral is the same for $e$ as for $e_1$. At $d = 1$ the
sphere is $\{\pm1\}$, $|\langle e,\eta\rangle| = 1$, and $\lambda_1 = 0$.

\emph{Clause (3), the tail integral.} For $s > 0$ put
\[
  T(s) := \int_{1/s}^{\infty}k(r)\,\frac{dr}{r} .
\]
It is finite: it is at most $C_k$ for $s \le 1$ and at most $c\log s + C_k$ for $s \ge 1$. For
$s \ge 1$,
\begin{equation}\label{eq:profile-tail}
  T(s) = \int_{1/s}^{1}k(r)\,\frac{dr}{r} + C_k
       = c\log s - \int_{1/s}^{1}\bigl(c - k(r)\bigr)\frac{dr}{r} + C_k
       = c\log s - \log K(1/s) + C_k .
\end{equation}

\emph{Clause (3), the layer-cake identity.} For $0 < y < c$ put
$\tau(y) := \sup\{r > 0 : k(r) > y\}$. The set is not empty, because $\sup k = c > y$. It is
bounded, because $k(r) \to 0$ as $r \to \infty$ and $y > 0$. So $0 < \tau(y) < \infty$. As $k$ is
nonincreasing, the set contains every $r' < r$ with each of its points $r$, so it is
$(0,\tau(y))$ or $(0,\tau(y)]$; and $\tau$ is nonincreasing in $y$, hence Borel. Since
$0 \le k \le c$,
\[
  k(r) = \int_0^c \mathbf 1\{y < k(r)\}\,dy \qquad (r > 0),
\]
and the indicator is a Borel function of $(r,y)$, $k$ being monotone. For a Borel
$w : (0,\infty) \to [0,\infty)$, Tonelli's theorem gives
\[
  \int_0^\infty w(r)\,k(r)\,dr = \int_0^c\Bigl(\int_0^{\tau(y)}w(r)\,dr\Bigr)dy .
\]
Let $s > 0$. With $w(r) = (1 - \cos rs)/r$ the inner integral is $\Cin(s\,\tau(y))$, by the
substitution $u = rs$. With $w(r) = \mathbf 1\{r > 1/s\}/r$ it is $\log^+(s\,\tau(y))$. So
\[
  \varphi(s) = \int_0^c \Cin\bigl(s\,\tau(y)\bigr)\,dy, \qquad
  T(s) = \int_0^c \log^+\bigl(s\,\tau(y)\bigr)\,dy .
\]
Both are finite, and their integrands differ by $G(s\,\tau(y))$, which is bounded. Hence
\begin{equation}\label{eq:profile-layer-cake}
  \varphi(s) = T(s) + \int_0^c G\bigl(s\,\tau(y)\bigr)\,dy \qquad (s > 0) .
\end{equation}
This is an identity, with no remainder. It is where the monotonicity of $k$ enters the exponent.
Monotonicity is also used for $k \le c$ and $k(r) \to 0$, in clause (1), and in the first term
below.

\emph{Clause (3), the sphere average.} Let $\rho > 0$ and write
$\eta_1 := \langle e,\eta\rangle$ for $\eta \in S^{d-1}$. As $\varphi$ is even,
$F_d(\rho e) = \int\varphi(\rho|\eta_1|)\,\bar\sigma(d\eta)$. The set $\{\eta_1 = 0\}$ is
$\bar\sigma$-null: it is a great subsphere when $d \ge 2$ and empty when $d = 1$. By
\ref{eq:profile-layer-cake} at $s = \rho|\eta_1|$,
\begin{equation}\label{eq:profile-sphere}
  F_d(\rho e) - T(\rho) = \int_{S^{d-1}}\bigl[T(\rho|\eta_1|) - T(\rho)\bigr]\,\bar\sigma(d\eta)
   + \int_{S^{d-1}}\int_0^c G\bigl(\rho|\eta_1|\,\tau(y)\bigr)\,dy\,\bar\sigma(d\eta) ,
\end{equation}
where both integrals converge absolutely by the bounds that follow.

\emph{The first term.} Let $0 < t \le 1$. Then
$T(\rho) - T(\rho t) = \int_{1/\rho}^{1/(\rho t)}k(r)\,dr/r$ lies between
$k(1/(\rho t))\log(1/t)$ and $c\log(1/t)$. So $-c\log(1/t) \le T(\rho t) - T(\rho) \le 0$ for
every $\rho > 0$, and $T(\rho t) - T(\rho) \to c\log t$ as $\rho \to \infty$, because
$k(0{+}) = c$. At $t = |\eta_1|$ the bound $c\,|\log|\eta_1||$ is $\bar\sigma$-integrable by
clause (2). By dominated convergence the first term of \ref{eq:profile-sphere} tends to
$c\lambda_d$ as $\rho \to \infty$, and it lies in $[c\lambda_d, 0]$ for every $\rho > 0$.

\emph{The second term.} The integrand lies in $[-3, \tfrac{13}{4}]$. For every $y \in (0,c)$ and
every $\eta$ with $\eta_1 \ne 0$, $\rho|\eta_1|\tau(y) \to \infty$ as $\rho \to \infty$, because
$\tau(y) > 0$, so the integrand tends to $g_\infty$. The function
$(\eta,y) \mapsto G(\rho|\eta_1|\tau(y))$ is Borel on $S^{d-1} \times (0,c)$ and bounded: $\tau$
is Borel, being monotone, and $G$ is continuous on $[0,\infty)$. So the inner integral
$\int_0^c G(\rho|\eta_1|\tau(y))\,dy$ is a Borel function of $\eta$, and it lies in
$[-3c, \tfrac{13}{4}c]$ for every $\rho > 0$ and every $\eta$. For $\eta$ with $\eta_1 \ne 0$ it
tends to $c\,g_\infty$ as $\rho \to \infty$, by bounded convergence on $(0,c)$. By bounded
convergence for $\bar\sigma$, the second term tends to $c\,g_\infty$, and it lies in
$[-3c, \tfrac{13}{4}c]$ for every $\rho > 0$.

Adding the two, $F_d(\rho e) - T(\rho) \to c(\lambda_d + g_\infty) = c\varkappa_d$ as
$\rho \to \infty$, and $c(\lambda_d - 3) \le F_d(\rho e) - T(\rho) \le \tfrac{13}{4}c$ for every
$\rho > 0$. As $T(\rho) = \int_{1/\rho}^\infty k(r)\,dr/r$, this is the bound of clause (3) for
every $\rho > 0$. For $\rho \ge 1$, \ref{eq:profile-tail} replaces $T(\rho)$ by
$c\log\rho - \log K(1/\rho) + C_k$. This is the limit and the two-sided bound of clause (3) for
$\rho \ge 1$. The
quotient form follows by taking exponentials: the quotient is $\exp$ of minus the difference
between $F_d(\rho e) - (c\log\rho - \log K(1/\rho) + C_k)$ and its limit.

\emph{Clause (4).} $F_d$ is rotation invariant, since $\sigma_{d-1}$ is, so
$F_d(\xi) = F_d(|\xi|e)$. Let $\rho = |\xi| \ge 1$ and write
$F_d(\rho e) = c\log\rho - \log K(1/\rho) + C_k + \Delta(\rho)$. Then
$\rho^{c}e^{-F_d(\rho e)} = e^{-C_k - \Delta(\rho)}K(1/\rho)$, and the two-sided bound of clause
(3) on $\Delta(\rho)$ gives the display. By the display, $\rho^{c}e^{-F_d(\rho e)}$ is bounded on
$[1,\infty)$ exactly when $K$ is bounded on $(0,1]$, which by clause (1) is the condition on the
integral. In that case $K(1/\rho) \to K(0{+}) < \infty$ and $\Delta(\rho) \to c\varkappa_d$, which
gives the limit.
\end{proof}

\begin{proof}[Proof of Proposition~\ref{prop:origin-rate-profile-d}]
A density $p(x) = q(|x|)$ with $q$ nonincreasing exists by
Proposition~\ref{prop:isotropic-radial-monotone}(2), as in the proof of Lemma~\ref{prop:origin-rate-d}. From here on
$q$ is any nonincreasing function on $(0,\infty)$ such that $x \mapsto q(|x|)$ is a density of
$\mu$. Let $e$ be a unit vector and write $\chi(\rho) := \widehat\mu(\rho e) = e^{-F_d(\rho e)}$
for $\rho > 0$. Since $0 < c \le d < \infty$, Proposition~\ref{prop:isotropic-exponent-profile}(3) applies and
gives $\chi(\rho)/(A\,\rho^{-c}K(1/\rho)) \to 1$ as $\rho \to \infty$. At $\rho = 1/s$,
\begin{equation}\label{eq:profile-transform-at-s}
  \frac{\widehat\mu(e/s)}{A\,s^{c}\,K(s)} \to 1 \qquad (s \downarrow 0) .
\end{equation}

\emph{Clause (1).} Let $c < d$. By Lemma~\ref{prop:origin-rate-d}(1),
$q(s)/(R(d,c)\,s^{-d}\,\widehat\mu(e/s)) \to 1$ and
$\mu(\{|x| \le s\})/(R(d,c)\,|S^{d-1}|\,c^{-1}\,\widehat\mu(e/s)) \to 1$. Multiplying each by
\ref{eq:profile-transform-at-s} gives the two limits, with $Q = R(d,c)\,A$.

\emph{Clause (2).} By clause (1), $q(s)\,s^{d-c}/(Q\,K(s)) \to 1$, and $Q > 0$. By
Proposition~\ref{prop:isotropic-exponent-profile}(1), as $s \downarrow 0$ the factor $K(s)$ increases to the
finite number $K(0{+}) \ge 1$ if $\int_0^1(c - k_d(u))\,du/u < \infty$, and $K(s) \to \infty$ if
the integral diverges. In the second case $q(s)\,s^{d-c} \to \infty$, and no bound
$q(s) \le C s^{c-d}$ holds near $0$. In the first case $q(s)\,s^{d-c} \to Q\,K(0{+})$, which is
finite and positive, so there is $s_0 > 0$ with
$q(s)\,s^{d-c} \le 2Q\,K(0{+})$ for $0 < s \le s_0$. For $s \ge s_0$, as $q$ is nonincreasing,
$1 \ge \mu(\{|x| < s\}) \ge q(s)\,V_d\,s^{d}$, so
$q(s)\,s^{d-c} \le s^{-c}/V_d \le s_0^{-c}/V_d$. So
$C := \max\{2Q\,K(0{+}),\,s_0^{-c}/V_d\}$ serves for every $s > 0$.

\emph{Clause (3), the truncated integral of the transform.} Let $c = d$. In polar coordinates the
$I(s)$ of Lemma~\ref{prop:origin-rate-d}(2) is
\[
  I(s) = \int_{\{|\xi| \le 1/s\}}\widehat\mu(\xi)\,d\xi = |S^{d-1}|\,P(1/s), \qquad
  P(R) := \int_0^R \chi(\rho)\,\rho^{d-1}\,d\rho .
\]
Put $f(\rho) := \chi(\rho)\,\rho^{d-1}$ and $h(\rho) := A\,K(1/\rho)/\rho$ for $\rho \ge 1$. Both
are positive and continuous, and $f(\rho)/h(\rho) = \chi(\rho)\,\rho^{d}/(A\,K(1/\rho)) \to 1$
as $\rho \to \infty$, by Proposition~\ref{prop:isotropic-exponent-profile}(3) with $c = d$. The substitution
$y = 1/\rho$ gives, for $R \ge 1$,
\[
  H(R) := \int_1^R h(\rho)\,d\rho = A\int_{1/R}^{1}K(y)\,\frac{dy}{y} = A\,\mathcal{L}(1/R) ,
\]
and $H(R) \ge A\log R \to \infty$, because $K \ge 1$. Let $0 < \varepsilon < 1$ and choose
$R_0 \ge 1$ with $|f - h| \le \varepsilon h$ on $[R_0,\infty)$. For $R \ge R_0$,
\[
  \bigl|P(R) - H(R)\bigr| \le P(1) + \int_1^{R_0}|f - h|\,d\rho + \varepsilon\,H(R) ,
\]
and $P(1) \le 1/d$ because $\chi \le 1$. Dividing by $H(R)$, which tends to infinity,
$|P(R)/H(R) - 1| \le 2\varepsilon$ for all large $R$. So $P(R)/(A\,\mathcal{L}(1/R)) \to 1$ as
$R \to \infty$, that is
\[
  \frac{I(s)}{|S^{d-1}|\,A\,\mathcal{L}(s)} \to 1 \qquad (s \downarrow 0) .
\]

\emph{Clause (3), the two limits.} By Lemma~\ref{prop:origin-rate-d}(2),
$q(s)/((2\pi)^{-d}I(s)) \to 1$ and $\mu(\{|x| \le s\})/((2\pi)^{-d}V_d\,s^{d}\,I(s)) \to 1$.
Multiplying each by the last display gives the two limits, with
$B = (2\pi)^{-d}|S^{d-1}|A$.

\emph{Clause (3), the logarithm.} Let $0 < s \le 1$. On $(s,1)$, $1 \le K(y) \le K(s)$, so
$\log(1/s) \le \mathcal{L}(s) \le K(s)\log(1/s)$. Let $0 < y_0 < 1$. For $0 < s < y_0$,
$\mathcal{L}(s) \ge \int_s^{y_0}K(y)\,dy/y \ge K(y_0)\log(y_0/s)$, and
$\log(y_0/s)/\log(1/s) \to 1$ as $s \downarrow 0$. So for every number $\kappa < K(y_0)$,
$\mathcal{L}(s)/\log(1/s) \ge \kappa$ for all small $s$. Now use
Proposition~\ref{prop:isotropic-exponent-profile}(1) at $c = d$. If $\int_0^1(d - k_d(u))\,du/u < \infty$,
then $K(y_0)$ increases to the finite number $K(0{+})$ as $y_0 \downarrow 0$, so the lower bound
holds for every $\kappa < K(0{+})$; and $\mathcal{L}(s) \le K(s)\log(1/s) \le K(0{+})\log(1/s)$.
Hence $\mathcal{L}(s)/\log(1/s) \to K(0{+})$, and with the first limit of the clause
$q(s)/\log(1/s) \to B\,K(0{+})$, which is finite and positive. If the integral diverges, then
$K(y_0) \to \infty$ as $y_0 \downarrow 0$, so the lower bound holds for every real $\kappa$:
$\mathcal{L}(s)/\log(1/s) \to \infty$, and with the first limit $q(s)/\log(1/s) \to \infty$,
because $B > 0$.
\end{proof}

% Q-0448, the joining step: the proof of Proposition prop:isotropic-stages, in a file of its
% own. It stood in section A.4 of the paper (appendix-isotropic.tex) while the proposition was
% proved in print only; Lean checks it since Q-0447.
% Q-0448: the proof of Proposition prop:isotropic-stages, the three proofs of record of the
% blueprint nodes lem:isotropic-stage-levy-measure, prop:isotropic-stage-atom and
% prop:isotropic-stage-density in turn (records/affine/transcribe-q0448-nodes.py says what the
% transcription changes; pinned). It is a file of its own so that it moves in one line: while
% the proposition was proved in print only, appendix-isotropic.tex input it (section A.4); since
% the Lean of queue item Q-0447 is merged (the joining step of Q-0448), supp-proofs.tex does.
% The closing paragraph of each proof of record, which Q-0447 rewrote (what is machine-checked,
% where the Lean takes another step), is not printed, as for every proof of the paper. One
% sentence of the second node's closing paragraph is printed, after item (d) of clause (2), and
% follows the blueprint: it says "on R^d, d >= 1" since Q-0447.
% The first paragraph states the lemma as its node does, for a datum and two numbers.
% proof of lem:isotropic-stage-levy-measure
\begin{proof}[Proof of Proposition~\ref{prop:isotropic-stages}]
\emph{Clause (1).} Write $F_d(\xi) := F(|\xi|)$ for the exponent of the datum. The clause is
proved for two numbers $0 < \ell < L$ in the place of $\chi(s)$ and $\chi(t)$. Put
$\kappa(r) := k_d(r/L) - k_d(r/\ell)$ for $r > 0$, and let $\nu$ be the measure on
$\RR^d \setminus \{0\}$ with the density $\kappa(|x|)\,\bigl(|S^{d-1}|\,|x|^d\bigr)^{-1}$. Three
assertions are proved: (i) $\kappa \ge 0$, the measure $\nu$ is rotation invariant, and
$\int(|x|^2 \wedge 1)\,\nu(dx) < \infty$; (ii) for every $\xi \in \RR^d$,
\[
  F_d(L\xi) - F_d(\ell\xi) = \tfrac12 a\,(L^2 - \ell^2)\,|\xi|^2
     + \int_{\RR^d}\bigl(1 - \cos\langle\xi,x\rangle\bigr)\,\nu(dx) ;
\]
(iii) $\nu(\RR^d \setminus \{0\}) = c\,\log(L/\ell)$, in $[0,\infty]$.

\emph{The dilation of a datum.} For $\beta > 0$ put $k^{(\beta)}(r) := k_d(r/\beta)$. It is
nonnegative and nonincreasing, and
\[
  \int_0^1 k^{(\beta)}(r)\,r\,dr = \beta^2\int_0^{1/\beta}k_d(u)\,u\,du < \infty, \qquad
  \int_1^\infty k^{(\beta)}(r)\,\frac{dr}{r} = \int_{1/\beta}^\infty k_d(u)\,\frac{du}{u} < \infty ,
\]
by the two conditions on $k_d$ and because $k_d$ is bounded on the interval between $1$ and
$1/\beta$. So $(a\beta^2, k^{(\beta)})$ is an isotropic origin datum in dimension $d$. Let
$\nu^{(\beta)}$ be its \Levy{} measure, of density $k_d(|x|/\beta)\,(|S^{d-1}||x|^d)^{-1}$, and
$\nu_d$ that of $(a,k_d)$. The substitution $y = \beta x$, $dy = \beta^d dx$, gives
$\int g(\beta x)\,\nu_d(dx) = \int g\,d\nu^{(\beta)}$ for every nonnegative Borel $g$. With
$g(y) = 1 - \cos\langle\xi,y\rangle$,
\begin{equation}\label{eq:dilated-exponent}
  F_d(\beta\xi) = \tfrac12 a\beta^2|\xi|^2
     + \int_{\RR^d}\bigl(1 - \cos\langle\xi,y\rangle\bigr)\,\nu^{(\beta)}(dy) ,
\end{equation}
so the exponent of $(a\beta^2, k^{(\beta)})$ is $\xi \mapsto F_d(\beta\xi)$, and the integral is
finite, as $\int(1 - \cos\langle\xi,y\rangle)\,d\nu^{(\beta)}
\le \int(|\xi|^2|y|^2/2 \wedge 2)\,d\nu^{(\beta)} < \infty$ by \ref{eq:polar-decomposition} and
the two conditions on $k^{(\beta)}$.

\emph{Assertion (i).} For $r > 0$, $r/L < r/\ell$, and $k_d$ is nonincreasing, so
$\kappa(r) = k^{(L)}(r) - k^{(\ell)}(r) \ge 0$. The density of $\nu$ is the difference of the
densities of $\nu^{(L)}$ and $\nu^{(\ell)}$, so
\begin{equation}\label{eq:stage-measure-sum}
  \nu + \nu^{(\ell)} = \nu^{(L)}
\end{equation}
as measures on $\RR^d \setminus \{0\}$. In particular $\nu \le \nu^{(L)}$, and
$\int(|x|^2 \wedge 1)\,d\nu \le \int(|x|^2 \wedge 1)\,d\nu^{(L)}$, which is finite: by
\ref{eq:polar-decomposition} it is $\int_0^1 k^{(L)}(r)\,r\,dr + \int_1^\infty k^{(L)}(r)\,dr/r$.
The density of $\nu$ is a function of $|x|$ and Lebesgue measure is invariant under
$\Orth(d)$, so $\nu$ is.

\emph{Assertion (ii).} Let $g(x) := 1 - \cos\langle\xi,x\rangle \ge 0$. By
\ref{eq:dilated-exponent} the integrals of $g$ against $\nu^{(L)}$ and $\nu^{(\ell)}$ are finite,
so \ref{eq:stage-measure-sum} gives
$\int g\,d\nu = \int g\,d\nu^{(L)} - \int g\,d\nu^{(\ell)}$. By \ref{eq:dilated-exponent} at
$\beta = L$ and at $\beta = \ell$ this is
$F_d(L\xi) - F_d(\ell\xi) - \tfrac12 a(L^2 - \ell^2)|\xi|^2$.

\emph{Assertion (iii).} Let $0 < \varepsilon < A$ with $\varepsilon/\ell \le A/L$, and
$B := \{x : \varepsilon < |x| \le A\}$. By \ref{eq:polar-decomposition} for the datum
$(a\beta^2, k^{(\beta)})$ and $g = \mathbf 1_B$, and the substitution $u = r/\beta$,
$\nu^{(\beta)}(B) = \int_\varepsilon^A k_d(r/\beta)\,dr/r
= \int_{\varepsilon/\beta}^{A/\beta}k_d(u)\,du/u$, which is finite. By
\ref{eq:stage-measure-sum}, and since the intervals $[\varepsilon/L, A/L]$ and
$[\varepsilon/\ell, A/\ell]$ overlap in $[\varepsilon/\ell, A/L]$,
\[
  \nu(B) = \int_{\varepsilon/L}^{A/L}k_d(u)\,\frac{du}{u} - \int_{\varepsilon/\ell}^{A/\ell}k_d(u)\,\frac{du}{u}
         = \int_{\varepsilon/L}^{\varepsilon/\ell}k_d(u)\,\frac{du}{u}
           - \int_{A/L}^{A/\ell}k_d(u)\,\frac{du}{u} .
\]
Each of the last two integrals is over an interval whose endpoints have the ratio $L/\ell$, so
$\int du/u = \log(L/\ell)$ over it. As $k_d$ is nonincreasing, the first lies between
$k_d(\varepsilon/\ell)\log(L/\ell)$ and $k_d(\varepsilon/L)\log(L/\ell)$, and both bounds tend
to $c\log(L/\ell)$ as $\varepsilon \downarrow 0$, also when $c = \infty$. The second is at most
$k_d(A/L)\log(L/\ell)$, which tends to $0$ as $A \to \infty$: a nonincreasing nonnegative
function with $\int_1^\infty k_d(r)\,dr/r < \infty$ tends to $0$ at infinity. With
$\varepsilon = 1/n$ and $A = n$ for $n \ge (L/\ell)^{1/2}$, so that
$\varepsilon/\ell \le A/L$, the sets $B$ increase to $\RR^d \setminus \{0\}$, so
$\nu(\RR^d \setminus \{0\}) = \lim_n\nu(B) = c\log(L/\ell)$. Since $\log(L/\ell) > 0$, this is
finite exactly when $c$ is.

\emph{The clause.} A gauge is increasing with $\chi(0) = 0$, so
$0 < \chi(s) < \chi(t)$. By \ref{eq:isotropic-classification} and $F(|\zeta|) = F_d(\zeta)$,
$-\log\widehat\mu_{s,t}(\xi) = F_d(\chi(t)\xi) - F_d(\chi(s)\xi)$, and assertions (i), (ii) and
(iii) with $\ell = \chi(s)$, $L = \chi(t)$ are the clause.

\emph{Clause (2).} Write $\nu := \nu_{s,t}$, $\ell := \chi(s)$ and $L := \chi(t)$. By
clause (1), $\nu$ is a finite measure of mass
$m = c\log(L/\ell)$, rotation invariant, hence symmetric, with a density, and by
\ref{eq:stage-transform} with $a = 0$,
$\widehat\mu_{s,t}(\xi) = \exp[-\int(1 - \cos\langle\xi,x\rangle)\,\nu(dx)]$. If $c$ were $0$,
the nonincreasing nonnegative function $k_d$ would vanish on $(0,\infty)$, the exponent of the
datum $(0,k_d)$ would be $0$, and $F$ would vanish on $[0,\infty)$. So $c > 0$, $m > 0$ and
$0 < e^{-m} < 1$, and $e^{-m} = (\ell/L)^c = \lambda^c$.

\emph{Item (a).} Put $\pi := e^{-m}\sum_{n \ge 0}\nu^{*n}/n!$. The measure $\nu^{*n}$ has the
mass $m^n$, so $\pi$ is a probability measure. The transform of $\nu^{*n}$ is $\widehat\nu^{\,n}$
and $|\widehat\nu| \le m$, so $\widehat\pi = e^{-m}\sum_n\widehat\nu^{\,n}/n!
= \exp(\widehat\nu - m)$. As $\nu$ is symmetric, $\widehat\nu(\xi) = \int\cos\langle\xi,x\rangle
\,\nu(dx)$, and $\widehat\nu(\xi) - m = -\int(1 - \cos\langle\xi,x\rangle)\,\nu(dx)$. So
$\widehat\pi = \widehat\mu_{s,t}$, and $\mu_{s,t} = \pi$, a probability measure being determined
by its transform.

\emph{Item (b).} Let $M$ be a finite measure and $N$ a Lebesgue-null Borel set. Then
$(\nu * M)(N) = \int\nu(N - y)\,M(dy) = 0$, because $N - y$ is Lebesgue-null and $\nu$ has a
density. So $\nu^{*n} = \nu * \nu^{*(n-1)}$ is absolutely continuous for every $n \ge 1$, and so
is $\eta_{s,t}$, a countable sum of such measures. Its mass is
$e^{-m}\sum_{n \ge 1}m^n/n! = 1 - e^{-m}$. The term $n = 0$ of item (a) is $e^{-m}\delta_0$.
So $\mu_{s,t} = \lambda^c\delta_0 + \eta_{s,t}$, and
$\mu_{s,t}(\{0\}) = \lambda^c + \eta_{s,t}(\{0\}) = \lambda^c$, the origin being a Lebesgue-null
set. A measure with a density gives the origin no mass, and $\lambda^c > 0$, so $\mu_{s,t}$ has
no density.

\emph{Item (c).} The convolution $\mu * f$ is additive in the measure $\mu$, and
$\delta_0 * f = f$. With item (b),
$\mu_{s,t} * f = \lambda^c f + \eta_{s,t} * f$, and $\eta_{s,t} * f = p_{s,t} * f$.

\emph{Item (d).} Suppose that $\mu_{s,t}$ is self-decomposable, fix $q \in (0,1)$, and let
$\rho$ be a residual: a probability measure with
$\widehat\mu_{s,t}(\xi) = \widehat\mu_{s,t}(q\xi)\,\widehat\rho(\xi)$ for every $\xi$, that is
$\mu_{s,t} = (qI)_*\mu_{s,t} * \rho$. By item (b),
$(qI)_*\mu_{s,t} = \lambda^c\delta_0 + (qI)_*\eta_{s,t}$, and $(qI)_*\eta_{s,t}$ is absolutely
continuous. As in item (b), $((qI)_*\eta_{s,t} * \rho)(\{0\}) = 0$. So the mass of the origin
under the right side is $\lambda^c\rho(\{0\})$, and under the left side it is $\lambda^c$. As
$\lambda^c > 0$, $\rho(\{0\}) = 1$, that is $\rho = \delta_0$. Then
$\widehat\mu_{s,t}(\xi) = \widehat\mu_{s,t}(q\xi)$ for every $\xi$, and by iteration
$\widehat\mu_{s,t}(\xi) = \widehat\mu_{s,t}(q^n\xi) \to \widehat\mu_{s,t}(0) = 1$, the transform
being continuous. So $\mu_{s,t} = \delta_0$, which contradicts $\mu_{s,t}(\{0\}) = \lambda^c < 1$.
The argument of item (d) shows more than the item: on $\RR^d$, $d \ge 1$, a probability measure that is the
sum of a point mass of weight in $(0,1)$ at the origin and an absolutely continuous measure is
not self-decomposable.

\emph{Clause (3).} Write $\nu := \nu_{s,t}$, $\ell := \chi(s)$, $L := \chi(t)$ and $v := a(L^2 - \ell^2) \ge 0$. By
clause (1), $\nu$ gives the origin no mass, is symmetric, and has
$\int(|x|^2 \wedge 1)\,d\nu < \infty$: it is a symmetric \Levy{} measure.

\emph{A symmetric triplet.} Let $\nu_0$ be a symmetric \Levy{} measure on $\RR^d$ and
$v_0 \ge 0$. By \ledger{A3}(iii) there is an infinitely divisible law with the triplet
$(v_0I, \nu_0, 0)$, and its transform is
\begin{equation}\label{eq:symmetric-triplet}
  \exp\Bigl[-\tfrac12 v_0|\xi|^2 - \int_{\RR^d}\bigl(1 - \cos\langle\xi,x\rangle\bigr)\,\nu_0(dx)\Bigr] .
\end{equation}
Indeed the integrand of the representation, $e^{i\langle\xi,x\rangle} - 1
- i\langle\xi,x\rangle\mathbf 1_D(x)$, has the real part $\cos\langle\xi,x\rangle - 1$ and an
imaginary part that is an odd function of $x$, integrable against $\nu_0$ because it is at most
$|\xi|^2|x|^2/2$ in modulus on $D$ (as $|\sin u - u| \le u^2/2$ for every real $u$) and at most
$1$ off $D$; its integral against the symmetric
measure $\nu_0$ is $0$.

\emph{Item (a).} By \ref{eq:symmetric-triplet} with $(v_0,\nu_0) = (v,\nu)$ and
\ref{eq:stage-transform}, the infinitely divisible law with the triplet $(vI, \nu, 0)$ has the
transform of $\mu_{s,t}$, so it is $\mu_{s,t}$. The polar profile and the mass of $\nu$ are
those of clause (1).

\emph{Item (b).} Let $a = 0$ and $c = \infty$, so that $\nu$ has infinite mass
(clause (1)). For $\varepsilon > 0$ let $\nu^{>}_\varepsilon$ and
$\nu^{\le}_\varepsilon$ be the restrictions of $\nu$ to $\{|x| > \varepsilon\}$ and to
$\{0 < |x| \le \varepsilon\}$. The first is a finite measure, because
$\int(|x|^2 \wedge 1)\,d\nu < \infty$; it is symmetric and has a density; and its mass
$m_\varepsilon$ increases to $\nu(\RR^d \setminus \{0\}) = \infty$ as $\varepsilon \downarrow 0$.
The second is a symmetric \Levy{} measure. Let $\mu^\varepsilon$ be the law of
\ref{eq:symmetric-triplet} for $(0, \nu^{\le}_\varepsilon)$, and put
$\pi_\varepsilon := e^{-m_\varepsilon}\sum_{n \ge 0}(\nu^{>}_\varepsilon)^{*n}/n!$. As in the
proof of items (a) and (b) of clause (2), which uses of the measure only that it is
finite, symmetric and has a density, $\pi_\varepsilon$ is a probability measure with the
transform $\exp[-\int(1 - \cos\langle\xi,x\rangle)\,\nu^{>}_\varepsilon(dx)]$, and
$\pi_\varepsilon = e^{-m_\varepsilon}\delta_0 + \eta_\varepsilon$ with $\eta_\varepsilon$
absolutely continuous. The transforms of $\pi_\varepsilon$ and $\mu^\varepsilon$ multiply to
\ref{eq:stage-transform} with $a = 0$, so $\mu_{s,t} = \pi_\varepsilon * \mu^\varepsilon$. Let
$N$ be a Lebesgue-null Borel set. Then
$(\eta_\varepsilon * \mu^\varepsilon)(N) = \int\eta_\varepsilon(N - y)\,\mu^\varepsilon(dy) = 0$,
so
\[
  \mu_{s,t}(N) = e^{-m_\varepsilon}\mu^\varepsilon(N) \le e^{-m_\varepsilon} .
\]
Letting $\varepsilon \downarrow 0$ gives $\mu_{s,t}(N) = 0$. So $\mu_{s,t}$ is absolutely
continuous, and the origin is a Lebesgue-null set.

\emph{Item (c).} Let $a > 0$. A gauge is increasing, so $v = a(L^2 - \ell^2) > 0$. By
\ref{eq:stage-transform}, $0 < \widehat\mu_{s,t}(\xi) \le e^{-v|\xi|^2/2}$, so
$\int|\widehat\mu_{s,t}(\xi)|\,|\xi|^n\,d\xi < \infty$ for every integer $n \ge 0$. By \ledger{A1}
at $n$, $\mu_{s,t}$ has a density $g_n$ of class $C^n$ whose partial derivatives of orders
$0,\dots,n$ tend to $0$ at infinity. Two continuous densities of one measure are equal, so the
$g_n$ are one function $p$. It is of class $C^n$ for every $n$, and each of its partial
derivatives is continuous and tends to $0$ at infinity, hence is bounded.

\emph{The last sentence.} If $a = 0$ and $c < \infty$, item (b) of clause (2) gives
the point mass. If $a = 0$ and $c = \infty$, or $a > 0$, items (b) and (c) give a density,
which excludes a point mass.
\end{proof}


\subsection{The dimension filtration}\label{supp:proofs-filtration}

The subsection has the proof of the lemma on the shadow, then the separating family: three
lemmas on its critical values, the proofs of these and of the lemma on a delay datum, the
top of the filtration, and the proof of Theorem~\ref{thm:separating-family-critical-values}. It
ends with a second proof of Proposition~\ref{prop:profile-one-dimension-down}(3).

\begin{proof}[Proof of Lemma~\ref{lem:shadow-polar-profile}]
\emph{The integral over the last coordinate.} Let $h : (0,\infty) \to [0,\infty)$ be Borel and
put $\bar h(r) := C_d\int_{\RR}h(r\sqrt{1+u^2})\,(1+u^2)^{-(d+1)/2}\,du \in [0,\infty]$ for
$r > 0$. Let $g \ge 0$ be Borel on $\RR^d$ with $g(0) = 0$, and write $z = (x,y)$ with
$x \in \RR^d$ and $y \in \RR$. By Tonelli's theorem,
\[
  \int_{\RR^{d+1}} g(x)\,\frac{h(|z|)}{|S^{d}|\,|z|^{d+1}}\,dz
   = \int_{\RR^d} g(x)\Bigl[\int_{\RR}
       \frac{h\bigl(\sqrt{|x|^2+y^2}\bigr)}{|S^{d}|\,(|x|^2+y^2)^{(d+1)/2}}\,dy\Bigr]dx .
\]
For $x \ne 0$ put $r = |x|$ and substitute $y = ru$. The inner integral becomes
\[
  (|S^{d}|\,r^{d})^{-1}\int_{\RR}h(r\sqrt{1+u^2})\,(1+u^2)^{-(d+1)/2}\,du =
  \bar h(r)\,(|S^{d-1}|\,r^{d})^{-1} .
\]
As $g(0) = 0$,
\begin{equation}\label{eq:shadow-identity}
  \int_{\RR^{d+1}} g(\Pi z)\,\frac{h(|z|)}{|S^{d}|\,|z|^{d+1}}\,dz
   = \int_{\RR^d} g(x)\,\frac{\bar h(|x|)}{|S^{d-1}|\,|x|^{d}}\,dx ,
\end{equation}
both sides in $[0,\infty]$.

\emph{Clause (1).} Take $h = \mathbf 1_{(1,2)}$ and $g = 1$ off the origin in
\ref{eq:shadow-identity}. In polar coordinates of $\RR^{d+1}$ the left side is
$\int_1^2 dt/t = \log 2$; the set $\{\Pi z = 0\}$ is a line, which is null because $d + 1 \ge 2$.
In polar coordinates of $\RR^d$, and by Tonelli's theorem, the right side is
\[
\begin{aligned}
  \int_0^\infty \bar h(r)\,\frac{dr}{r}
   &= C_d\int_{\RR}(1+u^2)^{-(d+1)/2}\Bigl[\int_0^\infty
       \mathbf 1_{(1,2)}\bigl(r\sqrt{1+u^2}\bigr)\frac{dr}{r}\Bigr]du \\
   &= C_d\int_{\RR}(1+u^2)^{-(d+1)/2}\,du\cdot\log 2 ,
\end{aligned}
\]
the inner integral being over an interval whose endpoints have the ratio $2$. So
$C_d\int_{\RR}(1+u^2)^{-(d+1)/2}du = 1$. As $k'$ is nonincreasing and nonnegative,
$0 \le k'(r\sqrt{1+u^2}) \le k'(r)$, so $0 \le \bar k(r) \le k'(r) < \infty$.

\emph{Clause (2).} For each $u$, $r \mapsto k'(r\sqrt{1+u^2})$ is nonincreasing, so $\bar k$ is.
Let $r_0 > 0$. A nonincreasing function has at most countably many discontinuities, and for each
of them, $t$, at most two real $u$ have $r_0\sqrt{1+u^2} = t$. So for all $u$ outside a countable
set, $k'(r\sqrt{1+u^2}) \to k'(r_0\sqrt{1+u^2})$ as $r \to r_0$. For $r \ge r_0/2$ the integrand
of $\bar k(r)$ is at most $k'(r_0/2)\,(1+u^2)^{-(d+1)/2}$, which is integrable. By dominated
convergence $\bar k(r) \to \bar k(r_0)$. As $r \downarrow 0$, $k'(r\sqrt{1+u^2})$ increases to
$k'(0{+})$ for every $u$, so by monotone convergence and clause (1),
$\bar k(r) \to k'(0{+})$, in $[0,\infty]$.

\emph{Clause (3).} By Definition~\ref{def:isotropic-origin-datum}, $\nu'$ has the density
$k'(|z|)\,(|S^{d}|\,|z|^{d+1})^{-1}$ on $\RR^{d+1} \setminus \{0\}$. Let $\bar\nu$ be the measure
on $\RR^d \setminus \{0\}$ with density $\bar k(|x|)\,(|S^{d-1}|\,|x|^{d})^{-1}$. Then
\ref{eq:shadow-identity} with $h = k'$ reads $\int g(\Pi z)\,\nu'(dz) = \int g\,d\bar\nu$ for
every nonnegative Borel $g$ vanishing at the origin. The set $\{z \ne 0 : \Pi z = 0\}$ lies in a
line, which is Lebesgue-null in $\RR^{d+1}$ because $d + 1 \ge 2$, and $\nu'$ has a density. So
$\nu'$ gives it the mass $0$, and the image of $\nu'$ under $\Pi$ gives the origin no mass. Take
$g(x) = |x|^2 \wedge 1$. As
$|\Pi z| \le |z|$,
\[
\begin{aligned}
  \int_0^1 \bar k(r)\,r\,dr + \int_1^\infty \bar k(r)\,\frac{dr}{r}
   &= \int_{\RR^d}\bigl(|x|^2 \wedge 1\bigr)\,\bar\nu(dx)
   \le \int_{\RR^{d+1}}\bigl(|z|^2 \wedge 1\bigr)\,\nu'(dz) \\
   &= \int_0^1 k'(r)\,r\,dr + \int_1^\infty k'(r)\,\frac{dr}{r} < \infty ,
\end{aligned}
\]
the first equality by polar coordinates in $\RR^d$ and the last by
\ref{eq:polar-decomposition} in dimension $d + 1$. With clauses (1) and (2), $\bar k$ is finite,
nonnegative and nonincreasing and satisfies the two conditions of
Definition~\ref{def:isotropic-origin-datum}. So $(a,\bar k)$ is an isotropic origin datum in dimension $d$,
and its \Levy{} measure is $\bar\nu$. For $\xi \in \RR^d$ take $g(x) = 1 - \cos\langle\xi,x\rangle$.
Since $\langle(\xi,0),z\rangle = \langle\xi,\Pi z\rangle$ and $|(\xi,0)| = |\xi|$,
\[
  F'\bigl((\xi,0)\bigr) = \tfrac12 a|\xi|^2
     + \int_{\RR^{d+1}}\bigl(1 - \cos\langle\xi,\Pi z\rangle\bigr)\,\nu'(dz)
   = \tfrac12 a|\xi|^2 + \int_{\RR^d}\bigl(1 - \cos\langle\xi,x\rangle\bigr)\,\bar\nu(dx) ,
\]
which is the exponent \ref{eq:isotropic-exponent} of $(a,\bar k)$ at $\xi$. If
$F'(\zeta) = F(|\zeta|)$, the left side is $F(|\xi|)$.
\end{proof}

% Q-0414, round 3: the lemmas on which the proof of the theorem on the separating family rests.
Lemmas~\ref{lem:threshold-monotone} and~\ref{lem:threshold-limit} give the strict decrease of the
critical values $w_d$ and bounds of their order, and Lemma~\ref{lem:threshold-asymptotic} gives
their limit.

% shared with blueprint lem:threshold-monotone
\begin{lemma}[the critical values are strictly decreasing]\label{lem:threshold-monotone}
For a real parameter $d \ge 1$ put $m = (d+1)/2$, $\lambda_d = 2^{-d/2}$ and
$h_d(z) = z^m(e^{-z} - \lambda_d e^{-z/2})$, and let $M_d = \sup_{z>0}h_d(z)$ as in
\ref{eq:threshold-constant}. Then $M_d \in (0,\infty)$ and the supremum is attained at a point of
$(0,d\log 2)$. Consequently $w_d = \Gamma(m)/(2M_d)$ is a well-defined positive number, and when
$d$ is an integer, $k^{(w)}_d$ is nonincreasing on $(0,\infty)$ exactly when $w \le w_d$. Moreover
\[
  M_{d+1} > \sqrt{m}\;M_d \qquad\text{and}\qquad w_{d+1} < w_d
\]
for every $d \ge 1$. In particular $w_1 > w_2 > w_3 > \dots > 0$, so
$(w_{d+1}, w_d\,]$ is a nonempty interval for every integer $d \ge 1$.
\end{lemma}

% shared with blueprint lem:threshold-limit
\begin{lemma}[the critical values tend to $0$ at the rate $d^{-1/2}$]\label{lem:threshold-limit}
For a real parameter $d \ge 1$ let $m$, $h_d$, $M_d$ and $w_d = \Gamma(m)/(2M_d)$ be as in
Lemma~\ref{lem:threshold-monotone}, and put
\[
  \varepsilon_d := 2^{-d/2}e^{m/2} = e^{1/4}\Bigl(\frac e4\Bigr)^{d/4},
  \qquad K := \sqrt{\pi/2} + \sqrt\pi + 4 .
\]
Then for every $d \ge 1$
\[
  w_d \;\ge\; \frac{e^{-1/2}}{2\sqrt m},
\]
and for every $d \ge 1$ with $\varepsilon_d < 1$
\[
  w_d \;\le\; \frac{K}{2(1-\varepsilon_d)\sqrt m} .
\]
Since $e < 4$, $\varepsilon_d$ decreases strictly to $0$ as $d \to \infty$. Hence
$w_d \to 0$, and $w_d$ is of exact order $d^{-1/2}$: $\sqrt m\,w_d \ge e^{-1/2}/2$ for every
$d \ge 1$, and $\sqrt m\,w_d \le K$ for every $d$ with $\varepsilon_d \le \tfrac12$, which is every
large $d$. In particular, for every fixed $w > 0$ there is a $D$ with $w > w_d$ for all
$d \ge D$.
\end{lemma}

Neither constant of the lemma is the limit of $\sqrt m\,w_d$. The next lemma gives it.

% shared with blueprint lem:threshold-asymptotic+22be5066299c
\begin{lemma}[the limit of the critical values]\label{lem:threshold-asymptotic}
For a real parameter $d \ge 1$ let $m = (d+1)/2$, $M_d$, $w_d = \Gamma(m)/(2M_d)$ and
$\varepsilon_d = 2^{-d/2}e^{m/2}$ be as in Lemma~\ref{lem:threshold-limit}, and put
\[
  S_m := \Gamma(m)\,m^{1/2-m}\,e^{m} .
\]
\begin{enumerate}
\item \emph{(Stirling's formula.)} $S_m \to \sqrt{2\pi}$ as $m \to \infty$ through the real
numbers.
\item \emph{(The critical value against the Gamma function.)} For every $d \ge 1$,
\[
  \sqrt m\,w_d \;\ge\; \tfrac12\,S_m ,
\]
and for every $d \ge 1$ with $\varepsilon_d < 1$,
\[
  \sqrt m\,w_d \;\le\; \frac{S_m}{2\,(1-\varepsilon_d)} .
\]
\item \emph{(The limit.)} As $d \to \infty$ through the real numbers,
\[
  \sqrt m\,w_d \to \sqrt{\pi/2} \qquad\text{and}\qquad \sqrt d\,w_d \to \sqrt\pi .
\]
\end{enumerate}
\end{lemma}

% Computed values: records/affine/reviews/verify-round1-scripts/a_critical_values-output.txt.
The limit is approached slowly, and not monotonically. Computed values of $\sqrt d\,w_d$ are
$12.61$ at $d = 1$ and $2.056$ at $d = 20$, against $\sqrt\pi \approx 1.7725$. They fall below
$\sqrt\pi$ between $d = 55$ and $d = 56$, are least at $d = 81$, where the value is $1.7661$, and
increase from there on as far as they were computed, to $d = 300$. The statements are
independent of these values.

\begin{proof}[Proof of Lemma~\ref{lem:threshold-monotone}]
\emph{The closed form and the criterion.} \prog{B}{App.~A}, restated (the paragraph \emph{the polar
profiles in closed form}, carried out there for every parameter and every integer $d \ge 1$). Let
$d \ge 1$ be an integer and $w > 0$, and write
$P(s,x) = \Gamma(s)^{-1}\int_0^x v^{s-1}e^{-v}dv$ for the regularized lower incomplete gamma
function. Split \ref{eq:polar-profile-d} at the two summands of $k^{(w)}_{\mathrm I}$. For the
power, $k_{\mathrm I}(r^2/2v) = (r^2/2v)^{-1/2}$ turns the integral into a Gamma integral and
gives $2^{3/2}r^{-1}\Gamma(m)/\Gamma(d/2)$. For the bump,
$\mathbf{1}_{[1,2]}(r^2/2v) = \mathbf{1}_{[r^2/4,\,r^2/2]}(v)$. Hence \prog{B}{(A.1)}:
\begin{equation}\label{eq:witness-polar}
  k^{(w)}_d(r) = \frac{2^{3/2}\,\Gamma(m)}{\Gamma(d/2)}\,\frac1r
   + 2w\Bigl[P\Bigl(\frac d2,\frac{r^2}2\Bigr) - P\Bigl(\frac d2,\frac{r^2}4\Bigr)\Bigr],
   \qquad r > 0 ,
\end{equation}
a smooth function on $(0,\infty)$. Since $\partial_x P(s,x) = x^{s-1}e^{-x}/\Gamma(s)$, the
derivative of $P(d/2, r^2/2b)$ in $r$ is $2\,(r\,\Gamma(d/2))^{-1}(r^2/2b)^{d/2}e^{-r^2/2b}$ for
$b > 0$. With $b = 1$ and $b = 2$, and $z = r^2/2$,
\[
  r^2\,\Gamma(d/2)\,\frac{d}{dr}k^{(w)}_d(r)
   = -2^{3/2}\Gamma(m) + 4wr\Bigl(z^{d/2}e^{-z} - \bigl(\tfrac z2\bigr)^{d/2}e^{-z/2}\Bigr) .
\]
As $r = \sqrt{2z}$, the last term is $4\sqrt2\,w\,z^m\bigl(e^{-z} - 2^{-d/2}e^{-z/2}\bigr)$. So the
derivative is nonpositive at $r$ if and only if \prog{B}{(A.2)} holds:
\begin{equation}\label{eq:threshold-criterion}
  w\,z^m\bigl(e^{-z} - 2^{-d/2}e^{-z/2}\bigr) \le \tfrac12\Gamma(m),
  \qquad z = \tfrac12 r^2 .
\end{equation}

\emph{The supremum, and the equivalence.} $h_d$ is smooth on $(0,\infty)$ with $h_d(0{+}) = 0$,
and $e^{-z} - \lambda_d e^{-z/2} > 0$ exactly when $e^{-z/2} > 2^{-d/2}$, that is exactly on
$(0, d\log 2)$, an interval that is nonempty because $d \ge 1$. So $h_d > 0$ somewhere,
$h_d \le 0$ on $[d\log 2,\infty)$, and $h_d(z) \le z^me^{-z} \to 0$ as $z \to \infty$. A positive
continuous function vanishing at both ends of $(0,\infty)$ attains its supremum, so
$M_d = h_d(z^*) \in (0,\infty)$ for some $z^* \in (0, d\log 2)$. For an integer $d$,
\ref{eq:threshold-criterion} reads $w\,h_d(z) \le \tfrac12\Gamma(m)$ at $z = r^2/2$. So the
criterion holds at every $r > 0$ exactly when $w M_d \le \tfrac12\Gamma(m)$, that is exactly when
$w \le w_d$. A differentiable function on $(0,\infty)$ is nonincreasing exactly when its
derivative is nonpositive throughout, which settles the equivalence for $k^{(w)}_d$.

\emph{The first-order condition.} Let $z^* \in (0,d\log 2)$ be a maximizer and put
$q := \lambda_d e^{z^*/2}$, so that $\lambda_d e^{-z^*/2} = q\,e^{-z^*}$. Since $z^* < d\log 2$,
$q \in (0,1)$. Then $M_d = (z^*)^m e^{-z^*}(1-q) > 0$. The maximizer is interior, so
$h_d'(z^*) = 0$, where
\[
  h_d'(z) = m z^{m-1}\bigl(e^{-z} - \lambda_d e^{-z/2}\bigr)
    + z^m\bigl(-e^{-z} + \tfrac12\lambda_d e^{-z/2}\bigr) .
\]
Dividing $h_d'(z^*) = 0$ by $(z^*)^{m-1}e^{-z^*} > 0$ turns it into $m(1-q) = z^*(1 - q/2)$. As
$1 - q/2 > 0$,
\[
  z^* = m\,\frac{1-q}{1-q/2} \qquad\text{and}\qquad
  \sqrt{z^*} = \sqrt{m}\,\sqrt{\frac{1-q}{1-q/2}} .
\]

\emph{Testing the next dimension at that point.} Since
$\lambda_{d+1} = 2^{-1/2}\lambda_d$ and $m_{d+1} = m + \tfrac12$,
\[
  h_{d+1}(z^*) = (z^*)^{m+1/2}\bigl(e^{-z^*} - 2^{-1/2}q\,e^{-z^*}\bigr)
   = \sqrt{z^*}\,\bigl(1 - q/\sqrt2\bigr)\,(z^*)^m e^{-z^*}
   = \sqrt{z^*}\,\frac{1 - q/\sqrt2}{1-q}\,M_d ,
\]
so that, substituting for $\sqrt{z^*}$,
\[
  M_{d+1} \;\ge\; h_{d+1}(z^*)
   \;=\; \sqrt{m}\,M_d\,\frac{1 - q/\sqrt2}{\sqrt{(1-q)(1-q/2)}} .
\]
The last fraction exceeds $1$, because its numerator and denominator are positive and
\[
  \bigl(1 - q/\sqrt2\bigr)^2 - (1-q)\bigl(1-q/2\bigr)
   = \Bigl(1 - \sqrt2\,q + \frac{q^2}{2}\Bigr)
     - \Bigl(1 - \frac32 q + \frac{q^2}{2}\Bigr)
   = \Bigl(\frac32 - \sqrt2\Bigr)q > 0
\]
for $q > 0$, because the constant $\tfrac32 - \sqrt2 = (1 - 2^{-1/2})^2$ is positive. (It is
the reciprocal of $6 + 4\sqrt2$.) Hence $M_{d+1} > \sqrt{m}\,M_d$.

\emph{The Gamma factor.} $\Gamma$ is log-convex on $(0,\infty)$: for $x,y > 0$ and
$\theta \in (0,1)$, H\"older's inequality applied to
$t^{\theta x + (1-\theta)y - 1}e^{-t} = (t^{x-1}e^{-t})^\theta (t^{y-1}e^{-t})^{1-\theta}$ with
exponents $1/\theta$ and $1/(1-\theta)$ gives
$\Gamma(\theta x + (1-\theta)y) \le \Gamma(x)^\theta\Gamma(y)^{1-\theta}$. At
$x = m$, $y = m+1$, $\theta = \tfrac12$ this is
$\Gamma(m + \tfrac12) \le \Gamma(m)^{1/2}\Gamma(m+1)^{1/2} = \sqrt{m}\,\Gamma(m)$, by
$\Gamma(m+1) = m\Gamma(m)$. Therefore
\[
  \frac{1}{2w_{d+1}} = \frac{M_{d+1}}{\Gamma(m + \frac12)}
   > \frac{\sqrt m\,M_d}{\Gamma(m+\frac12)}
   \ge \frac{\sqrt m\,M_d}{\sqrt m\,\Gamma(m)} = \frac{1}{2w_d} ,
\]
which is $w_{d+1} < w_d$. Every $w_d$ is positive, so the chain $w_1 > w_2 > \dots > 0$ and the
nonemptiness of $(w_{d+1},w_d\,]$ follow.
\end{proof}

\begin{proof}[Proof of Lemma~\ref{lem:threshold-limit}]
\emph{Stirling's bounds.} Let $m \ge 1$ and $\varphi(u) := u - \log(1+u)$ for $u > -1$, so that
$\varphi(0) = 0$ and $\varphi'(u) = u/(1+u)$. The substitution $t = m(1+u)$ in
$\Gamma(m+1) = \int_0^\infty t^m e^{-t}\,dt$ gives
\[
  \Gamma(m+1) = m^{m+1}e^{-m}\int_{-1}^\infty e^{-m\varphi(u)}\,du ,
  \qquad\text{so}\qquad
  \Gamma(m) = m^m e^{-m}\,I_m ,\quad I_m := \int_{-1}^\infty e^{-m\varphi(u)}\,du ,
\]
by $\Gamma(m+1) = m\Gamma(m)$. Four comparisons for $\varphi$ follow by comparing derivatives and
integrating from $0$. Put $\psi(u) := u^2/(2(1+u))$ for $u \ge 0$.
\begin{itemize}
\item On $(-1,0]$, $\varphi(u) \ge u^2/2$: there $\varphi'(s) = s/(1+s) \le s$, because
$s \le 0 < 1+s \le 1$. Integrate over $[u,0]$.
\item On $[0,\infty)$, $\varphi(u) \ge \psi(u)$: $\psi(0) = 0$ and
$\psi'(s) = (s^2+2s)/(2(1+s)^2) \le (2s^2+2s)/(2(1+s)^2) = \varphi'(s)$.
\item $\psi(u) \ge u^2/4$ on $[0,1]$, since $1+u \le 2$ there, and $\psi(u) \ge u/4$ on
$[1,\infty)$, since $u/(1+u) \ge \tfrac12$ there.
\item On $[0,\infty)$, $\varphi(u) \le u^2/2$, because $\varphi'(s) \le s$ there.
\end{itemize}
With $\int_0^\infty e^{-as^2}ds = \tfrac12\sqrt{\pi/a}$ for $a > 0$, the lower bounds for
$\varphi$, on $(-1,0]$, on $[0,1]$ and on $[1,\infty)$, give
\[
  I_m \le \int_{-\infty}^0 e^{-mu^2/2}\,du + \int_0^\infty e^{-mu^2/4}\,du
   + \int_1^\infty e^{-mu/4}\,du
  = \sqrt{\frac{\pi}{2m}} + \sqrt{\frac{\pi}{m}} + \frac4m e^{-m/4}
  \le \frac{K}{\sqrt m} ,
\]
using $4/m \le 4/\sqrt m$ for $m \ge 1$. The integrand is positive and $1/\sqrt m \le 1$, so the
upper bound $\varphi(u) \le u^2/2$ gives
\[
  I_m \ge \int_0^{1/\sqrt m} e^{-mu^2/2}\,du \ge \frac{1}{\sqrt m}\,e^{-1/2} .
\]
So $e^{-1/2}\,m^{m-1/2}e^{-m} \le \Gamma(m) \le K\,m^{m-1/2}e^{-m}$ for every $m \ge 1$, and
$m = (d+1)/2 \ge 1$ because $d \ge 1$.

\emph{$M_d$ from above.} Since $2^{-d/2}e^{-z/2} > 0$, $h_d(z) < z^m e^{-z}$, and
$z \mapsto z^m e^{-z}$ has derivative $z^{m-1}e^{-z}(m - z)$, so it is largest at $z = m$. Hence
$M_d \le m^m e^{-m}$, and by Lemma~\ref{lem:threshold-monotone} $M_d > 0$, so
\[
  w_d = \frac{\Gamma(m)}{2M_d} \ge \frac{e^{-1/2}\,m^{m-1/2}e^{-m}}{2\,m^m e^{-m}}
   = \frac{e^{-1/2}}{2\sqrt m} .
\]

\emph{$M_d$ from below.} Evaluating at $z = m$,
\[
  M_d \ge h_d(m) = m^m\bigl(e^{-m} - 2^{-d/2}e^{-m/2}\bigr) = m^m e^{-m}\,(1 - \varepsilon_d) ,
\]
and $\varepsilon_d = 2^{-d/2}e^{(d+1)/4} = e^{1/4}(e/4)^{d/4}$. When $\varepsilon_d < 1$ this is
positive, and
\[
  w_d = \frac{\Gamma(m)}{2M_d} \le \frac{K\,m^{m-1/2}e^{-m}}{2\,m^m e^{-m}(1-\varepsilon_d)}
   = \frac{K}{2(1-\varepsilon_d)\sqrt m} .
\]

\emph{The limit.} Since $e < 4$, $e/4 \in (0,1)$ and $\varepsilon_d$ decreases strictly to $0$.
In particular $\varepsilon_d \le \tfrac12$ for all large $d$, and there the upper bound reads
$\sqrt m\,w_d \le K$. Since $m = (d+1)/2 \to \infty$, $w_d \to 0$ follows. For fixed $w > 0$ the
upper bound falls below $w$ from some $D$ on, and so does $w_d$.
\end{proof}

\begin{proof}[Proof of Lemma~\ref{lem:threshold-asymptotic}]
From the proof of Lemma~\ref{lem:threshold-limit}: with $\varphi(u) := u - \log(1+u)$ for $u > -1$,
\[
  \Gamma(m) = m^m e^{-m}\,I_m , \qquad I_m := \int_{-1}^\infty e^{-m\varphi(u)}\,du ,
\]
so that $S_m = \sqrt m\,I_m$. Also $\varphi(u) \ge u^2/2$ on $(-1,0]$, $\varphi(u) \ge u^2/4$ on
$[0,1]$ and $\varphi(u) \ge u/4$ on $[1,\infty)$. And $0 < M_d \le m^m e^{-m}$ for every
$d \ge 1$, $M_d \ge m^m e^{-m}(1 - \varepsilon_d)$, and $\varepsilon_d \to 0$ as $d \to \infty$.

\emph{Clause (1).} The substitution $u = v/\sqrt m$ gives
\[
  S_m = \sqrt m\,I_m = \int_{\RR} f_m(v)\,dv ,
\]
where $f_m(v) := e^{-m\varphi(v/\sqrt m)}$ for $v > -\sqrt m$ and $f_m(v) := 0$ for
$v \le -\sqrt m$.
\emph{The pointwise limit.} Since $\varphi'(s) = s/(1+s)$ and $\varphi(0) = 0$,
$\varphi(u) - \tfrac12 u^2 = -\int_0^u s^2/(1+s)\,ds$. For $|u| < 1$ the denominator is at least
$1 - |u|$ on the interval of integration, so
\[
  \bigl|\varphi(u) - \tfrac12 u^2\bigr| \le \frac{|u|^3}{3\,(1 - |u|)} .
\]
Fix a real $v$. For $m > v^2$, the point $u = v/\sqrt m$ has $|u| < 1$, so $v > -\sqrt m$ and
\[
  \bigl|m\,\varphi(v/\sqrt m) - \tfrac12 v^2\bigr|
   \le \frac{|v|^3}{3\sqrt m\,\bigl(1 - |v|/\sqrt m\bigr)} ,
\]
which tends to $0$ as $m \to \infty$. Hence $f_m(v) \to e^{-v^2/2}$ for every real $v$.
\emph{The domination.} Let $m \ge 1$ and $v > -\sqrt m$, and apply the three comparisons at
$u = v/\sqrt m$. For $-\sqrt m < v \le 0$, $m\varphi(u) \ge v^2/2$. For $0 \le v \le \sqrt m$,
$m\varphi(u) \ge v^2/4$. For $v \ge \sqrt m$, $m\varphi(u) \ge \sqrt m\,v/4 \ge v/4$. So
\[
  0 \le f_m(v) \le g(v) := e^{-v^2/4} + e^{-|v|/4} \qquad (v \in \RR,\ m \ge 1),
\]
and $g$ is integrable on $\RR$. \emph{The limit.} By dominated convergence, along every sequence
$m_n \to \infty$, $S_{m_n} \to \int_{\RR}e^{-v^2/2}\,dv = \sqrt{2\pi}$. So $S_m \to \sqrt{2\pi}$ as
$m \to \infty$.

\emph{Clause (2).} $w_d = \Gamma(m)/(2M_d) = m^m e^{-m}I_m/(2M_d)$. As $0 < M_d \le m^m e^{-m}$,
$w_d \ge I_m/2$, that is $\sqrt m\,w_d \ge S_m/2$. If $\varepsilon_d < 1$, then
$M_d \ge m^m e^{-m}(1 - \varepsilon_d) > 0$ gives $w_d \le I_m/(2(1 - \varepsilon_d))$.

\emph{Clause (3).} As $d \to \infty$, $m = (d+1)/2 \to \infty$ and $\varepsilon_d \to 0$, so
$\varepsilon_d < 1$ for all large $d$, and both sides of clause (2) tend to $\sqrt{2\pi}/2 =
\sqrt{\pi/2}$ by clause (1). Hence $\sqrt m\,w_d \to \sqrt{\pi/2}$. Finally
$\sqrt d\,w_d = \sqrt{d/m}\,\sqrt m\,w_d$ and $d/m = 2d/(d+1) \to 2$, so
$\sqrt d\,w_d \to \sqrt2\,\sqrt{\pi/2} = \sqrt\pi$.
\end{proof}

\begin{proof}[Proof of Lemma~\ref{lem:delay-polar-profile}]
For $u > 0$ let $g_u(x) = (2\pi u)^{-d/2}e^{-|x|^2/2u}$ be the heat kernel of
Lemma~\ref{lem:positive-transform-origin}. It is even and its transform is $e^{-u|\xi|^2/2}$, so
\[
  1 - e^{-u|\xi|^2/2} = \int_{\RR^d}\bigl(1 - \cos\langle\xi,x\rangle\bigr)\,g_u(x)\,dx .
\]
Integrate against $k_{\mathrm I}(u)\,du/u$ and exchange the integrals by Tonelli's theorem, which
applies because the integrand is nonnegative:
\[
  F_{\mathrm I}\bigl(\tfrac12|\xi|^2\bigr) - \tfrac12 b_0|\xi|^2
   = \int_{\RR^d}\bigl(1 - \cos\langle\xi,x\rangle\bigr)\,f(x)\,dx ,
  \qquad f(x) := \int_0^\infty g_u(x)\,k_{\mathrm I}(u)\,\frac{du}{u} ,
\]
with $f$ measurable and valued in $[0,\infty]$. Fix $x \ne 0$ and put $r = |x|$. Under the
substitution $u = r^2/2v$, $du/u$ becomes $dv/v$ with the orientation reversed, and
$g_u(x) = \pi^{-d/2}r^{-d}v^{d/2}e^{-v}$. So, by \ref{eq:polar-profile-d} and
$|S^{d-1}| = 2\pi^{d/2}/\Gamma(d/2)$,
\[
  f(x) = \pi^{-d/2}r^{-d}\int_0^\infty v^{d/2-1}e^{-v}\,k_{\mathrm I}\Bigl(\frac{r^2}{2v}\Bigr)dv
   = \frac{k_d(r)}{|S^{d-1}|\,r^d} .
\]
This is the identity of exponents. The value $k_d(r)$ is finite. For fixed $x \ne 0$ the factor
$g_u(x)/u$ is bounded for $u \in (0,1]$, because $e^{-r^2/2u}$ decays faster than every power of
$1/u$ grows, and $g_u(x) \le (2\pi)^{-d/2}$ for $u \ge 1$. The two conditions on $k_{\mathrm I}$
then bound the integral defining $f(x)$ over $(0,1]$ and over $[1,\infty)$. Measurability of $k_d$
is Tonelli's theorem again. For the integrability conditions, $\int|x|^2g_u(x)\,dx = du$ and
$\int g_u = 1$ give $\int(|x|^2 \wedge 1)\,g_u(x)\,dx \le \min\{du, 1\}$. By Tonelli's theorem
\[
  \int_{\RR^d}(|x|^2 \wedge 1)\,f(x)\,dx \le d\int_0^1 k_{\mathrm I}(u)\,du
   + \int_1^\infty k_{\mathrm I}(u)\,\frac{du}{u} < \infty ,
\]
and in polar coordinates the left side is $\int_0^1 k_d(r)\,r\,dr + \int_1^\infty k_d(r)\,dr/r$.
If $k_d$ is also nonincreasing, $(b_0,k_d)$ satisfies Definition~\ref{def:isotropic-origin-datum}. Its
\Levy{} measure has the density $f$, and its exponent \ref{eq:isotropic-exponent} is
$F_{\mathrm I}(\tfrac12|\cdot|^2)$ by the identity above.
\end{proof}

% Q-0448 (the author's decision of 2026-10-13, on review D of the second external round;
% records/affine/reviews/VERIFY-round2-mathematics.md, items 15 and 16): Theorem
% thm:filtration-strict is not printed. What it had and no other printed statement has is the
% bottom of the chain, the blueprint node lem:filtration-top, printed here verbatim. Its proof
% is the text that stood as the last part of the proof of the theorem since Q-0424: the two
% inclusions, with the three steps (a), (b) and (c) stated and proved in full, where the proof
% of record cites three child nodes that this paper does not state (lem:lk-negative-definite,
% lem:bernstein-triplet-unique, lem:dilation-decreasing-profile). It was read against the proof
% of record step by step (the first inclusion; Steps 1 to 6 of the second), and the marker is
% pinned for that reason.
% Q-0450 (the scope audit of 2026-10-13, mismatch 15): "the top of the filtration", the words of
% the lemma's title, where the prose here, in section 2.6 and in Table tab:trust-base had "the
% bottom of the chain".
The top of the filtration, $\mathcal{A}_\infty = \mathcal{S}$, is restated from the cone paper as
Proposition~\ref{prop:line-filtration}(3). The formal development proves it, as the next lemma,
from Schoenberg's radial theorem, \ledger{A10} (Theorem~\ref{thm:cited-A10}).

% shared with blueprint lem:filtration-top+9701404270b6
\begin{lemma}[the top of the filtration is the subordinated slice]\label{lem:filtration-top}
In the notation of Definition~\ref{def:filtration-thresholds}, $\mathcal{A}_\infty = \mathcal{S}$.
\end{lemma}

% A paragraph of paper prose stands between the statement and its proof.
The route is that of the proof of \prog{B}{Prop.~4.9(4)} in the cone paper, which applies the
same theorem to the same two functions. Three steps, marked (a), (b) and (c), stand in the place
of three facts which that proof cites. Each is stated and proved after the two inclusions.

\begin{proof}[Proof of Lemma~\ref{lem:filtration-top}]
For
$\mathcal{S} \subseteq \mathcal{A}_\infty$, let $F = F_{\mathrm I}(\cdot^2/2)$ with $F_{\mathrm I}$
causally admissible. $F$ is admissible by \restated{B}{Lem.~4.2}{prop:line-bridge}, and since the delay profile of
$F_{\mathrm I}$ is nonincreasing, so is its polar profile \ref{eq:polar-profile-d} in every
dimension $d$; Lemma~\ref{lem:delay-polar-profile} and Proposition~\ref{prop:isotropic-polar-reading}(1), in its
converse direction, then put $F$ in $\mathcal{A}_d$.

For $\mathcal{A}_\infty \subseteq \mathcal{S}$, let $F \in \mathcal{A}_\infty$. Being admissible,
$F$ is continuous and nonnegative on $[0,\infty)$, with $F(0) = 0$. Put
$f(\sigma) = F(\sqrt{2\sigma})$ and, for $0 < c \le 1$, $g_c(\sigma) = f(\sigma) - f(c^2\sigma)$.
Fix a dimension $d \ge 1$. Proposition~\ref{prop:isotropic-polar-reading}(1), in its forward
direction, gives an isotropic origin datum $(a_d, k_d)$ in dimension $d$ with
\[
  F(|\eta|) = \tfrac12 a_d|\eta|^2 + \int_{\RR^d}\bigl(1 - \cos\langle\eta,x\rangle\bigr)\,\nu_d(dx) ,
\]
where $\nu_d$ is its \Levy{} measure, of density $k_d(|x|)/(|S^{d-1}|\,|x|^d)$. The image
$(c\,\cdot)_*\nu_d$ of $\nu_d$ under $x \mapsto cx$ has the density
$k_d(|x|/c)/(|S^{d-1}|\,|x|^d)$, so $(c\,\cdot)_*\nu_d \le \nu_d$, because $k_d$ is nonincreasing.
And $F(c|\eta|) = \tfrac12 a_dc^2|\eta|^2 + \int(1 - \cos\langle\eta,x\rangle)\,(c\,\cdot)_*\nu_d(dx)$.
By step (a), $\eta \mapsto F(|\eta|)$ and $\eta \mapsto F(|\eta|) - F(c|\eta|)$ are negative
definite on $\RR^d$. So are $\xi \mapsto f(|\xi|^2)$ and $\xi \mapsto g_c(|\xi|^2)$, which are
these two functions at the point $\sqrt2\,\xi$. They are continuous, and in dimension $0$ they
vanish. Also $f \ge 0$, and $g_c \ge 0$: for a function $g$ with $g(0) = 0$, the inequality of
\ledger{A10} with one point $\xi_1$, $|\xi_1|^2 = \sigma$, and $c_1 = 1$ reads $2g(\sigma) \ge 0$.
So \ledger{A10} applies to $f$ and to each $g_c$. It gives $a, b \ge 0$ and a measure $m$ on
$(0,\infty)$ with $\int(1 \wedge t)\,m(dt) < \infty$ and
$f(\sigma) = a + b\sigma + \int(1 - e^{-\sigma t})\,m(dt)$ for $\sigma > 0$, and likewise
$a_c$, $b_c$ and $m_c$ for $g_c$. Here $a = 0$, since $a \le f(\sigma)$ and
$f(\sigma) \to F(0) = 0$ as $\sigma \downarrow 0$. Let $D_\lambda m$ be the image of $m$ under
$t \mapsto \lambda t$, so that
$f(c^2\sigma) = bc^2\sigma + \int(1 - e^{-\sigma t})\,D_{c^2}m(dt)$. Then
$f = f(c^2\,\cdot) + g_c$ reads
\[
  b\sigma + \int\bigl(1 - e^{-\sigma t}\bigr)\,m(dt)
  = a_c + (bc^2 + b_c)\,\sigma + \int\bigl(1 - e^{-\sigma t}\bigr)\,(D_{c^2}m + m_c)(dt)
  \qquad (\sigma > 0),
\]
with both sides finite. By step (b), $m = D_{c^2}m + m_c$, so $D_{c^2}m \le m$. As $c$ ranges over
$(0,1]$, so does $c^2$: $D_\lambda m \le m$ for every $\lambda \in (0,1]$. Since
$\int(1 \wedge t^2)\,m(dt) \le \int(1 \wedge t)\,m(dt) < \infty$, step (c) gives
$m(dt) = k(t)\,t^{-1}dt$ with $k \ge 0$ nonincreasing. The integrals $\int_0^1 k(t)\,dt$ and
$\int_1^\infty k(t)\,t^{-1}dt$ are at most $\int(1 \wedge t)\,m(dt)$. So the function
$F_{\mathrm I}$ with drift $b$ and delay profile $k$ is causally admissible, and
$F_{\mathrm I} = f$ on $(0,\infty)$. Hence $F = F_{\mathrm I}(\cdot^2/2)$, and $F \in \mathcal{S}$.

\emph{Step (a): a \Levy--Khintchine form, and the difference of two with ordered data, is negative
definite.} Let $\alpha' \le \alpha$ be real numbers, and let $L' \le L$ be measures on $\RR^d$ such
that $\int_{\RR^d}(1 - \cos\langle\eta,x\rangle)\,L(dx) < \infty$ for every $\eta \in \RR^d$. Put
\[
  \psi(\eta) = \alpha|\eta|^2 + \int_{\RR^d}\bigl(1 - \cos\langle\eta,x\rangle\bigr)\,L(dx),
  \qquad
  \psi'(\eta) = \alpha'|\eta|^2 + \int_{\RR^d}\bigl(1 - \cos\langle\eta,x\rangle\bigr)\,L'(dx),
\]
and $\phi = \psi - \psi'$. Then for every $n$, all $\eta_1, \dots, \eta_n \in \RR^d$ and all
$c_1, \dots, c_n \in \CC$,
\[
  \sum_{j,k=1}^n \bigl(\phi(\eta_j) + \phi(\eta_k) - \phi(\eta_j - \eta_k)\bigr)\,c_j\,\overline{c_k}
  \ge 0 .
\]
The route uses this with $(\alpha, L) = (\tfrac12 a_d, \nu_d)$, and with $(\alpha', L') = (0, 0)$
for $F(|\eta|)$ and $(\alpha', L') = (\tfrac12 a_dc^2, (c\,\cdot)_*\nu_d)$ for
$F(|\eta|) - F(c|\eta|)$. For the proof,
first let the coefficients be real, $a_1, \dots, a_n \in \RR$. Write $\theta_j = \langle\eta_j,x\rangle$.
Since $\cos(\theta_j - \theta_k) = \cos\theta_j\cos\theta_k + \sin\theta_j\sin\theta_k$,
\[
  (1 - \cos\theta_j) + (1 - \cos\theta_k) - \bigl(1 - \cos(\theta_j - \theta_k)\bigr)
  = (1 - \cos\theta_j)(1 - \cos\theta_k) + \sin\theta_j\sin\theta_k ,
\]
and $|\eta_j|^2 + |\eta_k|^2 - |\eta_j - \eta_k|^2 = 2\langle\eta_j,\eta_k\rangle$. Each function
$x \mapsto 1 - \cos\langle\eta,x\rangle$ is integrable for $L$ by hypothesis, so the integrals may be
added, and
\[
  \sum_{j,k=1}^n a_j a_k\bigl(\psi(\eta_j) + \psi(\eta_k) - \psi(\eta_j - \eta_k)\bigr)
  = 2\alpha\Bigl|\sum_j a_j\eta_j\Bigr|^2 + \int_{\RR^d} P_a(x)\,L(dx),
\]
where
\[
  P_a(x) = \Bigl(\sum_j a_j\bigl(1 - \cos\langle\eta_j,x\rangle\bigr)\Bigr)^2
  + \Bigl(\sum_j a_j\sin\langle\eta_j,x\rangle\Bigr)^2 \ \ge\ 0 .
\]
$P_a$ is integrable for $L$, being a finite combination of the integrable functions above. Since
$L' \le L$, every function integrable for $L$ is integrable for $L'$, so the same identity holds
for $\psi'$, with $\alpha'$ and $L'$; and $\int P_a\,dL' \le \int P_a\,dL$, $P_a$ being nonnegative.
Subtracting,
\[
  \sum_{j,k} a_j a_k\bigl(\phi(\eta_j) + \phi(\eta_k) - \phi(\eta_j - \eta_k)\bigr)
  = 2(\alpha - \alpha')\Bigl|\sum_j a_j\eta_j\Bigr|^2
  + \int P_a\,dL - \int P_a\,dL' \ \ge\ 0 .
\]
Now let the $c_j$ be complex. The matrix $K_{jk} = \phi(\eta_j) + \phi(\eta_k) - \phi(\eta_j -
\eta_k)$ is real and symmetric, because $\phi$ is even. Writing $c_j = a_j + ib_j$,
\[
  \sum_{j,k} K_{jk}\,c_j\overline{c_k}
  = \sum_{j,k} K_{jk}(a_ja_k + b_jb_k) + i\sum_{j,k} K_{jk}(b_ja_k - a_jb_k) .
\]
The imaginary part vanishes, since exchanging $j$ and $k$ changes its sign and leaves $K$ unchanged,
and the real part is the sum of the two real cases just treated, so it is nonnegative.

\emph{Step (b): the representation is unique.} Let $a, b, a', b' \ge 0$ and let $m, m'$ be measures
on $(0,\infty)$ such that, for every $\sigma > 0$,
\[
  a + b\sigma + \int_{(0,\infty)}\bigl(1 - e^{-\sigma t}\bigr)\,m(dt)
  = a' + b'\sigma + \int_{(0,\infty)}\bigl(1 - e^{-\sigma t}\bigr)\,m'(dt) < \infty .
\]
Then $a = a'$, $b = b'$ and $m = m'$. For the representations of \ledger{A10} this is the last
sentence of Theorem~3.2 of \sscite{schilling2012bernstein} (p.~21), which \ledger{A10} leaves out.
For the proof, write $G(\sigma)$ for the common value. For $t > 0$ and $\sigma > 0$,
\[
  \bigl(1 - e^{-\sigma t}\bigr) + \bigl(1 - e^{-(\sigma+2)t}\bigr)
  + e^{-\sigma t}\bigl(1 - e^{-t}\bigr)^2 = 2\bigl(1 - e^{-(\sigma+1)t}\bigr),
\]
as expanding the square shows, and $b\sigma + b(\sigma + 2) = 2b(\sigma + 1)$. Integrating against
$m$, every term being nonnegative,
\[
  G(\sigma) + G(\sigma + 2) + \int_{(0,\infty)} e^{-\sigma t}\bigl(1 - e^{-t}\bigr)^2\,m(dt)
  = 2G(\sigma + 1),
\]
and the same with $a'$, $b'$ and $m'$. All the values of $G$ are finite, so the measures
$n(dt) = (1 - e^{-t})^2\,m(dt)$ and $n'(dt) = (1 - e^{-t})^2\,m'(dt)$ on $(0,\infty)$ have the same
Laplace transform $\int e^{-\sigma t}\,n(dt) = 2G(\sigma+1) - G(\sigma) - G(\sigma+2)$ at every
$\sigma > 0$, and it is finite. Such a transform determines the measure: $n_1(dt) := e^{-t}\,n(dt)$
is finite, its image under $t \mapsto e^{-t}$ is a finite measure on $(0,1)$ whose $k$-th moment is
$\int e^{-(k+1)t}\,n(dt)$, and likewise for $n'$; so the two images have the same moments, finite
measures on $[0,1]$ are determined by their moments (by the Weierstrass approximation theorem), and
$n_1 = n'_1$, so $n = n'$. The density $(1 - e^{-t})^2$ is positive and finite on
$(0,\infty)$, so dividing by it gives $m = m'$. Then the two integrals in the hypothesis are equal
and finite, so $a + b\sigma = a' + b'\sigma$ for every $\sigma > 0$; at $\sigma = 1$ and $\sigma = 2$
this gives $b = b'$ and $a = a'$.

\emph{Step (c): a measure that decreases under dilation has a nonincreasing profile.} Let $m$ be a
measure on $(0,\infty)$ with $\int_{(0,\infty)}(1 \wedge t^2)\,m(dt) < \infty$, and suppose that
its image $D_\lambda m$ under $t \mapsto \lambda t$ is at most $m$ for every $\lambda \in (0,1]$.
Then $m(dt) = k(t)\,t^{-1}dt$ for a nonincreasing $k : (0,\infty) \to [0,\infty)$. This is the
step from the dilation inequality to the profile in the proof of \prog{V}{Lem.~7.1}, where it is
not a numbered statement. For the proof,
let $\tilde m$ be the image of $m$ under $\log$, a measure on $\RR$. Its rays are finite:
$\tilde m((\theta,\infty)) = m((e^\theta,\infty)) < \infty$, by $\int(1 \wedge t^2)\,m(dt) <
\infty$. Dilation by $\lambda = e^{-h}$, $h > 0$, becomes translation by $-h$, so the hypothesis
reads $\tilde m(A + h) \le \tilde m(A)$ for every Borel $A$ and every $h > 0$: $\tilde m$
decreases under every right shift. Put
\[
  \varphi(\theta) := \sup_{n \ge 0}\, 2^n\,\tilde m\bigl((\theta,\theta + 2^{-n}]\bigr) \in [0,\infty] .
\]
Each term is nonincreasing in $\theta$, by the shift property, and the sequence is nondecreasing in
$n$, because splitting $(\theta,\theta+2\delta]$ at its midpoint and shifting the right half gives
$\tilde m((\theta,\theta+2\delta]) \le 2\,\tilde m((\theta,\theta+\delta])$; so $\varphi$ is
nonincreasing. For $a < b$, monotone convergence and Tonelli's theorem give
$\int_a^b \varphi = \lim_n 2^n\int|(a,b] \cap [x - 2^{-n}, x)|\,\tilde m(dx)$, which lies between
$\tilde m((a + 2^{-n}, b])$ and $\tilde m((a, b + 2^{-n}])$; both tend to $\tilde m((a,b])$, the
second because the rays are finite. So $\tilde m(d\theta) = \varphi(\theta)\,d\theta$. And $\varphi$
is finite at every point: $\varphi(\theta) = \infty$ would force $\varphi = \infty$ on
$(\theta - 1, \theta]$, $\varphi$ being nonincreasing, and so
$\tilde m((\theta - 1,\theta]) = \infty$, which the finite rays exclude.
Returning to $t = e^\theta$, $m(dt) = k(t)\,t^{-1}dt$ with $k(t) := \varphi(\log t)$, finite,
nonnegative and nonincreasing.
\end{proof}

\begin{proof}[Proof of Theorem~\ref{thm:separating-family-critical-values}]
\emph{Clause (1), $w \le w_D$ gives membership.} Lemma~\ref{lem:delay-polar-profile} applies to
$D^{(w)}$ in dimension $D$, since it reads no monotonicity. It gives that the polar profile
$k^{(w)}_D$ is finite and satisfies the two conditions of Definition~\ref{def:isotropic-origin-datum}, and
that
\[
  F^{(w)}(|\xi|) = \int_{\RR^D}\bigl(1 - \cos\langle\xi,x\rangle\bigr)\,
     \frac{k^{(w)}_D(|x|)}{|S^{D-1}|\,|x|^D}\,dx , \qquad \xi \in \RR^D .
\]
If $w \le w_D$, $k^{(w)}_D$ is nonincreasing (Lemma~\ref{lem:threshold-monotone}). So
$(0,k^{(w)}_D)$ is an isotropic origin datum in dimension $D$ with exponent
$\xi \mapsto F^{(w)}(|\xi|)$, and $F^{(w)} \in \mathcal{A}_D$ by
Proposition~\ref{prop:isotropic-polar-reading}(2).

\emph{Clause (1), $w > w_D$ excludes it.} Let $w > w_D$ and suppose
$F^{(w)} \in \mathcal{A}_D$. By Proposition~\ref{prop:isotropic-polar-reading}(2),
$\xi \mapsto F^{(w)}(|\xi|)$ is the exponent of an isotropic origin datum $(a,k_D)$ in dimension
$D$, with \Levy{} measure $\nu_D$ and $k_D$ nonincreasing. Let $\nu'$ be the measure on
$\RR^D \setminus \{0\}$ with density $k^{(w)}_D(|x|)\,(|S^{D-1}|\,|x|^D)^{-1}$. By polar
coordinates in $\RR^D$, which need no monotonicity of $k^{(w)}_D$, and the two integrability
conditions,
\[
  \int\bigl(|x|^2 \wedge 1\bigr)\,\nu'(dx)
   = \int_0^1 k^{(w)}_D(r)\,r\,dr + \int_1^\infty k^{(w)}_D(r)\,\frac{dr}{r} < \infty ,
\]
and by the display $F^{(w)}(|\xi|) = \int(1 - \cos\langle\xi,x\rangle)\,\nu'(dx)$. By
\ledger{A3}(iii) there are
infinitely divisible laws with the triplets $(0,\nu',0)$ and $(aI,\nu_D,0)$. Both \Levy{}
measures are symmetric, so the odd parts of the integrands integrate to $0$, and both laws have
the characteristic function $e^{-F^{(w)}(|\xi|)}$. So they are one law, and by the uniqueness
of the triplet, \ledger{A3}(ii), $a = 0$ and $\nu' = \nu_D$. The two radial densities then agree
almost everywhere on $\RR^D$, so $k^{(w)}_D = k_D$ almost everywhere on $(0,\infty)$, by polar
coordinates. On the other hand $M_D$ is attained at some $z^* > 0$
(Lemma~\ref{lem:threshold-monotone}), and $w > w_D$ means $w\,M_D > \tfrac12\Gamma(m)$. So the
criterion of Lemma~\ref{lem:threshold-monotone} fails at $z^*$: the smooth function $k^{(w)}_D$ has a
positive derivative at $r_0 = \sqrt{2z^*}$, and is continuous and strictly increasing on an
interval around $r_0$. Take $r < s$ in that interval. Points $u_n \downarrow r$ and
$v_n \uparrow s$ at which $k^{(w)}_D = k_D$ give
$k^{(w)}_D(u_n) = k_D(u_n) \ge k_D(v_n) = k^{(w)}_D(v_n)$, and by continuity
$k^{(w)}_D(r) \ge k^{(w)}_D(s)$, a contradiction. So $F^{(w)} \notin \mathcal{A}_D$.

\emph{Clause (2).} This is Lemma~\ref{lem:threshold-monotone}: $w_{d+1} < w_d$ for every $d \ge 1$,
and every $w_d$ is positive.

\emph{Clause (3).} By Lemma~\ref{lem:threshold-asymptotic}(3), $\sqrt d\,w_d \to \sqrt\pi$ as
$d \to \infty$ through the real numbers, hence through the integers.

\emph{Clause (4).} The critical values decrease strictly
(Lemma~\ref{lem:threshold-monotone}). Let $w \in (w_{d+1}, w_d\,]$. For $D \le d$,
$w_D \ge w_d \ge w$. For $D \ge d + 1$, $w_D \le w_{d+1} < w$. If $w > w_1$, then $w > w_D$ for
every $D$. Finally $w_D \to 0$ (Lemma~\ref{lem:threshold-limit}), so for $w \in (0,w_1]$ there is a
largest integer $d \ge 1$ with $w_d \ge w$, and it is the one $d$ with
$w \in (w_{d+1}, w_d\,]$.
With clause (1), this gives the first three sentences of the clause. For the last, let $d \ge 1$
be an integer. The interval $(w_{d+1}, w_d\,]$ is nonempty by clause (2), and for $w$ in it
$F^{(w)} \in \mathcal{A}_d \setminus \mathcal{A}_{d+1}$. The inclusion
$\mathcal{A}_{d+1} \subseteq \mathcal{A}_d$ is part of Proposition~\ref{prop:profile-one-dimension-down}(1).
\end{proof}

% Q-0448 (the author's decision of 2026-10-13; VERIFY-round2-mathematics.md, item 16): the
% blueprint node thm:no-cin-ray-in-the-plane is not printed as a theorem, being clause (3) of
% Proposition prop:profile-one-dimension-down. Its proof of record is kept, verbatim, as a
% second proof of that clause. The heading carries no \ref, so that `linkage check` does not
% take it for the proof of the proposition; with no marker the text below is not compared with
% the blueprint (it was verbatim at the commit of this item). The node stays at the blueprint
% and in Lean (no_cin_ray_in_the_plane, cin_ray_scaled).
\paragraph{A second proof of Proposition~\ref{prop:profile-one-dimension-down}(3).} The clause
says that for $\tau > 0$ and $b > 0$ the $\Cin$ ray $F(\omega) = b\,\Cin(\tau|\omega|)$, the
admissible exponent with Gaussian coefficient $0$ and folded profile
$b\,\mathbf{1}_{(0,\tau)}$ (Proposition~\ref{prop:line-extreme-rays}), belongs to no
$\mathcal{A}_d$ with $d \ge 2$. The proof above takes it from clause (2) and the nesting of the
classes. The proof below projects from $\RR^d$ to the line in one step, compares the shadow on
the line of a datum in dimension $d$ with the step, and does not use the nesting. It is the
route of the formal statement \texttt{no\_cin\_ray\_in\_the\_plane}.

\begin{proof}[Second proof of clause (3)]
Let $d \ge 2$ and suppose $F \in \mathcal{A}_d$: then $\xi \mapsto F(|\xi|)$ is the exponent of a
self-decomposable law on $\RR^d$.

\emph{The datum.} By Proposition~\ref{prop:isotropic-polar-reading}(1), in its forward direction,
$\xi \mapsto F(|\xi|)$ is the exponent $F_d$ of an isotropic origin datum $(a,k_d)$ in dimension
$d$. Write $\nu := \nu_d$ for its \Levy{} measure. By \ref{eq:isotropic-exponent},
\begin{equation}\label{eq:cin-ray-exponent}
  F(|\xi|) = \tfrac12 a|\xi|^2
   + \int_{\RR^d}\bigl(1 - \cos\langle\xi,x\rangle\bigr)\,\nu(dx), \qquad \xi \in \RR^d ,
\end{equation}
and $\nu$ is rotation invariant, with the radial density
$f_d(r) = k_d(r)\bigl(|S^{d-1}|r^d\bigr)^{-1}$, where $k_d$ is nonincreasing and finite at every
$r > 0$ (Definition~\ref{def:isotropic-origin-datum}).

\emph{The shadow on a line.} Let $\Pi(x) = x_1$ and let $\bar\nu$ be the image of $\nu$ under
$\Pi$, restricted to $\RR \setminus \{0\}$. Since $|x_1| \le |x|$,
$\int(1 \wedge y^2)\,\bar\nu(dy) \le \int(1 \wedge |x|^2)\,\nu(dx) < \infty$, so $\bar\nu$ is a
\Levy{} measure on $\RR$. It is symmetric, because $\nu$ is invariant under $x \mapsto -x$. Take
$\xi = (\omega,0,\dots,0)$ in \ref{eq:cin-ray-exponent}. The integrand vanishes where $x_1 = 0$,
so
\[
  \tfrac12 a\omega^2 + \int_{\RR}\bigl(1 - \cos(\omega y)\bigr)\,\bar\nu(dy) = b\,\Cin(\tau|\omega|),
  \qquad \omega \in \RR .
\]
The right side is the exponent of the datum $(0, b\,\mathbf{1}_{(0,\tau)})$ in dimension $1$, whose
\Levy{} measure is $\nu_1(dy) := b\,\mathbf{1}_{\{0<|y|<\tau\}}\,(2|y|)^{-1}dy$, read through
\ref{eq:polar-decomposition} at $d = 1$: there $S^0 = \{\pm1\}$ and
$\int_{\RR}(1 - \cos\omega y)\,\nu_1(dy) = b\int_0^\tau(1-\cos\omega r)\,dr/r = b\,\Cin(\tau|\omega|)$,
which is \restated{B}{Lem.~3.1}{prop:line-extreme-rays} scaled by $b$. As $\bar\nu$ is symmetric, the exponent of the triplet
$(a,\bar\nu,0)$ reduces to the left side, so by \ledger{A3}(iii) the law with that triplet has
characteristic function $e^{-b\Cin(\tau|\omega|)}$, the same as the law of the datum, whose triplet
is $(0,\nu_1,0)$. By the uniqueness in \ledger{A3}(ii), $a = 0$ and $\bar\nu = \nu_1$.

\emph{The shadow determines the support.} As $\nu_1$ does not charge $\{|y| > \tau\}$,
$\nu(\{x : |x_1| > \tau\}) = \bar\nu(\{|y| > \tau\}) = 0$. A reflection carrying $e_1$ to a unit
vector $u$ preserves $\nu$, so $\nu(\{x : |\langle u,x\rangle| > \tau\}) = 0$ for every unit $u$.
Take a dense sequence $(u_n)$ in $S^{d-1}$. If $|x| > \tau$, some $u_n$ has
$|u_n - x/|x|| < (|x|-\tau)/|x|$, and then $\langle u_n,x\rangle \ge |x| - |u_n - x/|x||\,|x| >
\tau$. So $\{|x| > \tau\}$ is covered by countably many null slabs, and
$\nu(\{|x| > \tau\}) = 0$. By Lemma~\ref{lem:polar-decomposition} the image of $\nu$ under
$x \mapsto |x|$ is $k_d(r)\,dr/r$ on $(0,\infty)$, so $k_d = 0$ almost everywhere on
$(\tau,\infty)$, and since $k_d$ is nonincreasing and nonnegative, $k_d(r) = 0$ for every
$r > \tau$.

\emph{The edge.} Let $0 < \varepsilon \le \tfrac12$ and consider the slab
$S_\varepsilon = \{x : \tau(1-\varepsilon) < |x_1| < \tau\}$. By \ref{eq:polar-decomposition},
\[
  \nu(S_\varepsilon) = \frac{1}{|S^{d-1}|}\int_{S^{d-1}}\sigma_{d-1}(d\eta)\int_0^\infty
    \mathbf 1_{S_\varepsilon}(r\eta)\,k_d(r)\,\frac{dr}{r} .
\]
A point $r\eta$ of $S_\varepsilon$ with $k_d(r) \ne 0$ has $r \le \tau$, by the previous step, so
$r \ge r|\eta_1| > \tau(1-\varepsilon)$, and $|\eta_1| > \tau(1-\varepsilon)/r \ge 1-\varepsilon$:
$\eta$ lies in the cap $K_\varepsilon = \{\eta : |\eta_1| > 1-\varepsilon\}$ and
$r \in (\tau(1-\varepsilon), \tau]$, an interval of length $\tau\varepsilon$ on which
$k_d(r)/r \le M := 2k_d(\tau/2)/\tau$. Hence
\[
  \nu(S_\varepsilon) \le \frac{\sigma_{d-1}(K_\varepsilon)}{|S^{d-1}|}\,M\,\tau\varepsilon .
\]
On the other side, $\nu(S_\varepsilon) \ge \bar\nu(\{\tau(1-\varepsilon) < |y| < \tau\}) =
\nu_1(\{\tau(1-\varepsilon) < |y| < \tau\}) = b\int_{\tau(1-\varepsilon)}^\tau dr/r \ge
(b/\tau)\,\tau\varepsilon$. Dividing by $\tau\varepsilon$,
$b/\tau \le M\,\sigma_{d-1}(K_\varepsilon)/|S^{d-1}|$ for every $\varepsilon \in (0,\tfrac12]$.
As $\varepsilon \downarrow 0$ the caps decrease to the two poles $\{|\eta_1| = 1\} = \{\pm e_1\}$,
which lie on the great subsphere $\{\eta_2 = 0\}$ because $d \ge 2$, and which are therefore
$\sigma_{d-1}$-null (a great subsphere is null, the cone over it lying in a hyperplane). So $\sigma_{d-1}(K_\varepsilon) \to 0$, and $b/\tau \le 0$, a contradiction.
Hence $F \notin \mathcal{A}_d$.
\end{proof}
